\documentclass[11pt]{article}
\usepackage[T1]{fontenc}\usepackage{lmodern}\usepackage[margin=.82in]{geometry}\usepackage{microtype}
\usepackage{amsmath,amssymb,amsthm,amscd,mathtools,mathrsfs}\usepackage{booktabs,longtable,array,tabularx,multirow}
\usepackage{graphicx,float,xcolor,enumitem,fancyhdr,hyperref}\usepackage[nameinlink,noabbrev]{cleveref}
\usepackage{listings,caption,bm,needspace,placeins,siunitx,pdflscape,setspace,etoolbox,url}
\graphicspath{{figures/}}\allowdisplaybreaks\setstretch{1.03}\setlist{nosep,leftmargin=*}
\makeatletter\renewcommand*\l@subsection{\@dottedtocline{2}{1.5em}{3.35em}}\makeatother
\hypersetup{pdftitle={Projector-Seed Calculus for Integer Sequences: Spectral, Local, Adelic, and Polyphase Transfer},pdfauthor={Komorebi P.},pdfsubject={Unified realization, spectral transfer, local arithmetic, and polyphase refinement theory for structured integer sequences},colorlinks=true,linkcolor=black,citecolor=black,urlcolor=black}
\pagestyle{fancy}\fancyhf{}\fancyhead[L]{Projector-Seed Calculus for Integer Sequences}\fancyhead[R]{Komorebi P.}\fancyfoot[C]{\thepage}\setlength{\headheight}{14pt}
\newtheorem{theorem}{Theorem}[section]\newtheorem{proposition}[theorem]{Proposition}\newtheorem{lemma}[theorem]{Lemma}\newtheorem{corollary}[theorem]{Corollary}
\theoremstyle{definition}\newtheorem{definition}[theorem]{Definition}\newtheorem{example}[theorem]{Example}\newtheorem{construction}[theorem]{Construction}\newtheorem{algorithm}[theorem]{Algorithm}\newtheorem{conjecture}[theorem]{Conjecture}\newtheorem{problem}[theorem]{Problem}
\theoremstyle{remark}\newtheorem{remark}[theorem]{Remark}\newtheorem{boundary}[theorem]{Boundary}\newtheorem{certificate}[theorem]{Computational Certificate}\newtheorem{classical}[theorem]{Classical Input}\newtheorem{principle}[theorem]{Structural Principle}
\newcommand{\bbC}{\mathbb C}\newcommand{\bbF}{\mathbb F}\newcommand{\bbN}{\mathbb N}\newcommand{\bbQ}{\mathbb Q}\newcommand{\bbR}{\mathbb R}\newcommand{\bbZ}{\mathbb Z}\newcommand{\Zp}{\mathbb Z_p}\newcommand{\Qp}{\mathbb Q_p}
\newcommand{\cA}{\mathcal A}\newcommand{\cB}{\mathcal B}\newcommand{\cC}{\mathcal C}\newcommand{\cD}{\mathcal D}\newcommand{\cE}{\mathcal E}\newcommand{\cF}{\mathcal F}\newcommand{\cG}{\mathcal G}\newcommand{\cH}{\mathcal H}\newcommand{\cK}{\mathcal K}\newcommand{\cL}{\mathcal L}\newcommand{\cM}{\mathcal M}\newcommand{\cO}{\mathcal O}\newcommand{\cP}{\mathcal P}\newcommand{\cQ}{\mathcal Q}\newcommand{\cR}{\mathcal R}\newcommand{\cS}{\mathcal S}\newcommand{\cT}{\mathcal T}
\newcommand{\Id}{\operatorname{Id}}\newcommand{\ord}{\operatorname{ord}}\newcommand{\rank}{\operatorname{rank}}\newcommand{\linrank}{\operatorname{linrank}}\newcommand{\Span}{\operatorname{span}}\newcommand{\Ker}{\operatorname{Ker}}\newcommand{\diag}{\operatorname{diag}}\newcommand{\sech}{\operatorname{sech}}\newcommand{\vp}{v_p}\newcommand{\vTwo}{v_2}\newcommand{\fall}[2]{(#1)_{#2}}\newcommand{\norm}[1]{\left\lVert#1\right\rVert}\newcommand{\Had}{\mathbin{\odot}}\newcommand{\Dir}{\mathbin{\ast_D}}\newcommand{\Analysis}{\mathscr A}\newcommand{\Synthesis}{\mathscr S}
\newcolumntype{L}[1]{>{\raggedright\arraybackslash}p{#1}}\newcolumntype{Y}{>{\raggedright\arraybackslash}X}
\lstset{language=Python,basicstyle=\ttfamily\scriptsize,breaklines=true,frame=single,showstringspaces=false,columns=fullflexible,keepspaces=true,numbers=left,numberstyle=\tiny,numbersep=6pt,tabsize=4}
\title{\textbf{Projector-Seed Calculus for Integer Sequences}\\[.55em]\large A Unified Realization Theory with Spectral, Local, Adelic, and Polyphase Transfer}
\author{Komorebi P.}\date{}
\begin{document}\maketitle
\begin{abstract}
This paper develops a self-contained projector-seed realization theory for structured integer sequences and extends it through three theorem-level transfer layers. Exact residue projectors first separate lossless representation from finite structural generation and from unrestricted intended-rule inference. A finite typed grammar covers linear states, ordinary and divided-power generating functions, convolution and transform systems, automatic digit operators, nonlinear fixed points, filtered products, multiplicative arithmetic functions, residue refinement, local coordinates, and completion. Every well-founded specification has exact coefficient semantics, a finite evaluator, and an exact packet lift.

The first transfer layer is spectral. For a persistent C-finite sequence $a_n=\sum_\lambda P_\lambda(n)\lambda^n$, the spectral torsion--Fourier theorem computes the exact minimal recurrence, reduced rational denominator, Hankel rank, and state dimension of every arithmetic subsequence $a_{Nm+j}$. Root-of-unity torsion is the only source of recurrence-order reduction. One finite torsion signature classifies all packet complexities.

The second transfer layer is local and adelic. Odd-prime and dyadic Euler coordinates are unified by one source-kernel convolution while retaining distinct integral geometries. The standard-window period is proved by finite-field Fourier transfer. The dyadic kernel is evaluated exactly, giving $v_2(\kappa_{2,m})=4m-1$ and the optimal truncation modulus $2^{4M+3}$. Carmichael--Mahler transfer converts finite modular state dynamics into exact local analytic expansions and simultaneous congruence ladders. This yields odd-prime depleted Euler interpolation, direct dyadic Euler interpolation, and finite adelic synchronization.

The third transfer layer is polyphase. Mixed-radix residue refinement factors exactly, Cauchy products acquire a carry algebra, and affine reindexing $x=Mq+s$ acts on Mahler coefficients through an explicit finite kernel $\Theta_{M,s}$. These kernels compose exactly, preserve every monotone valuation ladder, become Frobenius-sparse at prime-power radix, and amplify local precision under repeated $p$-power refinement. The resulting character-polyphase lattice separates torsion phase, $p$-adic digit address, and principal-unit Mahler depth.

A proof-producing compiler implements the realization, spectral, local, and polyphase stages with exact, modular, and $p$-adic execution. Arbitrary operator storage is excluded, and no claim is made that every database entry has a canonical compact presentation or that every displayed family is historically new. The companion verification suite performs \textbf{87{,}915 exact checks with zero failures}.
\end{abstract}
\begin{center}\fbox{\begin{minipage}{.94\textwidth}\textbf{Claim boundary.} The paper proves exact realization, packet lift, spectral torsion--Fourier transfer, finite-field standard-window periodicity, the sharp dyadic Euler ladder, Carmichael--Mahler state transfer, the stated local interpolation formulas, exact polyphase refinement, Mahler-kernel composition, valuation preservation, and prime-power precision amplification under the printed hypotheses. It does not claim that a finite prefix determines its intended continuation, that every sequence has a short legal grammar, or that polyphase decomposition alone forces simplification.\end{minipage}}\end{center}
\begin{center}\fbox{\begin{minipage}{.94\textwidth}\textbf{Local and polyphase transfer law.} A finite character address resolves torsion. A $p$-power polyphase address resolves local digits. The remaining coarse variable has a Mahler expansion whose coefficients carry explicit valuation weights. Refinement is exact, its affine kernel composes, and every principal-unit refinement increases the guaranteed precision slope.\end{minipage}}\end{center}
\tableofcontents\listoffigures\listoftables\clearpage

% ===== BEGIN sections/01_scope.tex =====
\part{Foundations of Sequence Realization}
\section{The sequence problem and the meaning of universality}
An integer sequence is at once a list, a coefficient object, an orbit, a solution of equations, and a possible compression target.  These descriptions are not interchangeable.

\begin{definition}[Three reconstruction levels]
Let $a=(a_n)_{n\ge0}$ be a sequence over a ring $R$.
\begin{enumerate}[label=\textup{(\Roman*)}]
\item \emph{Lossless representation}: encode all coefficients in exact reconstructive coordinates.
\item \emph{Finite structural realization}: give finitely many seeds, typed operators, and relations that generate every coefficient.
\item \emph{Intended-rule discovery}: infer from incomplete data the particular realization intended by an external source.
\end{enumerate}
\end{definition}
Level I is universal.  Level II is the substantive target.  Level III is a model-selection problem and cannot be solved from finite terms alone.

\begin{figure}[htbp]\centering\includegraphics[width=.98\textwidth]{01_sequence_architecture.pdf}\caption{The sequence branch of the projector-seed calculus.  A finite structural specification feeds exact projection, transfer, refinement, completion, and reconstruction.}\end{figure}
\begin{figure}[htbp]\centering\includegraphics[width=.93\textwidth]{02_three_reconstruction_levels.pdf}\caption{Three meanings of reconstruction.  The realization theorem addresses Level II; the spectral extension additionally transfers exact complexity invariants through arithmetic refinement.}\end{figure}

\subsection{The target class}
No useful theorem can literally say that every displayed integer list has a short explanation.  The appropriate target is the class of \emph{finitely specified structured sequences}: sequences supplied by a finite recurrence, differential system, automaton, algebraic or compositional system, transform rule, finitely described filtered product, multiplicative prime-power rule, or finite local coordinate system.
\[
 \boxed{\text{finite seeds}+\text{restricted reusable operators}+\text{explicit relations}\Longrightarrow\text{exact infinite sequence}.}
\]

\begin{boundary}[No hidden table operators]
A presentation is not structurally finite when an operator, relation, oracle, or coefficient table contains infinitely many independent target values.  Infinite products and inverse limits are admitted only when their factor family has a finite rule and every coefficient stabilizes after finitely many factors.
\end{boundary}

\subsection{From realization to transfer}
Realization alone answers how an infinite coefficient object is generated.  A transfer theorem must do more: it must carry a global invariant through a systematic transformation and determine a nontrivial invariant of every transformed object.

The complete linear transfer extension developed in \cref{sec:spectral-transfer} has the form
\[
 \boxed{
 \begin{aligned}
 \text{global roots and Jordan-polynomial data}
 &\longrightarrow
 \text{torsion classes under }\lambda\mapsto\lambda^N\\
 &\longrightarrow
 \text{packet Fourier phases}\\
 &\longrightarrow
 \text{exact minimal recurrence and rational denominator}\\
 &\longrightarrow
 \text{exact Hankel rank and finite packet signature}.
 \end{aligned}}
\]
It converts one finite global spectral specification into a finite signature that classifies infinitely many arithmetic subsequences.  Two further layers now follow the same pattern. Finite-field character data and modular state dynamics transfer to local Mahler packets and explicit valuation ladders. Affine polyphase refinement transfers those Mahler coordinates through finite kernels whose composition, residual sparsity, and precision behavior are exact.

\subsection{Heterogeneous examples}
\begin{center}
\begin{tabular}{L{.22\textwidth}L{.30\textwidth}L{.38\textwidth}}
\toprule
Family&Primary mechanism&Structural object\\
\midrule
Fibonacci/Lucas&finite linear state&powers of a companion operator and torsion-free packet signature\\
Torsion-degenerate recurrences&spectral transfer&root-of-unity aliasing and packet Fourier cancellation\\
Bell/Touchard&nonlinear composition&set partitions and Fa\`a di Bruno blocks\\
Catalan/Fuss--Catalan&contractive fixed point&rooted operator trees\\
Partitions&filtered infinite product&coefficientwise completion\\
Thue--Morse&finite digit automaton&typed base-$2$ paths\\
Euler/Bernoulli&reciprocal exponential sections&triangular inverse coordinates\\
Odd/dyadic Euler&weighted local transfer&valuation-filtered reconstruction\\
Local state packets&Carmichael--Mahler transfer&analytic branches and congruence ladders\\
Polyphase branches&affine Mahler kernels&mixed-radix digit refinement and precision transfer\\
Multiplicative functions&prime-power seeds&finite support on the factorization of $n$\\
\bottomrule
\end{tabular}
\end{center}
The common language does not claim that these mechanisms are identical.  It says they can be represented as finite seed systems, transported into exact packet coordinates, combined through typed operators, and evaluated by finite dependency graphs.

\subsection{Principal contributions}
The paper develops:
\begin{enumerate}
\item two coefficient channels and typed seed presentations;
\item exact residue, refinement, Fourier, and CRT packet laws;
\item a master realization theorem and exact packet compiler;
\item Hankel-rank linear complexity and canonical reachable-observable reduction;
\item the spectral torsion--Fourier transfer theorem, including exact packet minimality, rank-drop classification, finite torsion signatures, rank-count quasipolynomials, Fourier conservation, and packet uncertainty;
\item C-finite, P-recursive, hypergeometric, automatic, and regular systems;
\item Riordan, Sheffer, Appell, binomial, and Stirling transform orbits;
\item Bell composition and rooted-tree fixed points;
\item filtered Euler products and prime-power reconstruction;
\item congruential Euler families with separate odd-prime and dyadic local ladders;
\item finite-field standard-window transfer, sharp dyadic weights, Carmichael--Mahler state transfer, and finite adelic synchronization;
\item exact discrete and Mahler polyphase refinement, mixed-radix composition, Frobenius sparsity, and precision amplification;
\item a proof-producing sequence compiler, family atlas, exact verification suite, and claim-status ledger.
\end{enumerate}

\subsection{Logical architecture}
The paper is organized around five theorem levels.
\begin{enumerate}
\item \textbf{Realization}: a legal finite grammar has exact coefficient semantics.
\item \textbf{Packet compilation}: every legal construction lifts exactly to compatible residue systems.
\item \textbf{Spectral transfer}: persistent C-finite data determine every minimal arithmetic-subsequence recurrence and complexity profile.
\item \textbf{Local transfer}: finite character and modular-state data determine analytic Mahler branches and valuation ladders.
\item \textbf{Polyphase transfer}: mixed-radix affine kernels determine exact branch coordinates and propagate certified precision.
\end{enumerate}
The realization and packet levels apply throughout the declared grammar. The three transfer layers are established under their printed spectral, local-state, and Mahler hypotheses; analogous minimality laws remain open in broader nonlinear and holonomic strata.

\subsection{Status convention}
Theorems are proved under printed hypotheses.  A \emph{Classical Input} records a standard result, a \emph{Computational Certificate} records finite exact verification, a \emph{Boundary} blocks overreading, and a \emph{Conjecture} is not used as a premise of a proved result.  Historical priority is separate from internal correctness and remains an external audit task where explicitly noted.

% ===== END sections/01_scope.tex =====

% ===== BEGIN sections/02_foundations.tex =====
\section{Coefficient universes, seeds, and finite presentations}
\subsection{Ordinary and divided-power channels}
Let $R$ be commutative. A sequence has OGF $A(z)=\sum a_nz^n$. Over a $\bbQ$-algebra it has EGF $\widehat A(t)=\sum a_nt^n/n!$. Over an arbitrary ring use divided powers $t^{[n]}$ with $t^{[r]}t^{[s]}=\binom{r+s}{r}t^{[r+s]}$. Ordinary multiplication gives Cauchy convolution; divided-power multiplication gives binomial convolution.
\begin{definition}[Coefficient object]
A coefficient object is a typed element of $R^{\bbN}$, $R[[z]]$, a divided-power series module, a finite-dimensional state module, an arithmetic-function module with Dirichlet convolution, or a filtered/local completion. The type declaration determines legal operations.
\end{definition}
\begin{definition}[Typed projector-seed presentation]
A presentation is a tuple $\mathfrak P=(\mathcal X,S,\mathcal O,\mathcal R,\Lambda,\tau)$ of finitely many object types, seeds, typed operators or finitely ruled operator families, relations, output maps, and an optional topology for completion.
\end{definition}
A realization is \emph{finite-state linear} if its state is finite-dimensional and evolution is linear. It is \emph{finite-grammar} if its syntax is finite even when its semantic value is an infinite formal series. Bell numbers are finite-grammar through $E'=E$, $B'=EB$, $E(0)=B(0)=1$, but are not C-finite.
\subsection{Generated hulls}
Let $F$ be a free typed operator algebra and $s$ a seed. Evaluation $\operatorname{ev}_s:F\to\mathcal X$ sends $W$ to $W(s)$.
\begin{proposition}[Universal seed presentation]
\[\boxed{F/\ker(\operatorname{ev}_s)\cong\operatorname{im}(\operatorname{ev}_s).}\]
\end{proposition}
\begin{proof}This is the first isomorphism theorem.\end{proof}
\subsection{Finite prefixes do not identify an intended rule}
\begin{theorem}[Finite-prefix non-identifiability]
For integers $a_0,\ldots,a_m$, let $P(n)=\sum_{r=0}^m\Delta^ra_0\binom nr$. For every $c\in\bbZ$,
\[\boxed{P_c(n)=P(n)+c\binom n{m+1}}\]
is integer-valued, agrees with the prefix for $0\le n\le m$, and gives a distinct continuation as $c$ varies.
\end{theorem}
\begin{proof}Newton interpolation gives the prefix. The added binomial polynomial vanishes through $m$ and equals one at $m+1$.\end{proof}
\begin{corollary}No data-only finite-prefix recognizer can guarantee a unique intended sequence without a prior grammar, probability model, complexity rule, or external definition.\end{corollary}
\subsection{Complexity}
For a fixed grammar define $K_{\mathfrak G}(a)$ as the minimum syntax size of a valid presentation. A more informative profile records linear state dimension, grammar size, maximal arity, packet width, local depth, dependency horizon, and certificate cost. It is grammar-relative; unrestricted program equivalence inherits undecidability.

% ===== END sections/02_foundations.tex =====

% ===== BEGIN sections/03_packets.tex =====
\section{Exact residue packets and refinement}
For $A(z)=\sum a_nz^n$, define $P_{N,j}A=\sum_{n\equiv j\pmod N}a_nz^n$.
\begin{theorem}[Exact residue reconstruction]
\[P_{N,i}P_{N,j}=\delta_{ij}P_{N,j},\qquad\sum_{j=0}^{N-1}P_{N,j}=I,\qquad A=\sum_{j=0}^{N-1}P_{N,j}A.\]
If $N\mid M$, then $P_{N,j}=\sum_{r\equiv j\,(N)}P_{M,r}$.
\end{theorem}
\begin{proof}Each coefficient index lies in exactly one residue class, and a class modulo $N$ is the disjoint union of compatible classes modulo $M$.\end{proof}
\begin{figure}[htbp]\centering\includegraphics[width=.93\textwidth]{03_residue_selector_matrix.pdf}\caption{A residue selector matrix. Every coefficient is selected once and only once.}\end{figure}
Define analysis and synthesis by $\Analysis_N(A)=(P_{N,0}A,\ldots,P_{N,N-1}A)$ and $\Synthesis_N(A_0,\ldots,A_{N-1})=\sum_jA_j$. They are inverse on packet-supported tuples.
\subsection{Fourier form}
When $N$ is invertible and a primitive root $\omega_N$ is available,
\[\boxed{P_{N,j}A(z)=\frac1N\sum_{r=0}^{N-1}\omega_N^{-jr}A(\omega_N^rz).}\]
The coefficient-selector definition remains valid when this Fourier formula is not integral.
\subsection{Packet operation laws}
Write $A_j=P_{N,j}A$, $B_j=P_{N,j}B$, indices modulo $N$.
\begin{theorem}[Exact packet laws]
\begin{align*}P_{N,k}(A+B)&=A_k+B_k,&P_{N,k}(z^rA)&=z^rA_{k-r},\\P_{N,k}(A')&=(A_{k+1})',&P_{N,k}(AB)&=\sum_{r+s\equiv k}A_rB_s,\\P_{N,k}(A\Had B)&=A_k\Had B_k,&P_{N,k}(A\circ B)&=P_{N,k}\sum_{n\ge0}a_nB^n,\quad B(0)=0.
\end{align*}
\end{theorem}
\begin{proof}Track degrees modulo $N$; products add packet addresses and Hadamard products preserve the common index.\end{proof}
Cyclic packet multiplication $(U\star_NV)_k=\sum_{r+s\equiv k}U_rV_s$ makes $\Analysis_N$ an algebra isomorphism onto the packet-supported product space.
\subsection{CRT factorization}
\begin{theorem}[CRT packets]
If $(M,N)=1$, then
\[\boxed{P_{MN,r}=P_{M,r\bmod M}P_{N,r\bmod N}.}\]
\end{theorem}
\begin{proof}This is the Chinese remainder theorem applied to coefficient addresses.\end{proof}
\subsection{Nonlinear conjugacy}
For any map $F:X\to Y$ with packet systems $\cP,\cQ$, define $F_{\cQ\leftarrow\cP}=\Analysis_{\cQ}F\Synthesis_{\cP}$.
\begin{theorem}[Exact nonlinear packet coordinates]
\[\boxed{\Synthesis_{\cQ}F_{\cQ\leftarrow\cP}\Analysis_{\cP}=F,\qquad(G\circ F)_{\cR\leftarrow\cP}=G_{\cR\leftarrow\cQ}\circ F_{\cQ\leftarrow\cP}.}\]
\end{theorem}
\begin{proof}Insert $\Synthesis\Analysis=I$.\end{proof}
\begin{boundary}Exact packets need not be independent. Every output packet may depend on all input packets.\end{boundary}

% ===== END sections/03_packets.tex =====

% ===== BEGIN sections/04_master_realization.tex =====
\section{The master sequence-realization theorem}
\subsection{Restricted grammar}
The grammar is generated by: finite seeds; linear combinations, shifts, differences, sums, differentiation and integration; Cauchy, divided-power, Hadamard, matrix and Dirichlet products; unit reciprocals and unit-triangular inverses; composition with zero inner constant; Riordan/Sheffer transforms; finite state and digit evolution; contractive formal fixed points; exact residue refinement; filtered sums/products whose weights tend to infinity; local Mahler or valuation-filtered coordinates; and exact affine polyphase refinement with finite mixed-radix kernels.
\begin{definition}[Well founded]
A finite specification is well founded when each implicit strongly connected component is a finite linear system with invertible lead, a unit-triangular coefficient system, a strict contraction in a declared topology, or a recurrence with sufficient initial data.
\end{definition}
For $F\in R[[z]]$, let $\ord_zF$ be the least nonzero degree.
\begin{theorem}[Formal fixed point]
If $\Phi:(zR[[z]])^d\to(zR[[z]])^d$ satisfies $\ord_z(\Phi(U)-\Phi(V))\ge\ord_z(U-V)+1$, then it has a unique fixed point. Iteration from zero stabilizes the first $m$ coefficient layers after at most $m$ steps.
\end{theorem}
\begin{proof}Successive differences gain one order. Completeness gives a limit; applying the same inequality to two fixed points forces equality.\end{proof}
\begin{lemma}[Finite-prefix closure]
Exact finite-prefix evaluation is preserved by the grammar operations. In a filtered product $\prod_mF_m$ with $F_m\equiv1\pmod{z^{w(m)}}$ and $w(m)\to\infty$, only finitely many factors affect any requested prefix.
\end{lemma}
\begin{proof}Each explicit operation has a finite degree dependency. Unit inversion is triangular, formal composition uses finitely many powers per degree, recurrences and digit paths are finite, contractions stabilize, and deep filtered factors are invisible.\end{proof}
\begin{theorem}[Master projector-seed sequence realization]
Every well-founded finite specification determines a unique semantic coefficient object, a finite typed projector-seed presentation, an exact finite-prefix evaluator, exact modular/local evaluators where operations are defined, exact polyphase branch evaluators, and a finite proof-dependency graph reducing generated identities to primitive relations and closure laws.
\end{theorem}
\begin{proof}Topologically evaluate explicit components. Solve implicit blocks by linear inversion, triangular recursion, recurrence evolution, or contraction. Filtered components converge coefficientwise. Uniqueness holds componentwise.\end{proof}
\begin{theorem}[Exact packet compiler]
For every well-founded specification $\mathfrak P$ and modulus $N$, replacing each sequence node by $\Analysis_N$ and each operator by its exact packet conjugate produces a finite packet specification $\mathfrak P^{[N]}$. Synthesis reconstructs the original object, all equations lift exactly, refinement commutes with evaluation, and fixed points, composition, inversion, and filtered limits are preserved.
\end{theorem}
\begin{proof}Use $F^{[N]}=\Analysis_NF\Synthesis_N$, the nonlinear conjugacy theorem, explicit packet laws, and continuity of finite projectors.\end{proof}
\begin{corollary}[Broad structured coverage]
The grammar contains every C-finite sequence, every P-recursive sequence with initial data, every finite hypergeometric term rule, every automatic or finite linear digit sequence, every finite Riordan/Sheffer/Appell/binomial/Stirling expression, every admissible finite algebraic or contractive system, every effectively filtered product, every multiplicative function with a finite prime-power rule, and finite compositions, residue refinements, local Mahler packets, and affine polyphase refinements of these classes.
\end{corollary}
\begin{figure}[htbp]\centering\includegraphics[width=.97\textwidth]{04_sequence_grammar_closure.pdf}\caption{The restricted grammar. It is broad, but it excludes arbitrary infinite tables disguised as operators.}\end{figure}
\begin{boundary}This is a compiler theorem from structural specifications, not a theorem that finitely many terms reveal the specification.\end{boundary}

% ===== END sections/04_master_realization.tex =====


% ===== BEGIN sections/05_linear_recurrences.tex =====
\part{Linear Realizations and Spectral Transfer}
\section{Finite linear states, rational series, and recurrence packets}

Finite linear evolution is the most rigid part of the sequence calculus.  It is also the point at which seed rank, recurrence order, rational generating functions, Hankel rank, and residue-jump operators can be compared exactly.

\subsection{Finite linear realization}
\begin{definition}[Linear seed system]
A linear realization over a field $K$ is a triple $(C,v,\lambda)$ with $C\in M_d(K)$, $v\in K^d$, and $\lambda\in (K^d)^*$, producing
\[
 a_n=\lambda C^n v\qquad(n\geq0).
\]
The reachable space is $\mathcal R=\Span\{C^nv:n\geq0\}$ and the observable quotient removes those $x\in\mathcal R$ for which $\lambda C^nx=0$ for all $n$.
\end{definition}

\begin{theorem}[Cayley--Hamilton recurrence]
If $a_n=\lambda C^nv$ and
\[
 \chi_C(x)=x^d-c_{d-1}x^{d-1}-\cdots-c_1x-c_0,
\]
then
\[
 \boxed{a_{n+d}=c_{d-1}a_{n+d-1}+\cdots+c_1a_{n+1}+c_0a_n.}
\]
\end{theorem}
\begin{proof}
Apply $\chi_C(C)=0$ to $C^nv$ and then apply $\lambda$.
\end{proof}

\begin{theorem}[Equivalent C-finite descriptions]
For a sequence over a field $K$, the following are equivalent.
\begin{enumerate}[label=\textup{(\roman*)}]
\item The sequence satisfies a homogeneous constant-coefficient recurrence from some point onward.
\item Its ordinary generating function is rational.
\item It admits a finite-dimensional linear realization.
\end{enumerate}
With initial transients included in the state, ``from some point onward'' may be replaced by ``for all sufficiently supported indices.''
\end{theorem}
\begin{proof}
A recurrence gives a companion state.  A companion state gives a rational resolvent
\[
 A(z)=\lambda(I-zC)^{-1}v.
\]
Conversely, clearing the denominator of a rational series gives a constant-coefficient recurrence.
\end{proof}

\subsection{Minimality and Hankel rank}
Let $H(a)=(a_{i+j})_{i,j\geq0}$ be the infinite Hankel matrix and let $H_{r,s}$ denote its finite northwest blocks.

\begin{theorem}[Minimal linear dimension equals Hankel rank]
For a sequence $a$ over a field,
\[
 \boxed{\min\{d:a_n=\lambda C^nv\}=\rank H(a),}
\]
where the right side is the supremum of the ranks of finite Hankel blocks.  In particular, $a$ is C-finite exactly when $H(a)$ has finite rank.
\end{theorem}
\begin{proof}
If $a_n=\lambda C^nv$, then
\[
 a_{i+j}=(\lambda C^i)(C^jv),
\]
so every finite Hankel block factors through $K^d$ and has rank at most $d$.

Conversely, suppose $H(a)$ has finite rank $r$.  Let $W$ be the span of the columns
\[
 h_j=(a_j,a_{j+1},a_{j+2},\ldots)^T.
\]
The left shift sends $h_j$ to $h_{j+1}$ and therefore induces a linear map $C:W\to W$.  Take $v=h_0$ and let $\lambda$ read the first coordinate.  Then $a_n=\lambda C^nv$.  Minimality follows from the lower bound.
\end{proof}

\begin{corollary}[Canonical reachable-observable reduction]
Every finite linear seed system has a minimal realization obtained by restricting to its reachable space and quotienting by its unobservable subspace.  Any two minimal realizations of the same sequence are similar.
\end{corollary}
\begin{proof}
The reduction does not change $\lambda C^nv$.  Both minimal systems identify with the shift action on the column space of $H(a)$.
\end{proof}

This gives the first canonical complexity invariant in the sequence calculus:
\[
 \operatorname{linrank}(a)=\rank H(a).
\]
It is invariant under similarity, independent of a chosen recurrence presentation, and computable from a sufficiently large exact Hankel block once a finite-rank certificate is supplied.

\subsection{Jump operators and residue subsequences}
\begin{theorem}[Universal jump realization]
Let $a_n=\lambda C^nv$.  For a modulus $N\geq1$ and residue $0\leq j<N$, define $u_m^{(j)}=a_{Nm+j}$.  Then
\[
 \boxed{u_m^{(j)}=\lambda (C^N)^m C^jv.}
\]
Consequently every residue subsequence is generated by the same jump operator $J_N=C^N$ and satisfies the characteristic recurrence of $J_N$.
\end{theorem}
\begin{proof}
Use $C^{Nm+j}=(C^N)^mC^j$ and Cayley--Hamilton for $C^N$.
\end{proof}

\begin{corollary}[Packet recurrence compiler]
If
\[
 \chi_{C^N}(x)=x^d-q_{d-1}^{(N)}x^{d-1}-\cdots-q_0^{(N)},
\]
then for every $j$,
\[
 u_{m+d}^{(j)}=\sum_{r=0}^{d-1}q_r^{(N)}u_{m+r}^{(j)}.
\]
The coefficients depend on $N$ but not on the packet address $j$.
\end{corollary}

\subsection{The exact-minimality problem}
Residue splitting never increases the dimension required by the inherited jump realization, but the inherited characteristic polynomial of $C^N$ need not be minimal for an individual packet.  Two effects must be resolved:
\begin{enumerate}
\item distinct characteristic roots may acquire the same $N$th power;
\item the packet seed phase $C^jv$ may cancel part or all of the merged Jordan data.
\end{enumerate}
The next section solves both effects exactly.  It transfers the global characteristic roots and their polynomial multiplicity data into the minimal recurrence, reduced rational denominator, Hankel rank, and complete periodic rank profile of every arithmetic subsequence.

\subsection{Fibonacci and Lucas systems}
Let
\[
 C=\begin{pmatrix}1&1\\1&0\end{pmatrix},\qquad
 v=\begin{pmatrix}1\\0\end{pmatrix}.
\]
Then
\[
 C^n v=\begin{pmatrix}F_{n+1}\\F_n\end{pmatrix}.
\]
The trace and determinant of $C^N$ are
\[
 \operatorname{tr}(C^N)=L_N,\qquad \det(C^N)=(-1)^N.
\]

\begin{theorem}[Universal Fibonacci packet law]
For every $N\geq1$ and $0\leq j<N$, the sequence $u_m=F_{Nm+j}$ satisfies
\[
 \boxed{u_{m+2}=L_Nu_{m+1}-(-1)^Nu_m.}
\]
The companion Lucas packet $v_m=L_{Nm+j}$ satisfies the same recurrence.
\end{theorem}
\begin{proof}
The characteristic polynomial of $C^N$ is
\[
 x^2-L_Nx+(-1)^N.
\]
Apply the jump theorem.
\end{proof}

For $N=3$ this becomes
\[
 u_{m+2}=4u_{m+1}+u_m.
\]
The three packet seeds are
\[
 (F_0,F_3,F_6,\ldots),\quad(F_1,F_4,F_7,\ldots),\quad(F_2,F_5,F_8,\ldots),
\]
and all evolve by the same jump matrix $C^3$.

\begin{figure}[htbp]
\centering
\includegraphics[width=.82\textwidth]{05_fibonacci_jump_packets.pdf}
\caption{Fibonacci residue packets for modulus three.  Their initial states differ, while their transfer operator is the common jump $C^3$.}
\end{figure}

\subsection{General Lucas sequences and higher recurrences}
Let $U_{n+2}=PU_{n+1}-QU_n$ with companion matrix
\[
 C_{P,Q}=\begin{pmatrix}P&-Q\\1&0\end{pmatrix}.
\]
Write $V_N=\operatorname{tr}(C_{P,Q}^N)$.  Since $\det C_{P,Q}=Q$,
\[
 \boxed{U_{N(m+2)+j}=V_NU_{N(m+1)+j}-Q^N U_{Nm+j}.}
\]
Thus Pell, Jacobsthal, Mersenne-type, and general Lucas pairs are one packet theorem, not isolated formulas.

For a recurrence of order $d$, the coefficients of the jump recurrence are the elementary symmetric functions of the $N$th powers of the roots, with the usual Jordan-polynomial corrections when roots repeat.  Equivalently they are obtained exactly from $C^N$ without introducing algebraic roots.

\begin{algorithm}[Exact linear packet compilation]
Given a companion matrix $C$, modulus $N$, and desired packet $j$:
\begin{enumerate}
\item compute $J=C^N$ by exact binary exponentiation;
\item compute the packet seed $v_j=C^jv$;
\item minimize the reachable-observable triple $(J,v_j,\lambda)$ if desired;
\item compute $\chi_J$ or the minimal polynomial on the packet state;
\item emit the recurrence, initial packet values, and a Cayley--Hamilton certificate.
\end{enumerate}
\end{algorithm}

\subsection{Polynomial and quasipolynomial sequences}
A polynomial sequence $a_n=P(n)$ of degree $r$ satisfies
\[
 \Delta^{r+1}a_n=0.
\]
It therefore has a unipotent linear realization of dimension at most $r+1$ in the Newton basis $\binom nr$.  A quasipolynomial of period $N$ is a finite sum of polynomial residue packets, hence has a finite realization combining a cyclic address state with unipotent polynomial states.

\begin{proposition}[Quasipolynomial normal form]
A sequence is quasipolynomial with period dividing $N$ and degree at most $r$ exactly when it lies in the span of
\[
 n\longmapsto \binom{\lfloor n/N\rfloor}{q}\mathbf 1_{n\equiv j\,(N)},
 \qquad 0\leq q\leq r,
 \quad 0\leq j<N.
\]
\end{proposition}
\begin{proof}
Restrict to each residue class and use the Newton basis for polynomial functions of the packet index.
\end{proof}

\subsection{What the linear theory does and does not cover}
Finite linear states capture rational generating functions and every constant-coefficient recurrence, but not Bell numbers, partition numbers, or generic algebraic series.  Those families require larger finite grammars.  The value of the projector-seed architecture is that the same packet, refinement, transport, and certification language continues after finite-dimensional linearity ends.

% ===== END sections/05_linear_recurrences.tex =====

% ===== BEGIN sections/06_spectral_torsion_fourier.tex =====
\section{Spectral torsion--Fourier transfer and exact packet complexity}
\label{sec:spectral-transfer}

The universal jump realization shows that every arithmetic subsequence of a C-finite sequence is generated by the power $C^N$ of one global state operator.  That inherited realization is generally not minimal.  The purpose of this section is to determine the minimal packet recurrence exactly and to identify the unique mechanism by which its order can fall.

The result transfers one finite collection of global spectral data into the complete family of arithmetic-subsequence recurrences, Hankel ranks, rational denominators, and packet-complexity profiles.  The transfer has two components:
\[
 \boxed{
 \text{decimation by }N
 =
 \text{torsion aliasing under }\lambda\mapsto\lambda^N
 +
 \text{finite Fourier phase selection by }j.
 }
\]

\subsection{Persistent C-finite sequences and active modes}
Let $K$ be a field of characteristic zero and let $\overline K$ be an algebraic closure.

\begin{definition}[Persistent C-finite sequence]
A C-finite sequence is called \emph{persistent} when its minimal recurrence polynomial has nonzero constant term.  Equivalently, every active characteristic root is nonzero, or every minimal state transition operator is invertible.
\end{definition}

The adjective excludes only finite-support nilpotent transients.  Those are restored in \cref{subsec:finite-transient-extension}.  The zero sequence has rank zero in every packet and is omitted from the active-mode notation below.

\begin{theorem}[Active exponential-polynomial decomposition]
\label{thm:active-modal-decomposition}
Every nonzero persistent C-finite sequence has a unique expansion
\[
 \boxed{
 a_n=\sum_{\lambda\in\Lambda}P_\lambda(n)\lambda^n,
 }
\]
where $\Lambda\subset\overline K^\times$ is finite, its elements are distinct, and every $P_\lambda\in\overline K[x]$ is nonzero.  If
\[
 d_\lambda=\deg P_\lambda,
\]
then the minimal recurrence polynomial over $\overline K$ is
\[
 \boxed{
 \chi_a(X)=\prod_{\lambda\in\Lambda}(X-\lambda)^{d_\lambda+1},
 }
\]
and
\[
 \boxed{
 \rho(a):=\operatorname{linrank}(a)
 =\rank H(a)
 =\deg\chi_a
 =\sum_{\lambda\in\Lambda}(d_\lambda+1).
 }
\]
\end{theorem}

\begin{proof}
Put a minimal state operator into Jordan form over $\overline K$.  A Jordan block of size $r$ at eigenvalue $\lambda\neq0$ contributes a polynomial of degree at most $r-1$ times $\lambda^n$.  Reachability and observability remove every inactive leading coefficient.  Conversely, the sequences $n^q\lambda^n$ for distinct pairs $(\lambda,q)$ are linearly independent and have minimal annihilators $(E-\lambda)^{q+1}$, where $E$ is the forward shift.  Pairwise coprimality for distinct $\lambda$ gives the product formula and the dimension identity.  Uniqueness follows from the same linear independence.
\end{proof}

\begin{remark}[Descent to the base field]
The displayed factorization is over $\overline K$.  Because $a_n\in K$ and the active decomposition is unique, the collection of active modes is Galois stable.  Hence the product defining $\chi_a$ descends to the monic minimal recurrence polynomial in $K[X]$.
\end{remark}

\subsection{Torsion classes and packet-phase polynomials}
For $N\geq1$ and any address $j\geq0$, define the decimated sequence
\[
 a^{[N,j]}_m=a_{Nm+j}.
\]
The standard residue packets use $0\leq j<N$, but allowing all $j\geq0$ makes the periodicity theorem transparent.

\begin{definition}[$N$-torsion equivalence]
For active roots $\lambda,\eta\in\Lambda$, write
\[
 \boxed{
 \lambda\sim_N\eta
 \quad\Longleftrightarrow\quad
 \lambda^N=\eta^N
 \quad\Longleftrightarrow\quad
 (\lambda/\eta)^N=1.
 }
\]
Thus the equivalence classes are the fibers of the power map $\lambda\mapsto\lambda^N$ on the active spectrum.
\end{definition}

For each class $C\in\Lambda/{\sim_N}$ choose a representative $\lambda_C$.  Write
\[
 \zeta_\lambda=\frac{\lambda}{\lambda_C}
 \qquad(\lambda\in C).
\]
Then $\zeta_\lambda^N=1$.

\begin{definition}[Packet-phase polynomial]
For $C\in\Lambda/{\sim_N}$ and $j\geq0$, define
\[
 \boxed{
 R_{C,j}(x)=\sum_{\lambda\in C}\zeta_\lambda^jP_\lambda(x).
 }
\]
Changing the representative multiplies $R_{C,j}$ by a nonzero scalar and does not change whether it vanishes or what its degree is.
\end{definition}

The polynomial $R_{C,j}$ is a finite Fourier combination of all modes that become spectrally indistinguishable after the $N$-jump.  Its degree records the Jordan information surviving at packet address $j$.

\begin{figure}[htbp]
\centering
\includegraphics[width=.97\textwidth]{16_spectral_torsion_transfer.pdf}
\caption{The spectral transfer mechanism.  The power map merges roots along torsion classes, packet addresses apply finite Fourier phases inside each class, and the surviving polynomial degrees give the exact minimal packet recurrence.}
\label{fig:spectral-transfer}
\end{figure}

\subsection{The spectral torsion--Fourier packet transfer theorem}

\begin{theorem}[Spectral torsion--Fourier packet transfer]
\label{thm:spectral-torsion-fourier}
Let $a_n=\sum_{\lambda\in\Lambda}P_\lambda(n)\lambda^n$ be a persistent C-finite sequence over a characteristic-zero field.  For every $N\geq1$ and $j\geq0$ the following statements hold.

\begin{enumerate}[label=\textup{(\arabic*)}]
\item \textbf{Exact modal transfer.}
\[
 \boxed{
 a_{Nm+j}
 =
 \sum_{C\in\Lambda/{\sim_N}}
 \lambda_C^j R_{C,j}(Nm+j)(\lambda_C^N)^m.
 }
\]

\item \textbf{Exact minimal recurrence.}  The minimal characteristic polynomial of $a^{[N,j]}$ over $\overline K$ is
\[
 \boxed{
 \chi_{N,j}(X)
 =
 \prod_{\substack{C\in\Lambda/{\sim_N}\\R_{C,j}\not\equiv0}}
 (X-\lambda_C^N)^{1+\deg R_{C,j}}.
 }
\]
This polynomial is Galois invariant and descends to $K[X]$.

\item \textbf{Exact packet complexity.}
\[
 \boxed{
 \rho_a(N,j)
 :=\operatorname{linrank}(a^{[N,j]})
 =\rank H(a^{[N,j]})
 =
 \sum_{\substack{C\in\Lambda/{\sim_N}\\R_{C,j}\not\equiv0}}
 (1+\deg R_{C,j}).
 }
\]

\item \textbf{Complete rank-drop criterion.}
\[
 \boxed{
 \rho_a(N,j)=\rho(a)
 \quad\Longleftrightarrow\quad
 \lambda\mapsto\lambda^N
 \text{ is injective on }\Lambda.
 }
\]
If one torsion collision occurs, then $\rho_a(N,j)<\rho(a)$ for every address $j$.

\item \textbf{Exact denominator transfer.}  If
\[
 U_{N,j}(w)=\sum_{m\geq0}a_{Nm+j}w^m,
\]
then its reduced denominator over $\overline K$ is
\[
 \boxed{
 Q_{N,j}(w)
 =
 \prod_{\substack{C\in\Lambda/{\sim_N}\\R_{C,j}\not\equiv0}}
 (1-\lambda_C^Nw)^{1+\deg R_{C,j}}.
 }
\]

\item \textbf{Exact periodicity.}  Let
\[
 \boxed{
 L_a=
 \operatorname{lcm}\left\{
 \operatorname{ord}(\lambda/\eta):
 \lambda,\eta\in\Lambda,\ \lambda\neq\eta,
 \ \lambda/\eta\text{ a root of unity}
 \right\},
 }
\]
with $L_a=1$ when the set is empty.  Then
\[
 \boxed{
 \rho_a(N+L_a,j)=\rho_a(N,j),
 \qquad
 \rho_a(N,j+L_a)=\rho_a(N,j).
 }
\]
Thus $\rho_a(N,j)$ depends only on $(N\bmod L_a,j\bmod L_a)$.

\item \textbf{Finite classification.}  The map
\[
 \boxed{
 \mathfrak T_a:(\bbZ/L_a\bbZ)^2\longrightarrow\bbZ_{\geq0},
 \qquad
 \mathfrak T_a(\overline N,\overline j)=\rho_a(N,j),
 }
\]
is a well-defined finite invariant that determines the minimal recurrence order of every arithmetic subsequence of $a$.
\end{enumerate}
\end{theorem}

\subsection{Proof by classwise modal reduction}
The proof is separated into local lemmas so that each certificate can be checked independently.

\begin{lemma}[Classwise transfer identity]
\label{lem:classwise-transfer}
For every $C\in\Lambda/{\sim_N}$,
\[
 \sum_{\lambda\in C}P_\lambda(Nm+j)\lambda^{Nm+j}
 =
 \lambda_C^j(\lambda_C^N)^mR_{C,j}(Nm+j).
\]
\end{lemma}

\begin{proof}
For $\lambda=\lambda_C\zeta_\lambda$ with $\zeta_\lambda^N=1$,
\[
 \lambda^{Nm+j}
 =\lambda_C^{Nm+j}\zeta_\lambda^{Nm+j}
 =\lambda_C^j(\lambda_C^N)^m\zeta_\lambda^j.
\]
Sum over $\lambda\in C$.
\end{proof}

\begin{lemma}[Minimal annihilator of one polynomial mode]
\label{lem:polynomial-mode-annihilator}
Let $\mu\neq0$ and let $Q(m)$ be a nonzero polynomial of degree $d$.  The sequence $u_m=Q(m)\mu^m$ has minimal shift annihilator
\[
 (E-\mu)^{d+1}.
\]
\end{lemma}

\begin{proof}
A direct calculation gives
\[
 (E-\mu)(Q(m)\mu^m)=\mu^{m+1}\Delta Q(m).
\]
Each application lowers the polynomial degree by one.  Therefore the $(d+1)$st power annihilates the sequence while the $d$th power does not.
\end{proof}

\begin{lemma}[Coprime modal synthesis]
\label{lem:coprime-modal-synthesis}
If $u_m=\sum_{r=1}^sQ_r(m)\mu_r^m$ with distinct nonzero $\mu_r$ and nonzero $Q_r$, then the minimal annihilator of $u$ is
\[
 \prod_{r=1}^s(E-\mu_r)^{1+\deg Q_r}.
\]
\end{lemma}

\begin{proof}
The indicated product annihilates every term.  The primary factors for distinct $\mu_r$ are pairwise coprime.  Applying the product of all factors except the $r$th isolates a nonzero polynomial multiple of $\mu_r^m$, which still requires the full exponent $1+\deg Q_r$ by \cref{lem:polynomial-mode-annihilator}.  No factor or exponent can be removed.
\end{proof}

\begin{proof}[Proof of \cref{thm:spectral-torsion-fourier}]
Statement (1) follows by summing \cref{lem:classwise-transfer} over all torsion classes.  Distinct classes have distinct $N$th powers.  Since substitution $x\mapsto Nm+j$ preserves the degree of a nonzero polynomial in characteristic zero, \cref{lem:polynomial-mode-annihilator,lem:coprime-modal-synthesis} give statements (2) and (3).  The reduced rational denominator is the reciprocal form of the same minimal shift polynomial, proving (5).

For (4), if every class is a singleton then $R_{\{\lambda\},j}=P_\lambda$, so every modal degree survives and $\rho_a(N,j)=\rho(a)$.  Conversely, suppose a class $C$ contains at least two roots.  Its contribution to the global rank is
\[
 \sum_{\lambda\in C}(d_\lambda+1),
\]
while its packet contribution is at most
\[
 1+\max_{\lambda\in C}d_\lambda,
\]
or zero.  The latter is strictly smaller.  Hence every packet has a strict total rank drop.

For (6), every torsion ratio $\zeta$ has order dividing $L_a$.  Therefore
\[
 \zeta^{N+L_a}=1\Longleftrightarrow\zeta^N=1,
\]
so $\sim_{N+L_a}=\sim_N$.  The phase factors also satisfy $\zeta^{j+L_a}=\zeta^j$.  The degree formula in (3) proves both periodicities.  Statement (7) is immediate.
\end{proof}

\subsection{Rank-deficit decomposition}
The theorem identifies two distinct sources of simplification: unavoidable spectral merging and address-dependent Fourier cancellation.

For a torsion class $C$, define
\[
 r_\lambda=d_\lambda+1,
 \qquad
 r_C^{\max}=\max_{\lambda\in C}r_\lambda,
\]
and set
\[
 q_{C,j}=
 \begin{cases}
 1+\deg R_{C,j},&R_{C,j}\not\equiv0,\\
 0,&R_{C,j}\equiv0.
 \end{cases}
\]

\begin{theorem}[Collision--phase deficit decomposition]
\label{thm:deficit-decomposition}
The packet-rank deficit
\[
 \delta_a(N,j)=\rho(a)-\rho_a(N,j)
\]
decomposes exactly as
\[
 \boxed{
 \delta_a(N,j)
 =
 \sum_{C\in\Lambda/{\sim_N}}
 \left[
 \underbrace{\sum_{\lambda\in C}r_\lambda-r_C^{\max}}_{\text{torsion-collision deficit}}
 +
 \underbrace{r_C^{\max}-q_{C,j}}_{\text{phase-cancellation deficit}}
 \right].
 }
\]
The first term depends only on $N$.  The second depends on the packet address and measures cancellation of leading Jordan-polynomial data inside each merged class.
\end{theorem}

\begin{proof}
For every class, add and subtract $r_C^{\max}$ between its global contribution $\sum_{\lambda\in C}r_\lambda$ and its packet contribution $q_{C,j}$, then sum over classes.
\end{proof}

\begin{corollary}[Degenerate sampling moduli]
Let
\[
 \mathcal D_a=
 \left\{N\geq1:\rho_a(N,j)<\rho(a)\text{ for some, equivalently every, }j\right\}.
\]
Then
\[
 \boxed{
 \mathcal D_a
 =
 \bigcup_{\substack{\lambda\neq\eta\\\lambda/\eta\text{ torsion}}}
 \operatorname{ord}(\lambda/\eta)\,\bbN.
 }
\]
In particular, $\mathcal D_a$ is periodic modulo $L_a$.
\end{corollary}

\begin{corollary}[Nondegenerate C-finite sequences]
If no ratio of two distinct active characteristic roots is a root of unity, then $L_a=1$ and
\[
 \boxed{
 \rho_a(N,j)=\rho(a)
 \qquad(N\geq1,\\ j\geq0).
 }
\]
Every arithmetic subsequence preserves the global minimal recurrence order.
\end{corollary}

\subsection{Rational generating functions and root filters}
Let
\[
 A(z)=\sum_{n\geq0}a_nz^n,
 \qquad
 U_{N,j}(w)=\sum_{m\geq0}a_{Nm+j}w^m.
\]
The residue projector and the decimated OGF are related by
\[
 \boxed{
 P_{N,j}A(z)=z^jU_{N,j}(z^N).
 }
\]
If $K$ contains a primitive $N$th root of unity $\omega_N$, then
\[
 \boxed{
 z^jU_{N,j}(z^N)
 =
 \frac1N\sum_{r=0}^{N-1}\omega_N^{-jr}A(\omega_N^rz).
 }
\]
Thus the root filter is a Fourier analysis operator, while the theorem identifies exactly how the poles of the filtered rational function merge and cancel.

\begin{corollary}[Pole-transfer law]
Suppose the reduced OGF of a persistent C-finite sequence has poles $z=\lambda^{-1}$ of orders $d_\lambda+1$.  Under decimation by $N$, poles whose reciprocals have the same $N$th power merge.  The surviving pole at $w=\lambda_C^{-N}$ has exact order $1+\deg R_{C,j}$.
\end{corollary}

\begin{remark}
The denominator obtained from the inherited jump matrix $C^N$ is only an upper bound.  The reduced denominator $Q_{N,j}$ is obtained after both reachability--observability minimization and Fourier cancellation.
\end{remark}

\subsection{Finite torsion signature and quasipolynomial packet counts}
The exact double periodicity compresses infinitely many packet systems into finite data.

\begin{definition}[Torsion packet signature]
The finite table
\[
 \mathfrak T_a(r,s)=\rho_a(N,j),
 \qquad
 N\equiv r\pmod{L_a},
 \quad
 j\equiv s\pmod{L_a},
\]
is called the \emph{torsion packet signature} of $a$.
\end{definition}

\begin{theorem}[Packet-rank counting quasipolynomials]
\label{thm:rank-count-quasipolynomial}
For $q\geq0$, let
\[
 c_q(N)=\#\{0\leq j<N:\rho_a(N,j)=q\}.
\]
For each residue $s\pmod{L_a}$ there are integers $A_{q,s}$ and $B_{q,s}$ such that, whenever $N=tL_a+s$ with $0\leq s<L_a$,
\[
 \boxed{
 c_q(N)=tA_{q,s}+B_{q,s}.
 }
\]
Hence $c_q(N)$ is an exact degree-one quasipolynomial in $N$.
\end{theorem}

\begin{proof}
For fixed $N\bmod L_a$, the rank depends only on $j\bmod L_a$.  The interval $0\leq j<N$ contains $t$ complete address periods followed by a prefix of length $s$.  Count addresses of rank $q$ in the full period and in the prefix.
\end{proof}

\begin{corollary}
The number of packets of every fixed rank, including rank zero, is an exact quasipolynomial function of the modulus.  Once $N\geq L_a$, the minimum packet rank, maximum packet rank, and set of attained packet ranks depend only on $N\bmod L_a$ and are therefore eventually periodic.
\end{corollary}

\begin{figure}[htbp]
\centering
\includegraphics[width=.86\textwidth]{17_torsion_packet_signature.pdf}
\caption{A complete $6\times6$ torsion packet signature for a mixed order-two/order-three example.  One finite table determines every arithmetic-subsequence rank.}
\label{fig:torsion-signature}
\end{figure}

\subsection{Fourier information conservation}
A single packet can lose modes, but a full address cycle retains the original classwise data.

Let $T$ be a maximal torsion class.  Choose $\lambda_0\in T$, let $L_T$ be the exponent of the finite ratio group
\[
 \{\lambda/\lambda_0:\lambda\in T\}\subseteq\mu_{L_T},
\]
and write
\[
 \lambda=\lambda_0\zeta,
 \qquad
 P_\zeta=P_\lambda,
\]
with $P_\zeta=0$ for unused $\zeta\in\mu_{L_T}$.

\begin{theorem}[Classwise Fourier analysis and synthesis]
\label{thm:fourier-conservation}
Whenever $L_T\mid N$, all modes in $T$ belong to one $N$-torsion class and their phase polynomials are
\[
 \boxed{
 R_j(x)=\sum_{\zeta\in\mu_{L_T}}\zeta^jP_\zeta(x),
 \qquad0\leq j<L_T.
 }
\]
They reconstruct the original modal polynomials by the inverse finite Fourier transform
\[
 \boxed{
 P_\zeta(x)=\frac1{L_T}\sum_{j=0}^{L_T-1}\zeta^{-j}R_j(x).
 }
\]
Therefore the complete packet family conserves all global torsion-mode information even when individual packets have smaller rank.
\end{theorem}

\begin{proof}
The first formula is the definition of the packet-phase polynomial.  Orthogonality of characters of the cyclic group gives
\[
 \frac1{L_T}\sum_{j=0}^{L_T-1}(\xi/\zeta)^j
 =\begin{cases}1,&\xi=\zeta,\\0,&\xi\neq\zeta.\end{cases}
\]
Substitute the forward transform and interchange the finite sums.
\end{proof}

\begin{figure}[htbp]
\centering
\includegraphics[width=.92\textwidth]{18_fourier_mode_reconstruction.pdf}
\caption{Classwise information conservation.  Sampling merges torsion-related global modes, packet addresses retain their Fourier phases, and a complete address cycle reconstructs the original modal polynomials.}
\label{fig:fourier-reconstruction}
\end{figure}

\subsection{A packet uncertainty principle}
Fourier cancellation cannot suppress a coefficient layer on arbitrarily many packet addresses.  In this subsection assume the finite modal coefficient field is embedded in $\bbC$, as it is for integer and rational sequences.

Fix one torsion class of exponent $L$.  For a polynomial degree $d$, let $c_\zeta$ be the coefficient of $x^d$ in $P_\zeta(x)$.  Define
\[
 s_d=\#\{\zeta\in\mu_L:c_\zeta\neq0\}
\]
and let $t_d$ be the number of addresses $j\pmod L$ for which the coefficient of $x^d$ in $R_j(x)$ is nonzero.

\begin{theorem}[Packet Fourier uncertainty]
\label{thm:packet-uncertainty}
If the coefficient vector $(c_\zeta)$ is nonzero, then
\[
 \boxed{s_dt_d\geq L.}
\]
In particular, if the highest active degree is supported on $s$ global torsion modes, then at least $\lceil L/s\rceil$ packet addresses retain that highest degree.
\end{theorem}

\begin{proof}
Let
\[
 \widehat c_j=\sum_{\zeta\in\mu_L}c_\zeta\zeta^j.
\]
By Parseval for the unnormalized Fourier transform,
\[
 \sum_j|\widehat c_j|^2=L\sum_\zeta|c_\zeta|^2.
\]
For every $j$, Cauchy--Schwarz gives
\[
 |\widehat c_j|^2
 \leq
 s_d\sum_\zeta|c_\zeta|^2.
\]
Only $t_d$ terms on the left are nonzero, so
\[
 L\sum_\zeta|c_\zeta|^2
 \leq
 t_ds_d\sum_\zeta|c_\zeta|^2.
\]
Cancel the nonzero norm.
\end{proof}

\subsection{Basis-free state-space interpretation}
Let $(C,v,\ell)$ be the minimal reachable-observable realization of $a$.  For each packet define
\[
 J_N=C^N,
 \qquad
 v_{N,j}=C^jv.
\]
The inherited packet realization is $(J_N,v_{N,j},\ell)$.

\begin{corollary}[Exact reachable-observable dimension of the jump system]
The dimension of the reachable-observable reduction of $(J_N,v_{N,j},\ell)$ is exactly $\rho_a(N,j)$ from \cref{thm:spectral-torsion-fourier}.  Thus the torsion-phase formula computes the result of state-space minimization without constructing an infinite Hankel matrix.
\end{corollary}

\begin{remark}[Two exact computational routes]
There are two complementary packet compilers.
\begin{enumerate}
\item The \emph{spectral route} factors the minimal recurrence, detects torsion ratios, forms $R_{C,j}$, and outputs the reduced recurrence directly.
\item The \emph{matrix route} computes $J_N=C^N$, forms $v_{N,j}$, performs exact reachable-observable reduction, and extracts the minimal polynomial on the reduced state.
\end{enumerate}
The matrix route avoids algebraic factorization for one packet.  The spectral route exposes the finite signature and all moduli simultaneously.
\end{remark}

\subsection{Proof-producing transfer algorithm}

\begin{algorithm}[Spectral torsion packet compiler]
\label{alg:spectral-torsion-compiler}
Input a persistent C-finite sequence through a reduced rational OGF, a minimal recurrence, or a minimal linear realization.
\begin{enumerate}
\item Compute the active roots $\lambda$ and multiplicity degrees $d_\lambda$ over a splitting field.
\item Recover the active modal polynomials $P_\lambda(n)$ from the initial values or Jordan data.
\item Determine every root-of-unity ratio $\lambda/\eta$ and its order; form $L_a$.
\item For every $N\bmod L_a$, partition $\Lambda$ by $\lambda^N$.
\item For every $j\bmod L_a$, form the phase polynomials $R_{C,j}$.
\item Emit $\chi_{N,j}$, $Q_{N,j}$, $\rho_a(N,j)$, the deficit decomposition, and the finite signature.
\item Certify the result by modal substitution or by exact reachable-observable reduction of selected jump systems.
\end{enumerate}
\end{algorithm}

\begin{theorem}[Compiler completeness on the persistent C-finite stratum]
The algorithm computes the exact minimal recurrence, reduced rational denominator, Hankel rank, and finite torsion signature for every arithmetic subsequence of the input sequence.  No packet modulus outside one period $N\bmod L_a$ and no packet address outside one period $j\bmod L_a$ requires a new rank calculation.
\end{theorem}

\begin{proof}
Each output is one of the explicit quantities in \cref{thm:spectral-torsion-fourier,thm:deficit-decomposition}.  Exact periodicity reduces the infinite family to the finite signature.
\end{proof}

\subsection{Anchor examples}

\begin{example}[Fibonacci has no packet-rank drop]
Let
\[
 F_n=\frac{\alpha^n-\beta^n}{\alpha-\beta},
 \qquad
 \alpha=\frac{1+\sqrt5}{2},
 \quad
 \beta=\frac{1-\sqrt5}{2}.
\]
The ratio $\alpha/\beta$ is not a root of unity.  Hence $L_F=1$ and
\[
 \boxed{\rho_F(N,j)=2\quad\text{for all }N,j.}
\]
The exact packet polynomial is
\[
 (X-\alpha^N)(X-\beta^N)
 =X^2-L_NX+(-1)^N,
\]
which recovers the universal Fibonacci packet law and proves its order is minimal for every packet.
\end{example}

\begin{example}[Complete mode cancellation]
For
\[
 a_n=2^n+(-2)^n,
\]
the torsion period is $L_a=2$.  If $N$ is odd, the two modes remain distinct and every packet has rank two.  If $N$ is even, they merge and
\[
 R_j=1+(-1)^j.
\]
Therefore
\[
 \boxed{
 \rho_a(N,j)=
 \begin{cases}
 2,&N\text{ odd},\\
 1,&N\text{ even and }j\text{ even},\\
 0,&N\text{ even and }j\text{ odd}.
 \end{cases}}
\]
\end{example}

\begin{example}[Cancellation of the leading Jordan layer]
Let
\[
 a_n=n2^n+(n+1)(-2)^n.
\]
The global minimal polynomial is $(X-2)^2(X+2)^2$, so $\rho(a)=4$.  For even $N$,
\[
 R_j(x)=x+(-1)^j(x+1).
\]
Hence
\[
 R_j(x)=
 \begin{cases}
 2x+1,&j\text{ even},\\
 -1,&j\text{ odd}.
 \end{cases}
\]
It follows that
\[
 \boxed{
 \rho_a(N,j)=
 \begin{cases}
 4,&N\text{ odd},\\
 2,&N\text{ even and }j\text{ even},\\
 1,&N\text{ even and }j\text{ odd}.
 \end{cases}}
\]
The odd packet address cancels the entire degree-one Jordan layer but leaves a constant mode.
\end{example}

\begin{example}[A cubic periodic selector]
Let $\omega^3=1$, $\omega\neq1$, and
\[
 a_n=1+\omega^n+\omega^{2n}
 =\begin{cases}3,&3\mid n,\\0,&3\nmid n.\end{cases}
\]
The global rank is three and $L_a=3$.  If $3\nmid N$, all modes remain distinct.  If $3\mid N$, all modes merge and
\[
 R_j=1+\omega^j+\omega^{2j}.
\]
Therefore
\[
 \boxed{
 \rho_a(N,j)=
 \begin{cases}
 3,&3\nmid N,\\
 1,&3\mid N\text{ and }3\mid j,\\
 0,&3\mid N\text{ and }3\nmid j.
 \end{cases}}
\]
\end{example}

\begin{example}[Mixed torsion orders and a complete finite signature]
Consider the integer sequence
\[
 a_n=2^n+(-2)^n+3^n(\omega^n+\omega^{2n}),
 \qquad\omega^3=1.
\]
The active roots split into an order-two torsion class $\{2,-2\}$ and an order-three torsion class $\{3\omega,3\omega^2\}$.  Thus $L_a=6$ and the global rank is four.  The exact rank table is
\[
 \boxed{
 \rho_a(N,j)=
 \begin{cases}
 4,&N\equiv1,5\pmod6,\\
 3,&N\equiv3\pmod6,\\
 3,&N\equiv2,4\pmod6\text{ and }j\text{ even},\\
 2,&N\equiv2,4\pmod6\text{ and }j\text{ odd},\\
 2,&N\equiv0\pmod6\text{ and }j\text{ even},\\
 1,&N\equiv0\pmod6\text{ and }j\text{ odd}.
 \end{cases}}
\]
The entire infinite packet-rank family is therefore represented by the $6\times6$ table in \cref{fig:torsion-signature}.
\end{example}

\begin{figure}[htbp]
\centering
\includegraphics[width=.94\textwidth]{19_packet_rank_profiles.pdf}
\caption{Representative packet-rank profiles.  Torsion-free spectra preserve rank, while order-two and order-three torsion produce exact periodic rank collapse and address-dependent cancellation.}
\label{fig:packet-rank-profiles}
\end{figure}
\FloatBarrier

\subsection{Finite-support transients}
\label{subsec:finite-transient-extension}
The persistence hypothesis removes only nilpotent state components.

\begin{theorem}[Finite-transient extension]
Every C-finite sequence has a unique decomposition
\[
 \boxed{a=h+b,}
\]
where $h$ has finite support and $b$ is persistent C-finite.  Equivalently, its rational OGF has a unique decomposition
\[
 \boxed{A(z)=H(z)+\frac{P(z)}{Q(z)},}
\]
where $H(z)$ is a polynomial, $P/Q$ is reduced and proper, $Q(0)=1$, and the recurrence associated with $P/Q$ has nonzero constant term.

For every $N,j$, the packet of $b$ is governed by \cref{thm:spectral-torsion-fourier}.  The packet of $h$ is an explicit finite-support correction.  Consequently the persistent packet signature is exactly periodic, and all deviations caused by zero characteristic roots occur in a finite computable prefix.
\end{theorem}

\begin{proof}
In a minimal state realization, split the state space into the generalized zero-eigenspace and the direct sum of nonzero generalized eigenspaces.  The zero-eigenspace is nilpotent and therefore contributes a finite-support sequence.  The remaining operator is invertible and gives $b$.  Uniqueness follows from the uniqueness of the primary decomposition.
\end{proof}

\subsection{Transfer-theorem status and scope}
The theorem resolves the exact packet-rank and periodicity questions for the complete C-finite stratum.  It does not claim the same formula for P-recursive, algebraic, automatic, or nonlinear fixed-point sequences.  In those strata, the analogous target is to identify the native global invariant whose restriction under packet refinement controls exact minimal complexity.

The significance of the result is therefore methodological as well as linear-algebraic:
\[
 \boxed{
 \text{global invariant}
 \longrightarrow
 \text{exact packet transfer law}
 \longrightarrow
 \text{finite classification of infinitely many refinements}.
 }
\]
This is the transfer pattern pursued in the later holonomic, nonlinear, transform, and local-arithmetic parts of the calculus.

% ===== END sections/06_spectral_torsion_fourier.tex =====

% ===== BEGIN sections/06_holonomic_hypergeometric.tex =====
\part{Holonomic, Digital, and Transform Realizations}
\section{Holonomic, hypergeometric, and Ap\'ery-type systems}

A P-recursive sequence need not have finite Hankel rank.  It nevertheless has a finite recurrence grammar because the transfer law depends polynomially on the index.

\subsection{P-recursive and D-finite equivalence}
\begin{definition}[P-recursive sequence]
A sequence $a_n$ over a characteristic-zero field is P-recursive when there exist polynomials $p_0(n),\ldots,p_r(n)$, not all zero, such that
\[
 \sum_{k=0}^{r}p_k(n)a_{n+k}=0
\]
for all sufficiently large $n$.
\end{definition}
\begin{definition}[D-finite series]
A formal series $A(z)$ is D-finite when it satisfies
\[
 \sum_{k=0}^{r}q_k(z)A^{(k)}(z)=0
\]
with polynomial coefficients $q_k(z)$, not all zero.
\end{definition}

\begin{classical}[Coefficient translation]
Over characteristic zero, a sequence is P-recursive exactly when its ordinary generating series is D-finite.  The translation is effective in both directions by comparing coefficients after replacing $n$ with the Euler operator $\theta=z\frac{d}{dz}$.
\end{classical}

\begin{proposition}[Finite nonautonomous state]
Assume $p_r(n)\neq0$ on the range of interest.  A P-recursive sequence has a state
\[
 v_n=(a_n,a_{n+1},\ldots,a_{n+r-1})^T
\]
with rational transfer
\[
 v_{n+1}=C(n)v_n.
\]
After clearing denominators, exact integer or modular evolution is possible whenever the leading coefficient is invertible in the chosen channel.
\end{proposition}

\subsection{Multisection closure}
\begin{theorem}[Holonomic residue closure]
If $A(z)$ is D-finite, then every residue section $P_{N,j}A(z)$ is D-finite.  Equivalently, every residue subsequence $a_{Nm+j}$ of a P-recursive sequence is P-recursive.
\end{theorem}
\begin{proof}
Over a field containing a primitive $N$th root of unity,
\[
 P_{N,j}A(z)=\frac1N\sum_{r=0}^{N-1}\omega_N^{-jr}A(\omega_N^rz).
\]
D-finite series are closed under algebraic substitution $z\mapsto\omega_N^rz$ and finite linear combinations.  The result descends to the base field because the section has base-field coefficients.
\end{proof}

\begin{remark}
The theorem proves existence of a packet recurrence.  It does not by itself give the smallest recurrence.  Elimination in the differential module, direct creative telescoping, or exact recurrence guessing followed by proof can produce a smaller packet operator.
\end{remark}

\begin{proposition}[Order bound by direct state jumping]
A first-order rational state system $v_{n+1}=C(n)v_n$ of dimension $r$ induces, on the $j$th $N$-section,
\[
 v_{N(m+1)+j}=J_{N,j}(m)v_{Nm+j},
\qquad
 J_{N,j}(m)=\prod_{s=0}^{N-1}C(Nm+j+s),
\]
where the product is ordered from right to left in increasing $s$.  Every scalar output is therefore P-recursive of order at most $r$ after denominator clearing.
\end{proposition}
\begin{proof}
Compose the $N$ consecutive transfers.  Eliminate the finite state by the cyclic-vector or determinant construction over the rational function field $K(m)$.
\end{proof}

\subsection{Hypergeometric seeds}
A hypergeometric term satisfies
\[
 \frac{a_{n+1}}{a_n}=R(n)
\]
for a rational function $R$.  It is a one-dimensional nonautonomous seed system.

\begin{proposition}[Hypergeometric packet ratio]
For $u_m=a_{Nm+j}$,
\[
 \boxed{\frac{u_{m+1}}{u_m}=\prod_{s=0}^{N-1}R(Nm+j+s).}
\]
Thus every residue subsequence of a hypergeometric term is hypergeometric in the packet index.
\end{proposition}
\begin{proof}
Telescope the $N$ successive ratios.
\end{proof}

Examples include factorials, binomial coefficients along affine rays, rising and falling factorials, central binomial coefficients, and many classical combinatorial kernels.  For example,
\[
 c_n=\binom{2n}{n},\qquad
 \frac{c_{n+1}}{c_n}=\frac{2(2n+1)}{n+1}.
\]
The Catalan numbers $C_n=c_n/(n+1)$ are also hypergeometric, although their nonlinear tree realization in \cref{sec:nonlinear} exposes a different structure.

\subsection{D-finite closure inside the grammar}
\begin{theorem}[Holonomic closure block]
The D-finite class is closed under addition, Cauchy product, differentiation, integration, Hadamard product, algebraic substitution into the argument, and multiplication by an algebraic series.  Each closure operation can be implemented as a finite module construction over $K(z)$.
\end{theorem}
\begin{proof}[Proof scheme]
Represent each series by the finite-dimensional $K(z)$-space spanned by its derivatives.  Direct sums handle addition, tensor products handle products, and the chain rule handles algebraic substitution.  Hadamard closure follows through the diagonal of a rationally constructed multivariate D-finite series.  The constructions yield annihilating operators by finite linear dependence.
\end{proof}

\begin{boundary}
D-finite series are generally not closed under unrestricted composition or unit reciprocal.  The sequence grammar permits those operations because it is broader than the holonomic subclass, but the output may leave that subclass.
\end{boundary}

\subsection{Ap\'ery-type recurrences}
Define the Ap\'ery numbers
\[
 A_n=\sum_{k=0}^{n}\binom nk^2\binom{n+k}{k}^2.
\]
They begin
\[
 1,5,73,1445,33001,819005,21460825,584307365,\ldots.
\]

\begin{theorem}[Ap\'ery recurrence]
The sequence satisfies
\[
 \boxed{(n+1)^3A_{n+1}=(34n^3+51n^2+27n+5)A_n-n^3A_{n-1}}
\]
for $n\geq1$, with $A_0=1$ and $A_1=5$.
\end{theorem}
\begin{proof}[Certificate route]
The binomial sum is a proper hypergeometric summand in $(n,k)$.  Creative telescoping produces the displayed recurrence with a telescoping certificate in $k$.  Summing over the finite support eliminates the certificate boundary.  Alternatively, the recurrence plus the two initial values is a complete independent sequence specification.
\end{proof}

The corresponding state is
\[
 \begin{pmatrix}A_{n+1}\\A_n\end{pmatrix}
 =
 \begin{pmatrix}
 \dfrac{34n^3+51n^2+27n+5}{(n+1)^3}&-\dfrac{n^3}{(n+1)^3}\\[1em]
 1&0
 \end{pmatrix}
 \begin{pmatrix}A_n\\A_{n-1}\end{pmatrix}.
\]
Residue packets are obtained by multiplying $N$ consecutive rational transfer matrices.

\begin{proposition}[Integrality as a separate invariant]
A rational state recurrence can generate integers even when its one-step matrix has denominators.  Integrality is not implied merely by the state realization; it requires an independent divisibility certificate, a combinatorial sum, a lattice-preservation proof, or an arithmetic recurrence argument.
\end{proposition}

This distinction is important for the compiler: it records both a semantic recurrence certificate and, when available, a coefficient-domain certificate.

\subsection{Diagonal and constant-term realizations}
Many integer sequences arise as diagonals or constant terms of rational functions.  A typical form is
\[
 a_n=[x_1^0\cdots x_d^0]R(x_1,\ldots,x_d)^n.
\]
This is a finite seed specification: multiplication by $R$ is the evolution operator and constant-term extraction is the output projector.  The resulting generating function is often D-finite, and residue packets correspond to replacing the step operator by its $N$th power while shifting the seed.

\begin{principle}[Multiple realizations are informative]
The same sequence may admit a hypergeometric sum, a diagonal, a differential equation, a recurrence, and a nonlinear combinatorial grammar.  These are not redundant.  Their comparison reveals different invariants: integrality, asymptotics, congruences, packet complexity, or positivity.
\end{principle}

\subsection{Exact modular execution}
For a P-recursive recurrence, modular evaluation at a prime $p$ is direct when the leading coefficient is nonzero modulo $p$.  Singular indices require one of three treatments:
\begin{enumerate}
\item use an equivalent recurrence whose leading coefficient is invertible;
\item carry a fraction-free lattice state before reduction;
\item treat the singular index as a refinement boundary and restart from exact values.
\end{enumerate}
The typed presentation must declare which treatment is valid.  A modular program that divides by a nonunit without such a declaration is not a proof-producing realization.

% ===== END sections/06_holonomic_hypergeometric.tex =====

% ===== BEGIN sections/07_automatic_sequences.tex =====
\section{Automatic, regular, and digit-path sequences}

Residue refinement becomes especially concrete for digit-defined sequences.  Instead of evolving one index at a time, a finite family of digit operators evolves through the radix expansion of the index.

\subsection{Finite digit realizations}
Fix a base $k\geq2$.  Write
\[
 n=(d_\ell\cdots d_1d_0)_k.
\]
\begin{definition}[Digit seed system]
A finite linear digit realization is
\[
 \boxed{a_n=u^TT_{d_\ell}\cdots T_{d_1}T_{d_0}v,}
\]
where $T_0,\ldots,T_{k-1}$ are fixed matrices and a consistent convention for leading zeros is declared.
\end{definition}
A deterministic finite automaton with output is the special case in which each $T_d$ is a zero-one transition matrix and $u$ reads the output state.

\begin{definition}[$k$-kernel]
The $k$-kernel of $a$ is
\[
 \mathcal K_k(a)=\{(a_{k^en+r})_{n\geq0}:e\geq0,\ 0\leq r<k^e\}.
\]
A finite-valued sequence is $k$-automatic when its $k$-kernel is finite.  A sequence is $k$-regular over a ring when the module spanned by its $k$-kernel is finitely generated.
\end{definition}

\begin{theorem}[Digit refinement theorem]
If $a$ has a digit realization, then every radix packet $n\mapsto a_{k^en+r}$ has the same digit operators and a refined seed obtained by the fixed digits of $r$.  Conversely, a finitely generated $k$-kernel produces a finite linear digit realization.
\end{theorem}
\begin{proof}
Appending the digits of $r$ to the expansion of $n$ multiplies the seed by a fixed word in the $T_d$.  Conversely, choose generators for the kernel module; the maps $b(n)\mapsto b(kn+d)$ act linearly on that finite module and become the digit matrices.
\end{proof}

\begin{corollary}[Exact radix packet compiler]
Compilation to modulus $k^e$ requires no elimination: the packet state is obtained by following the $e$ fixed low-order digit operators.  Refinement from $k^e$ to $k^{e+1}$ is one additional digit transition.
\end{corollary}

This is a literal projector refinement tree.  A node at depth $e$ represents one residue class modulo $k^e$; its $k$ children are the compatible classes modulo $k^{e+1}$.

\subsection{Thue--Morse}
Let $t_n$ be the parity of the number of ones in the binary expansion of $n$.  Then
\[
 \boxed{t_{2n}=t_n,\qquad t_{2n+1}=1-t_n,\qquad t_0=0.}
\]
With the affine state $(t_n,1)^T$, take
\[
 T_0=\begin{pmatrix}1&0\\0&1\end{pmatrix},\qquad
 T_1=\begin{pmatrix}-1&1\\0&1\end{pmatrix}.
\]
Reading binary digits from most significant to least significant gives an exact two-dimensional linear digit realization.

\begin{theorem}[Thue--Morse packet closure]
Every binary residue packet of Thue--Morse is either $t_n$ or $1-t_n$.  Hence the entire infinite binary refinement tree has only two semantic states.
\end{theorem}
\begin{proof}
Induct on the number of fixed low-order bits using the two recursions.  Each appended zero preserves the state and each appended one complements it.
\end{proof}

The verification suite checks the bit-count definition, both binary recursions, and all values through $4095$.

\subsection{Rudin--Shapiro and pattern counters}
Let $r_n=(-1)^{e_2(n)}$, where $e_2(n)$ counts overlapping occurrences of $11$ in the binary expansion.  A finite automaton needs only remember the preceding digit and the current parity.  The corresponding four-state digit system generates the Rudin--Shapiro sequence.  More generally, counting occurrences of a fixed finite digit word modulo $m$ requires only a suffix automaton crossed with a cyclic count state.

\begin{proposition}[Finite-pattern realization]
Let $W$ be a finite set of radix words and let the output depend on the number of occurrences of each word modulo a fixed integer.  Then the resulting sequence is automatic.
\end{proposition}
\begin{proof}
The state records the longest suffix that is a prefix of some word in $W$, together with the residue counters.  There are finitely many states and each new digit updates them deterministically.
\end{proof}

\subsection{Paperfolding, period-doubling, and substitutions}
Uniform substitutions induce digit operators on the finite set of letters.  If a fixed point of a $k$-uniform morphism is read by a finite coding, the resulting sequence is $k$-automatic.  Period-doubling, regular paperfolding, and many morphic parity sequences therefore have finite seed grammars.

\begin{construction}[Substitution-to-digit compiler]
For a $k$-uniform substitution $\sigma$, label each position $0,\ldots,k-1$ inside $\sigma(x)$.  The digit $d$ sends letter $x$ to the $d$th letter of $\sigma(x)$.  These $k$ maps are the digit operators.  The seed is the initial letter and the output map is the coding.
\end{construction}

\subsection{From automatic to regular}
Finite-valued automatic sequences are restrictive, but $k$-regular systems allow integer-valued outputs.  The binary sum-of-digits function $s_2(n)$ satisfies
\[
 s_2(2n)=s_2(n),\qquad s_2(2n+1)=s_2(n)+1.
\]
It is generated by the affine matrices
\[
 T_0=\begin{pmatrix}1&0\\0&1\end{pmatrix},\qquad
 T_1=\begin{pmatrix}1&1\\0&1\end{pmatrix}
\]
on $(s_2(n),1)^T$.  Similar systems generate digit sums, weighted digit sums, and many radix valuation statistics.

\begin{proposition}[Affine-to-linear lift]
Any finite-dimensional affine digit system $x\mapsto A_dx+b_d$ becomes linear after adjoining a constant coordinate:
\[
 \begin{pmatrix}x\\1\end{pmatrix}
 \longmapsto
 \begin{pmatrix}A_d&b_d\\0&1\end{pmatrix}
 \begin{pmatrix}x\\1\end{pmatrix}.
\]
\end{proposition}

\subsection{Base dependence and mixed-radix packets}
Automaticity is base-dependent.  A sequence can have a compact binary grammar and no comparable decimal grammar.  Projector packets themselves are base-neutral, but radix refinement chooses a preferred tree.

If two bases are multiplicatively independent, simultaneous compact realizations are highly restrictive.  The sequence calculus records this as an intersection of two digit-operator systems rather than silently treating all residue refinements as equivalent.

\begin{boundary}
A finite program that computes $a_n$ from the digits of $n$ is not necessarily automatic or regular.  The state dimension must remain finite and independent of the digit length.  Primality indicators, for example, are computable from digits but are not thereby finite-state digit sequences.
\end{boundary}

\subsection{Digit paths and ordinary residue projectors}
For an arbitrary modulus $N$ not equal to a radix power, one can combine the digit automaton with a finite residue tracker
\[
 r\longmapsto kr+d\pmod N.
\]
The product state generates both $a_n$ and the indicator of $n\equiv j\pmod N$.  This gives an explicit finite realization for every ordinary residue packet of an automatic sequence.

\begin{theorem}[Automatic residue packets]
If $a_n$ is $k$-automatic, then for every $N$ and $j$, the masked sequence
\[
 b_n=a_n\mathbf 1_{n\equiv j\,(N)}
\]
is $k$-automatic.  If $a_n$ is $k$-regular, then $b_n$ is $k$-regular.
\end{theorem}
\begin{proof}
Take the direct product of the original digit system with the finite residue automaton and apply the masked output map.
\end{proof}

\subsection{Certificate format}
A digital-sequence certificate consists of:
\begin{enumerate}
\item the radix and digit-reading convention;
\item the finite state set or matrices $T_d$;
\item the seed and output map;
\item an invariant proving leading-zero consistency;
\item equations identifying the refined packet states;
\item optional minimization by partition refinement or reachable-observable reduction.
\end{enumerate}
Unlike term matching, this certificate proves every coefficient simultaneously.

% ===== END sections/07_automatic_sequences.tex =====

% ===== BEGIN sections/08_transform_calculus.tex =====
\section{Transform calculus and transported packet systems}

Sequence databases are dense with transforms: binomial, inverse binomial, Stirling, Euler, Borel, Appell, Sheffer, Riordan, Hadamard, and compositions of these.  The projector-seed perspective treats an invertible transform as a change of coordinates and transports the packet projectors by conjugation.

\subsection{Triangular operators}
Let $T=(T_{n,k})_{n,k\geq0}$ be lower triangular with unit or invertible diagonal.  It acts on sequences by
\[
 (Ta)_n=\sum_{k=0}^{n}T_{n,k}a_k.
\]
Every requested prefix uses a finite matrix block.

\begin{theorem}[Unit-triangular inversion]
If $T=D+L$ with invertible diagonal $D$ and strictly lower triangular $L$, then on every finite prefix
\[
 \boxed{T^{-1}=\sum_{r\geq0}(-D^{-1}L)^rD^{-1},}
\]
where the sum truncates because a strictly lower triangular matrix is nilpotent on a finite block.
\end{theorem}
\begin{proof}
This is the finite Neumann expansion of $(I+D^{-1}L)^{-1}D^{-1}$ on each prefix.  The compatible finite inverses define the formal triangular inverse.
\end{proof}

\subsection{Riordan operators}
For ordinary generating functions define
\[
 \mathcal R_{g,f}[A](z)=g(z)A(f(z)),
\]
where $g(0)$ is a unit, $f(0)=0$, and $f'(0)$ is a unit.  Its matrix entries are
\[
 R_{n,k}=[z^n]g(z)f(z)^k.
\]

\begin{theorem}[Riordan group law]
\[
 \boxed{\mathcal R_{g,f}\mathcal R_{h,\ell}=\mathcal R_{g\,h\circ f,\,\ell\circ f}.}
\]
The inverse is
\[
 \boxed{\mathcal R_{g,f}^{-1}=\mathcal R_{1/(g\circ\bar f),\,\bar f},}
\]
where $\bar f$ is the compositional inverse of $f$.
\end{theorem}
\begin{proof}
Substitute directly:
\[
 g(z)h(f(z))A(\ell(f(z))).
\]
The inverse follows by composing with $\bar f$ and cancelling $g$.
\end{proof}

For exponential generating functions the same formula is interpreted in the divided-power basis.  This contains Appell and Sheffer transformations.

\subsection{Transported projectors}
\begin{theorem}[Transform transport]
Let $T$ be an invertible linear transform and $P_{N,j}$ the residue projectors.  Define
\[
 \boxed{Q_{N,j}=TP_{N,j}T^{-1}.}
\]
Then
\[
 Q_{N,i}Q_{N,j}=\delta_{ij}Q_{N,j},\qquad
 \sum_jQ_{N,j}=I,
\]
and for $b=Ta$,
\[
 \boxed{b=\sum_{j=0}^{N-1}Q_{N,j}b.}
\]
\end{theorem}
\begin{proof}
Conjugation preserves products, sums, and identities.  Moreover $Q_{N,j}b=TP_{N,j}a$.
\end{proof}

Thus an invertible transform does not destroy the packet decomposition; it changes what a packet looks like in transformed coordinates.  The original residue packet is sparse in coefficient space, while the transported packet may be a dense structured triangular combination.

\begin{corollary}[Refinement under transforms]
If $N\mid M$, then
\[
 Q_{N,j}=\sum_{r\equiv j\,(N)}Q_{M,r}.
\]
\end{corollary}

\begin{figure}[htbp]
\centering
\includegraphics[width=.9\textwidth]{06_transform_orbit_network.pdf}
\caption{A transform orbit of sequence presentations.  Invertible arrows carry packet systems by conjugation; noninvertible constructors remain typed exact maps.}
\end{figure}

\subsection{Binomial transform}
The binomial transform is
\[
 b_n=\sum_{k=0}^{n}\binom nk a_k,
\]
with inverse
\[
 a_n=\sum_{k=0}^{n}(-1)^{n-k}\binom nk b_k.
\]
For EGFs,
\[
 \widehat B(t)=e^t\widehat A(t).
\]
For OGFs,
\[
 B(z)=\frac{1}{1-z}A\!\left(\frac{z}{1-z}\right).
\]
It is therefore both an Appell multiplication in the divided-power channel and a Riordan operator in the ordinary channel.

\begin{figure}[htbp]
\centering
\includegraphics[width=.69\textwidth]{10_binomial_transform_matrix.pdf}
\caption{The lower-triangular binomial transform.  Every finite prefix is an exact finite coordinate change.}
\end{figure}

\begin{example}[Binomial transform of Fibonacci]
Let $b_n=\sum_k\binom nkF_k$.  Since the Fibonacci characteristic roots $\alpha,\beta$ are shifted to $1+\alpha,1+\beta$, the transformed sequence satisfies
\[
 b_{n+2}=3b_{n+1}-b_n,
\]
with $b_0=0,b_1=1$.  Hence $b_n=F_{2n}$.  The transform has converted one linear seed system into a similar two-dimensional system with a different evolution operator.
\end{example}

\subsection{Stirling transforms and Bell families}
The second-kind Stirling transform
\[
 b_n=\sum_{k=0}^{n}\left\{\begin{matrix}n\\k\end{matrix}\right\}a_k
\]
has EGF relation
\[
 \widehat B(t)=\widehat A(e^t-1).
\]
The first-kind transform is the inverse composition with $\log(1+t)$.  Taking $a_k=x^k$ gives the Touchard polynomials
\[
 T_n(x)=\sum_{k=0}^{n}\left\{\begin{matrix}n\\k\end{matrix}\right\}x^k,
\qquad
 \sum_{n\geq0}T_n(x)\frac{t^n}{n!}=e^{x(e^t-1)}.
\]
At $x=1$, $T_n(1)=B_n$ is the Bell sequence.

\subsection{Appell and Sheffer orbits}
An Appell family has EGF
\[
 G(t)e^{xt}=\sum_{n\geq0}P_n(x)\frac{t^n}{n!},
\]
with $G(0)$ a unit.  It obeys
\[
 \frac{d}{dx}P_n(x)=nP_{n-1}(x).
\]
A Sheffer family replaces $xt$ by $xF(t)$ with $F(0)=0$ and $F'(0)$ a unit.  The pair $(G,F)$ is an exponential Riordan operator.

\begin{theorem}[Appell seed reconstruction]
Let $g_n=P_n(0)$.  Then
\[
 \boxed{P_n(x)=\sum_{k=0}^{n}\binom nk g_{n-k}x^k.}
\]
Hence an Appell family is reconstructed from one scalar seed sequence by the binomial transform with the monomial seed.
\end{theorem}
\begin{proof}
Multiply the EGFs $G(t)$ and $e^{xt}$ and compare divided-power coefficients.
\end{proof}

Euler, Bernoulli, Genocchi, Hermite-type, and many generalized polynomial sequences appear as transform orbits of their scalar boundary sequence.

\subsection{Transform equivalence and invariants}
\begin{definition}[Transform orbit]
Fix a declared transform group $\mathfrak T$.  Two sequence presentations are transform-equivalent when $b=Ta$ for some $T\in\mathfrak T$.
\end{definition}
This relation is only useful after the group is restricted.  Allowing every invertible triangular matrix would put all unit-normalized sequences in one orbit and destroy content.

Meaningful restricted groups include:
\begin{itemize}
\item affine binomial transforms;
\item ordinary or exponential Riordan groups with algebraic or rational data;
\item Appell and Sheffer subgroups;
\item finite Fourier and residue coordinate changes;
\item local unit-triangular matrices with prescribed valuation weights.
\end{itemize}

Possible orbit invariants include differential transcendence class, singularity pattern, growth order, denominator spectrum, Hankel-rank behavior, packet complexity, and local valuation slopes.  Not all are preserved by every transform class, so the transform type is part of the theorem statement.

\subsection{Weighted triangular paths}
Suppose
\[
 T=D+\sum_{r\geq1}\pi^{w_r}L_r,
\]
where every $L_r$ lowers a filtration and $w_r>0$.  Expanding the inverse produces weighted operator paths.  A path of cumulative weight at least $M$ vanishes modulo $\pi^M$.

\begin{theorem}[Valuation-truncated triangular inverse]
If every off-diagonal path has additive nonnegative weight and only finitely many paths have weight below $M$ between fixed coefficient levels, then $T^{-1}$ modulo $\pi^M$ is a finite sum over paths of cumulative weight $<M$.
\end{theorem}
\begin{proof}
Expand the filtered Neumann series.  Every omitted term is divisible by $\pi^M$, while local nilpotence or coefficientwise filtration makes the retained set finite.
\end{proof}

This theorem is the abstract mechanism behind the odd-prime and weighted dyadic Euler truncations developed later.

\subsection{Transform diagrams}
A sequence identity is often best represented by a commuting diagram rather than one formula.  If $T$ is a transform, $P$ a projector, and $F$ an evolution operator, the compiler records
\[
 \begin{CD}
 a @>F>> a'\\
 @VT VV @VV T V\\
 b @>>TFT^{-1}> b'
 \end{CD}
 \qquad
 \begin{CD}
 a @>P>> a_P\\
 @VT VV @VV T V\\
 b @>>TPT^{-1}> b_P.
 \end{CD}
\]
These diagrams distinguish a true conjugacy from a mere numerical coincidence between initial terms.

% ===== END sections/08_transform_calculus.tex =====


% ===== BEGIN sections/09_nonlinear_composition_trees.tex =====
\part{Nonlinear, Product, and Arithmetic Sequence Systems}
\section{Composition, Bell partitions, and rooted-tree fixed points}\label{sec:nonlinear}

Finite structural generation extends far beyond linear recurrences.  Composition produces set partitions, polynomial fixed points produce rooted trees, and packet refinement labels every combinatorial constructor by residue address.

\subsection{Packet Fa\`a di Bruno calculus}
Let $F(y)=\sum_{r\geq0}f_r y^r/r!$ and let $G(z)$ be a formal series.  The classical derivative formula is
\[
 D^n(F\circ G)
 =
 \sum_{\pi\in\Pi_n}
 F^{(|\pi|)}(G)
 \prod_{B\in\pi}G^{(|B|)},
\]
where $\Pi_n$ is the set of partitions of $\{1,\ldots,n\}$.

\begin{theorem}[Packet Fa\`a di Bruno formula]
Let $G=\sum_{j=0}^{N-1}G_j$ be its exact residue decomposition.  Then every output packet of $F\circ G$ is obtained by projecting
\[
 \boxed{
 \sum_{\pi\in\Pi_n}
 F^{(|\pi|)}(G)
 \prod_{B\in\pi}
 \left(\sum_{j=0}^{N-1}G_j^{(|B|)}\right).
 }
\]
Equivalently, expand over packet-labeled set partitions and retain exactly those monomials whose total degree address equals the requested packet.
\end{theorem}
\begin{proof}
Insert the exact reconstruction $G=\sum_jG_j$ into the ordinary Fa\`a di Bruno formula, distribute the finite sums, and apply the residue projector.  No approximation is introduced.
\end{proof}

\begin{remark}
The Bell numbers appear twice: first as the number of set partitions indexing the derivative formula, and again as a coefficient sequence generated by a particular composition.  The two appearances are the same combinatorial mechanism viewed as operator syntax and as coefficient semantics.
\end{remark}

\subsection{Bell and Touchard families}
Define
\[
 \mathcal B(t)=\exp(e^t-1)=\sum_{n\geq0}B_n\frac{t^n}{n!}.
\]

\begin{theorem}[Bell realization]
The Bell numbers admit each of the following finite projector-seed presentations:
\begin{enumerate}[label=\textup{(\roman*)}]
\item the composition seed $\mathcal B=\exp\circ(e^t-1)$;
\item the differential system
\[
 E'=E,\qquad \mathcal B'=E\mathcal B,
 \qquad E(0)=\mathcal B(0)=1;
\]
\item the recurrence
\[
 \boxed{B_{n+1}=\sum_{k=0}^{n}\binom nk B_k,\qquad B_0=1;}
\]
\item the partition expansion $B_n=|\Pi_n|$.
\end{enumerate}
These presentations are equivalent under coefficient extraction and Fa\`a di Bruno expansion.
\end{theorem}
\begin{proof}
Differentiate $\mathcal B=e^{e^t-1}$ to obtain $\mathcal B'=e^t\mathcal B$.  Comparing divided-power coefficients gives the recurrence.  In Fa\`a di Bruno, all derivatives of $e^y$ at $y=0$ and of $e^t-1$ at $t=0$ equal one, so each set partition contributes one.
\end{proof}

The Touchard polynomials
\[
 T_n(x)=\sum_{k=0}^{n}\left\{\begin{matrix}n\\k\end{matrix}\right\}x^k
\]
have EGF
\[
 \sum_{n\geq0}T_n(x)\frac{t^n}{n!}=e^{x(e^t-1)}.
\]
They refine Bell numbers by the number of blocks, with $T_n(1)=B_n$.

\begin{proposition}[Bell packet differential system]
Write $E_r=P_{N,r}e^t$ and $B_s=P_{N,s}\mathcal B$.  Then
\[
 \boxed{(B_{k+1})'=\sum_{r+s\equiv k\,(N)}E_rB_s.}
\]
Together with $\sum_sB_s=\mathcal B$, this is an exact finite system for all Bell residue packets.
\end{proposition}
\begin{proof}
Project $\mathcal B'=E\mathcal B$ and use the derivative and cyclic-product packet laws.
\end{proof}

\begin{figure}[htbp]
\centering
\includegraphics[width=.8\textwidth]{07_bell_catalan_growth.pdf}
\caption{Bell and Catalan growth.  Their grammars are both finite, although one is compositional and the other is algebraic-tree based.}
\end{figure}

\subsection{Contractive polynomial systems and trees}
Let $R$ be a commutative ring and let $U(z)\in zR[[z]]^d$ satisfy
\[
 U=G(z,U),
\]
where every monomial of $G$ either has positive $z$-order or contains an unknown factor in a way that raises the coefficient filtration.

\begin{theorem}[Rooted-tree expansion]
Under the formal contraction hypothesis of the master realization theorem, the unique solution of $U=G(z,U)$ is the coefficientwise limit of iteration from zero.  Expanding the iteration produces a sum over rooted trees whose vertices are labeled by the monomials of $G$ and whose leaves are labeled by explicit seed terms.
\end{theorem}
\begin{proof}
Each substitution replaces an unresolved occurrence of $U$ by one monomial of $G$.  A finite iteration produces finite rooted expansion trees.  For a fixed coefficient degree, contraction bounds the required depth, so only finitely many trees contribute.  Their stable sum is the unique fixed point.
\end{proof}

\begin{corollary}[Packet-labeled tree expansion]
After residue compilation, every edge and vertex receives a packet address.  A tree contributes to packet $j$ exactly when the sum of its degree-address contributions is $j$ modulo $N$.
\end{corollary}

\subsection{Catalan}
Let
\[
 u=z+u^2,
 \qquad u(0)=0.
\]
Then
\[
 u(z)=\frac{1-\sqrt{1-4z}}2
 =\sum_{n\geq1}C_{n-1}z^n.
\]

\begin{theorem}[Catalan tree realization]
The coefficient $[z^n]u$ is the number of full ordered binary trees with $n$ leaves, hence
\[
 \boxed{C_n=\frac{1}{n+1}\binom{2n}{n},\qquad
 C_{n+1}=\sum_{k=0}^{n}C_kC_{n-k}.}
\]
\end{theorem}
\begin{proof}
The term $z$ is a leaf and the term $u^2$ joins an ordered pair of subtrees.  Coefficient extraction gives the convolution recurrence.  Lagrange inversion applied to $u=z/(1-u)$ gives the binomial formula.
\end{proof}

For the parity split $u=u_0+u_1$, where $u_j=P_{2,j}u$,
\[
 \boxed{u_0=u_0^2+u_1^2,\qquad u_1=z+2u_0u_1.}
\]

\begin{proof}
The square has even part $u_0^2+u_1^2$ and odd part $2u_0u_1$; the seed $z$ is odd.
\end{proof}

\begin{figure}[htbp]
\centering
\includegraphics[width=.86\textwidth]{08_catalan_even_odd_packets.pdf}
\caption{The two exact Catalan coefficient packets.  Their coupled fixed-point equations reconstruct the unsplit tree series.}
\end{figure}

\subsection{Lagrange inversion as a sequence compiler}
\begin{theorem}[One-variable Lagrange seed theorem]
Let $u=z\Phi(u)$ with $\Phi(0)$ a unit.  Then for $n\geq1$,
\[
 \boxed{[z^n]u(z)^k=\frac{k}{n}[w^{n-k}]\Phi(w)^n.}
\]
In particular,
\[
 [z^n]u(z)=\frac1n[w^{n-1}]\Phi(w)^n.
\]
\end{theorem}
\begin{proof}
Apply the formal residue identity to the change of variables $z=u/\Phi(u)$.  The coefficient formula follows from integration by parts in the formal residue calculus.
\end{proof}

This theorem compiles a nonlinear implicit seed into explicit coefficient extraction.  For $\Phi(w)=(1+w)^r$ it gives Fuss--Catalan numbers
\[
 \boxed{C_n^{(r)}=\frac{1}{rn+1}\binom{(r+1)n}{n}}
\]
for the solution of $C=1+zC^{r+1}$.

\subsection{Motzkin and Schr\"oder systems}
The Motzkin generating function is the unique unit solution of
\[
 M=1+zM+z^2M^2.
\]
Coefficient extraction gives
\[
 M_0=1,\qquad
 M_n=M_{n-1}+\sum_{k=0}^{n-2}M_kM_{n-2-k}.
\]
It begins
\[
 1,1,2,4,9,21,51,127,323,835,2188,\ldots.
\]

The large Schr\"oder series satisfies
\[
 S=1+zS+zS^2
\]
and begins
\[
 1,2,6,22,90,394,1806,\ldots.
\]
Both are finite polynomial fixed-point grammars, and every residue packet is obtained by cyclic packet convolution.

\begin{proposition}[Algebraic packet system]
Suppose $A_1,\ldots,A_d$ satisfy polynomial equations
\[
 F_i(z,A_1,\ldots,A_d)=0
\]
with a unique formal branch at the declared initial values.  For every modulus $N$, replacing each product by cyclic packet product and each series by its $N$ packets gives a finite polynomial system whose synthesis is the original branch.
\end{proposition}
\begin{proof}
This is the exact packet compiler applied to polynomial maps.  Uniqueness of the original branch and invertibility of analysis-synthesis identify the packet branch.
\end{proof}

\subsection{Combinatorial constructor grammar}
A large collection of labeled combinatorial classes is generated from a finite constructor set:
\[
 \operatorname{SEQ},\quad \operatorname{SET},\quad \operatorname{CYC},\quad
 +,\quad\times,\quad\circ,
\]
with the corresponding OGF or EGF operations.  When the constructor specification is finite and analytically or formally well founded, it is a finite projector-seed grammar.

\begin{principle}[Constructor-to-packet translation]
Disjoint union becomes packetwise addition, product becomes cyclic packet convolution, substitution becomes packet Fa\`a di Bruno, and recursive constructors become packet-labeled tree equations.  Thus residue refinement can be performed before or after combinatorial synthesis.
\end{principle}

\subsection{Ambiguity of multiple nonlinear specifications}
A nonlinear equation may have several formal or analytic branches.  The presentation must include sufficient initial data, a valuation condition, or a selected local branch.  For example, $u=z+u^2$ has two algebraic solutions, but only one lies in $zR[[z]]$.  Branch selection is part of the seed, not an informal afterthought.

\begin{boundary}
A finite polynomial equation does not guarantee integer coefficients, positivity, or a combinatorial interpretation.  Those are additional invariants requiring separate proofs.  The realization theorem guarantees exact semantics, not every desired arithmetic property.
\end{boundary}

% ===== END sections/09_nonlinear_composition_trees.tex =====

% ===== BEGIN sections/10_filtered_products.tex =====
\section{Filtered products, Euler transforms, and partition families}

Infinite products are legitimate finite-grammar objects when each coefficient depends on only finitely many factors.  This coefficient stabilization is the completion principle behind partitions and many product-defined sequences.

\subsection{Coefficientwise completion}
Let $R[[z]]$ carry the $z$-adic filtration.

\begin{theorem}[Filtered product completion]
Let $F_m(z)\in1+z^{w(m)}R[[z]]$ with $w(m)\to\infty$.  Then
\[
 \boxed{F(z)=\prod_{m\geq1}F_m(z)}
\]
exists uniquely in $R[[z]]$.  For every degree $n$, the coefficient $[z^n]F$ is already equal to the coefficient in the finite product over those $m$ with $w(m)\leq n$.
\end{theorem}
\begin{proof}
Factors with $w(m)>n$ are congruent to one modulo $z^{n+1}$ and cannot affect degrees through $n$.  Hence the finite products form a Cauchy system whose coefficients stabilize.
\end{proof}

\begin{corollary}[Exact finite-prefix algorithm]
To compute $F$ through degree $n$, retain only factors with $w(m)\leq n$ and truncate every intermediate product modulo $z^{n+1}$.
\end{corollary}

\subsection{Euler transform}
Let $(b_m)_{m\geq1}$ be integers or elements of a characteristic-zero ring and define
\[
 A(z)=\prod_{m\geq1}(1-z^m)^{-b_m}
 =\sum_{n\geq0}a_nz^n,
 \qquad a_0=1.
\]
Set
\[
 c_n=\sum_{d\mid n}d b_d.
\]

\begin{theorem}[Euler-transform recurrence]
The coefficients satisfy
\[
 \boxed{n a_n=\sum_{k=1}^{n}c_k a_{n-k}.}
\]
Conversely, if $a_0=1$, then the sequence $b_n$ is recovered by
\[
 \boxed{b_n=\frac1n\sum_{d\mid n}\mu(n/d)c_d,}
\]
where $c_n=[z^n]zA'(z)/A(z)$.
\end{theorem}
\begin{proof}
Logarithmic differentiation gives
\[
 \frac{zA'}A
 =\sum_{m\geq1}b_m\frac{mz^m}{1-z^m}
 =\sum_{n\geq1}\left(\sum_{d\mid n}d b_d\right)z^n.
\]
Multiplying by $A$ yields the coefficient recurrence.  M\"obius inversion recovers $n b_n$ from $c_n$.
\end{proof}

\begin{remark}
The product and recurrence are two exact seed presentations.  The product exposes combinatorial components and coefficient stabilization; the recurrence gives a direct sequential evaluator.
\end{remark}

\subsection{Ordinary partitions}
For unrestricted partitions, $b_m=1$:
\[
 P(z)=\prod_{m\geq1}\frac1{1-z^m}
 =\sum_{n\geq0}p(n)z^n.
\]
Then $c_n=\sigma_1(n)$ and
\[
 \boxed{n p(n)=\sum_{k=1}^{n}\sigma_1(k)p(n-k).}
\]
The first values are
\[
 1,1,2,3,5,7,11,15,22,30,42,56,77,101,135,176,\ldots.
\]

\begin{figure}[htbp]
\centering
\includegraphics[width=.91\textwidth]{09_partition_product_stabilization.pdf}
\caption{Coefficient stabilization for the partition product.  The coefficient of degree $n$ is final once factors through $m=n$ have been included.}
\end{figure}

\subsection{Distinct and odd parts}
The distinct-part series is
\[
 D(z)=\prod_{m\geq1}(1+z^m).
\]
Since
\[
 1+z^m=\frac{1-z^{2m}}{1-z^m},
\]
we obtain
\[
 D(z)=\prod_{m\text{ odd}}(1-z^m)^{-1}.
\]
Thus the number of partitions into distinct parts equals the number of partitions into odd parts.  In the Euler-transform exponent sequence, $b_m=1$ for odd $m$ and $b_m=0$ for even $m$.

\subsection{Colored, over-, and plane partitions}
For $r$-colored partitions,
\[
 P_r(z)=\prod_{m\geq1}(1-z^m)^{-r},
\]
so $b_m=r$ and $c_n=r\sigma_1(n)$.

Overpartitions have
\[
 \overline P(z)=\prod_{m\geq1}\frac{1+z^m}{1-z^m},
\]
which is again a finitely ruled product.

Plane partitions have MacMahon's product
\[
 \prod_{m\geq1}(1-z^m)^{-m},
\]
corresponding to $b_m=m$ and
\[
 c_n=\sum_{d\mid n}d^2=\sigma_2(n).
\]
Hence their coefficients $q(n)$ satisfy
\[
 n q(n)=\sum_{k=1}^{n}\sigma_2(k)q(n-k).
\]

\subsection{Products with signs and multiple species}
A general finite-rule product may have factors
\[
 (1-\varepsilon_m z^{w_m})^{-b_m},
 \qquad \varepsilon_m\in R.
\]
As long as $w_m\to\infty$, coefficient stabilization remains exact.  The logarithmic derivative becomes
\[
 \frac{zA'}A
 =\sum_m b_m w_m\frac{\varepsilon_m z^{w_m}}{1-\varepsilon_m z^{w_m}},
\]
which compiles into a divisor-like recurrence whenever the weights are effectively enumerable.

\subsection{Packet products and stabilization}
\begin{theorem}[Residue packet limit]
Let $F^{(M)}=\prod_{m\leq M}F_m$ converge coefficientwise to $F$.  Then for every $N,j$,
\[
 \boxed{P_{N,j}F=\lim_{M\to\infty}P_{N,j}F^{(M)}}
\]
coefficientwise, and the packet coefficient of degree $n$ stabilizes at the same finite stage as the unsplit coefficient.
\end{theorem}
\begin{proof}
Residue projectors are coefficient selectors and therefore continuous in the coefficientwise topology.
\end{proof}

The packet system can be evaluated by maintaining an $N$-tuple and multiplying each factor through cyclic packet convolution.  This may expose congruence cancellations that are invisible in the scalar product.

\subsection{Pentagonal acceleration as a secondary relation}
The reciprocal partition product also satisfies Euler's pentagonal identity
\[
 \prod_{m\geq1}(1-z^m)
 =\sum_{r\in\bbZ}(-1)^r z^{r(3r-1)/2}.
\]
Multiplying by $P(z)$ gives the sparse recurrence
\[
 p(n)=\sum_{r\neq0}(-1)^{r+1}p\!\left(n-\frac{r(3r-1)}2\right),
\]
with out-of-range terms zero.  The product grammar and sparse recurrence are two different exact compilations of the same object.

\begin{principle}[Redundant exact presentations]
When a sequence has both a dense logarithmic-derivative recurrence and a sparse functional recurrence, the compiler should retain both.  One may be better for proof transfer, another for computation, and a third for local congruences.
\end{principle}

\subsection{Partition-type asymptotics are not automatic}
The filtered product theorem proves existence and exact finite-prefix computation.  It does not by itself prove asymptotic formulas, modularity, or congruence families.  Such results require additional analytic or algebraic structures.  The seed calculus records these as extra operators and relations rather than treating them as consequences of formal completion alone.

\begin{boundary}
An infinite product is admitted only with an effective factor rule and a proof of coefficient stabilization.  Writing an arbitrary sequence as a product whose exponents encode all its values is storage universality, not a compact structural realization.
\end{boundary}

% ===== END sections/10_filtered_products.tex =====

% ===== BEGIN sections/11_arithmetic_functions.tex =====
\section{Multiplicative arithmetic functions and prime-power seeds}

Ordinary coefficient indices are additive under Cauchy product.  Arithmetic functions carry a second monoidal structure in which indices multiply and the natural product is Dirichlet convolution.

\subsection{Dirichlet sequence objects}
For arithmetic functions $f,g:\bbN_{\geq1}\to R$, define
\[
 (f\Dir g)(n)=\sum_{d\mid n}f(d)g(n/d).
\]
The identity is $\varepsilon(1)=1$ and $\varepsilon(n)=0$ for $n>1$.

\begin{theorem}[Dirichlet inverse]
If $f(1)$ is a unit, there is a unique $g=f^{-1_D}$ with $f\Dir g=\varepsilon$.  It is determined by
\[
 g(1)=f(1)^{-1},\qquad
 \boxed{g(n)=-f(1)^{-1}\sum_{\substack{d\mid n\\d<n}}f(n/d)g(d).}
\]
\end{theorem}
\begin{proof}
The coefficient at $n$ in $f\Dir g=\varepsilon$ is triangular with diagonal coefficient $f(1)$.
\end{proof}

The Dirichlet series channel, when convergent or formally interpreted prime by prime, sends convolution to multiplication:
\[
 \mathcal D_f(s)=\sum_{n\geq1}\frac{f(n)}{n^s},
 \qquad
 \mathcal D_{f\Dir g}=\mathcal D_f\mathcal D_g.
\]

\subsection{Prime-power reconstruction}
\begin{definition}[Multiplicative seed]
An arithmetic function $f$ is multiplicative when $f(1)=1$ and
\[
 f(mn)=f(m)f(n)\qquad((m,n)=1).
\]
Its local seeds are the sequences $f(p^e)$ indexed by primes $p$ and exponents $e\geq0$.
\end{definition}

\begin{theorem}[Prime-power seed theorem]
A multiplicative arithmetic function is reconstructed exactly by
\[
 \boxed{f(n)=\prod_{p^e\parallel n}f(p^e).}
\]
If the map $(p,e)\mapsto f(p^e)$ has a finite rule, then $f$ has a finite projector-seed grammar consisting of factorization, local evaluation, and finite product synthesis.
\end{theorem}
\begin{proof}
The prime-power factors of $n$ are pairwise coprime.  Repeated multiplicativity gives the product.
\end{proof}

This theorem separates a finite structural rule from the cost of integer factorization.  The semantic presentation can be compact even when the fastest evaluation algorithm is nontrivial.

\subsection{Classical local seeds}
The following table gives local rules.
\begin{center}
\begin{tabular}{L{.2\textwidth}L{.31\textwidth}L{.37\textwidth}}
\toprule
Function&Prime-power seed&Reconstruction feature\\
\midrule
M\"obius $\mu$&$1,-1,0,0,\ldots$&detects squarefreeness\\
Liouville $\lambda$&$(-1)^e$&parity of total valuation\\
Divisor count $\tau_r$&$\binom{e+r-1}{r-1}$&ordered $r$-fold divisor convolution\\
Sum of divisors $\sigma_k$&$1+p^k+\cdots+p^{ek}$&finite geometric local state\\
Euler totient $\varphi$&$1$ at $e=0$, $p^{e-1}(p-1)$ for $e\geq1$&unit count\\
Jordan $J_k$&$p^{ek}(1-p^{-k})$ for $e\geq1$&higher-dimensional unit count\\
Squarefree indicator $\mu^2$&$1,1,0,0,\ldots$&bounded local exponent\\
\bottomrule
\end{tabular}
\end{center}

Many of these local exponent sequences are C-finite or polynomial-exponential, so the arithmetic layer nests the linear sequence layer inside each prime packet.

\subsection{Euler products}
For multiplicative $f$,
\[
 \mathcal D_f(s)=\prod_p\left(\sum_{e\geq0}f(p^e)p^{-es}\right)
\]
whenever analytic convergence is valid, or as a formal product in independent prime variables.  Thus the same object has:
\begin{itemize}
\item a global coefficient sequence $f(n)$;
\item local prime-power seed sequences;
\item a Dirichlet-convolution algebra;
\item an Euler-product completion.
\end{itemize}

\begin{proposition}[Convolution preserves multiplicativity]
If $f$ and $g$ are multiplicative, then $f\Dir g$ is multiplicative, and at each prime
\[
 \boxed{(f\Dir g)(p^e)=\sum_{r=0}^{e}f(p^r)g(p^{e-r}).}
\]
\end{proposition}
\begin{proof}
Divisors of a product of coprime integers split uniquely.  At a prime power, divisors are $p^r$.
\end{proof}

Therefore Dirichlet convolution becomes ordinary Cauchy convolution in the exponent coordinate of each prime.  This is a local-global polyphase decomposition of the arithmetic index.

\subsection{Residue packets for multiplicative convolution}
Define
\[
 f_r(n)=f(n)\mathbf1_{n\equiv r\,(N)}.
\]

\begin{theorem}[Multiplicative residue packet law]
For every residue $k$ modulo $N$,
\[
 \boxed{P_{N,k}(f\Dir g)
 =\sum_{rs\equiv k\,(N)} f_r\Dir g_s.}
\]
\end{theorem}
\begin{proof}
In a divisor pair $n=d(n/d)$, the residue of $n$ is the product of the two divisor residues modulo $N$.  Group the divisor sum by those residues.
\end{proof}

Unlike Cauchy packets, whose addresses add, Dirichlet packet addresses multiply.  The address algebra is the finite multiplicative monoid $\bbZ/N\bbZ$, including zero divisors.  Restriction to units gives the finite unit group.

\subsection{Valuation packets}
For a fixed prime $p$, decompose positive integers by $e=v_p(n)$ and unit part $u=n/p^e$.  A multiplicative function factors as
\[
 f(n)=f(p^e)f(u).
\]
This is a two-coordinate seed system: an additive valuation coordinate and a unit-residue coordinate.  Refining $u$ modulo $p^L$ produces the local unit packets used in $p$-adic interpolation.

\begin{proposition}[Valuation convolution]
At a fixed prime, Dirichlet convolution of multiplicative functions becomes Cauchy convolution of the local exponent sequences:
\[
 h_e=\sum_{r=0}^{e}f_r g_{e-r}.
\]
The generating series
\[
 F_p(T)=\sum_{e\geq0}f(p^e)T^e
\]
therefore satisfies $H_p(T)=F_p(T)G_p(T)$.
\end{proposition}

\subsection{Completely multiplicative modes}
If $f$ is completely multiplicative, then $f(p^e)=f(p)^e$, a one-dimensional geometric local state.  Finite linear combinations of such modes behave like exponential sums on the multiplicative monoid.  Their finite differences along unit-group steps provide valuation gains, as developed in the Mahler section.

\subsection{Additive arithmetic functions}
An additive arithmetic function satisfies $a(mn)=a(m)+a(n)$ for coprime $m,n$ and reconstructs as
\[
 a(n)=\sum_{p^e\parallel n}a(p^e).
\]
Examples include $\omega(n)$, $\Omega(n)$, and $\log n$.  These use additive synthesis over prime-power packets rather than multiplicative synthesis.

\subsection{Hybrid ordinary and Dirichlet transforms}
A sequence may participate in both Cauchy and Dirichlet convolution.  For example, the divisor-sum sequence has an OGF, a Dirichlet series, an Euler product, and recurrence-transform descriptions.  There is no general algebra isomorphism between the two convolution structures.  A bridge theorem must state the extra kernel or transform that connects them.

\begin{boundary}
Prime-power reconstruction is exact for multiplicative functions, but it does not imply that arbitrary arithmetic functions are determined by prime powers.  Nor does a compact local rule automatically give a fast global evaluator without access to the factorization of $n$.
\end{boundary}

% ===== END sections/11_arithmetic_functions.tex =====

% ===== BEGIN sections/12_special_function_families.tex =====
\section{Exponential sections, reciprocal towers, and congruential Euler families}

Residue projectors applied to the exponential produce a universal finite differential system.  Signed sections solve $y^{(N)}=-y$; reciprocals produce generalized secant and tangent families; normalized ordinary sections produce a triangular family containing Euler and Bernoulli boundaries.

\subsection{Exponential residue sections}
For $0\leq j<N$, define
\[
 \boxed{\mathcal E_{N,j}(z)=\sum_{n\geq0}\frac{z^{Nn+j}}{(Nn+j)!}.}
\]
These are the residue packets of $e^z$.

\begin{theorem}[Root-filter representation]
If $\omega_N$ is a primitive $N$th root of unity, then
\[
 \boxed{\mathcal E_{N,j}(z)=\frac1N\sum_{r=0}^{N-1}\omega_N^{-jr}e^{\omega_N^rz}.}
\]
Moreover,
\[
 \mathcal E_{N,j}'=\mathcal E_{N,j-1}\quad(j>0),
 \qquad
 \mathcal E_{N,0}'=\mathcal E_{N,N-1},
\]
and every section satisfies $\mathcal E_{N,j}^{(N)}=\mathcal E_{N,j}$.
\end{theorem}
\begin{proof}
The root average selects exactly the degrees congruent to $j$ modulo $N$.  Termwise differentiation cycles the sections.
\end{proof}

Thus the $N$ packets form a fundamental system for the constant-coefficient equation $y^{(N)}=y$.

\subsection{Signed transfer basis}
Define
\[
 \boxed{\mathcal D_{N,j}(z)=\sum_{n\geq0}(-1)^n\frac{z^{Nn+j}}{(Nn+j)!},}
 \qquad
 \mathcal D_N=\mathcal D_{N,0}.
\]

\begin{theorem}[Signed derivative cycle]
For $1\leq j<N$,
\[
 \mathcal D_{N,j}'=\mathcal D_{N,j-1},
\]
while
\[
 \mathcal D_{N,0}'=-\mathcal D_{N,N-1}.
\]
Hence
\[
 \boxed{\mathcal D_{N,j}^{(N)}=-\mathcal D_{N,j}.}
\]
If $\mathcal R_N=\{\eta:\eta^N=-1\}$, then
\[
 \boxed{\mathcal D_N(z)=\frac1N\sum_{\eta\in\mathcal R_N}e^{\eta z}.}
\]
\end{theorem}
\begin{proof}
Differentiate the series and reindex.  One complete cycle introduces one sign.  The root average selects multiples of $N$ and contributes $\eta^{Nn}=(-1)^n$.
\end{proof}

The first cases are
\[
 \mathcal D_1=e^{-z},\qquad
 \mathcal D_2=\cos z,\qquad
 \mathcal D_4=\cos(z/\sqrt2)\cosh(z/\sqrt2).
\]

\subsection{Reciprocal secant tower}
Formally at the origin, or analytically away from zeros of $\mathcal D_N$, define
\[
 \boxed{\operatorname{Sec}_N^{[r]}(z)=\mathcal D_N(z)^{-r}}
\]
and
\[
 \operatorname{Sec}_N^{[r]}(z)=\sum_{n\geq0}P_{N,n}(r)\frac{z^{Nn}}{(Nn)!}.
\]

\begin{theorem}[Reciprocal semigroup]
\[
 \operatorname{Sec}_N^{[r]}\operatorname{Sec}_N^{[s]}
 =\operatorname{Sec}_N^{[r+s]},
\]
so
\[
 \boxed{P_{N,n}(r+s)=\sum_{q=0}^{n}\binom{Nn}{Nq}P_{N,q}(r)P_{N,n-q}(s).}
\]
For fixed $n$, $P_{N,n}(r)$ is a polynomial in $r$ of degree at most $n$.
\end{theorem}
\begin{proof}
The functional identity is the exponent law for the same formal logarithm.  Coefficient extraction gives block convolution.  The logarithmic differential equation
\[
 \mathcal D_N S'+r\mathcal D_N'S=0
\]
is triangular in block degree and proves the polynomial degree by induction.
\end{proof}

At $N=2$ and $r=1$, the coefficients are the positive secant numbers
\[
 1,1,5,61,1385,50521,\ldots.
\]

\subsection{Transfer tangent and higher Riccati law}
Define
\[
 \boxed{\operatorname{Tan}_N(z)=-\frac{\mathcal D_N'(z)}{\mathcal D_N(z)}.}
\]
Then
\[
 \frac{d}{dz}\operatorname{Sec}_N^{[r]}
 =r\operatorname{Tan}_N\operatorname{Sec}_N^{[r]}.
\]

Let $Y_N$ be the complete differential Bell polynomial determined by
\[
 \frac{d^N}{dz^N}e^{g(z)}
 =e^{g(z)}Y_N(g',g'',\ldots,g^{(N)}).
\]

\begin{theorem}[Higher Riccati equation]
\[
 \boxed{Y_N(-\operatorname{Tan}_N,-\operatorname{Tan}_N',\ldots,-\operatorname{Tan}_N^{(N-1)})=-1.}
\]
\end{theorem}
\begin{proof}
Write $\mathcal D_N=e^g$ locally.  Then $g'=-\operatorname{Tan}_N$.  Divide $\mathcal D_N^{(N)}=-\mathcal D_N$ by $\mathcal D_N$ and use the Bell-polynomial identity.
\end{proof}

For $N=2$, $\operatorname{Tan}_2=\tan z$ and the equation is
\[
 \operatorname{Tan}_2'=1+\operatorname{Tan}_2^2.
\]

\subsection{Congruential Euler packets}
The exponential section $\mathcal E_{N,j}$ begins with $z^j/j!$.  Remove that forced zero and invert the remaining unit.

\begin{definition}[Congruential Euler function]
\[
 \boxed{\mathcal C_{N,j}(z)=\frac{z^j/j!}{\mathcal E_{N,j}(z)}
 =\sum_{n\geq0}C_{N,j}(n)\frac{z^{Nn}}{(Nn)!}.}
\]
\end{definition}

\begin{theorem}[Triangular inverse system]
The coefficients are uniquely determined by $C_{N,j}(0)=1$ and
\[
 \boxed{\sum_{q=0}^{n}\binom{Nn+j}{Nq}C_{N,j}(q)=\delta_{n0}.}
\]
For $n\geq1$,
\[
 \boxed{C_{N,j}(n)=
 -\frac{\sum_{q=0}^{n-1}\binom{Nn+j}{Nq}C_{N,j}(q)}{\binom{Nn+j}{Nn}}.}
\]
\end{theorem}
\begin{proof}
Multiply $\mathcal C_{N,j}$ by $\mathcal E_{N,j}$ and extract the coefficient of $z^{Nn+j}/(Nn+j)!$.
\end{proof}

\begin{proposition}[Universal denominator clearing]
Let
\[
 D_{N,j}(n)=\prod_{r=1}^{n}\binom{Nr+j}{Nr}.
\]
Then
\[
 \boxed{D_{N,j}(n)C_{N,j}(n)\in\bbZ.}
\]
\end{proposition}
\begin{proof}
Induct through the triangular recurrence.  The product contains every prior diagonal denominator and the new one.
\end{proof}

\subsection{Euler, Bernoulli, and Genocchi boundaries}
At $(N,j)=(2,0)$,
\[
 \mathcal C_{2,0}(z)=\sech z
 =\sum_{n\geq0}E_n\frac{z^n}{n!},
\]
with even values
\[
 1,-1,5,-61,1385,-50521,\ldots.
\]
At $(N,j)=(2,1)$,
\[
 \mathcal C_{2,1}(z)=\frac{z}{\sinh z},
\]
and
\[
 \boxed{C_{2,1}(n)=2(1-2^{2n-1})B_{2n}.}
\]
Multiplication by simple exponential or scalar factors yields the standard Bernoulli and Genocchi generating functions.  The value of the construction is not that these boundary families are new, but that they are embedded in one $(N,j)$ triangular inverse family.

\subsection{Higher powers and block-Appell systems}
For integer $r\geq0$, define
\[
 \mathcal C_{N,j}^{[r]}(z)=\mathcal C_{N,j}(z)^r
 =\sum_{n\geq0}C_{N,j}^{[r]}(n)\frac{z^{Nn}}{(Nn)!}.
\]
Then
\[
 \boxed{C_{N,j}^{[r+s]}(n)=\sum_{q=0}^{n}\binom{Nn}{Nq}C_{N,j}^{[r]}(q)C_{N,j}^{[s]}(n-q).}
\]

Define polynomials by
\[
 \mathcal C_{N,j}(z)^r e^{xz}
 =\sum_{m\geq0}E_m^{(N,j;r)}(x)\frac{z^m}{m!}.
\]

\begin{theorem}[Block-Appell realization]
\[
 \boxed{E_m^{(N,j;r)}(x)=\sum_{0\leq n\leq m/N}\binom{m}{Nn}C_{N,j}^{[r]}(n)x^{m-Nn},}
\]
with
\[
 \frac{d}{dx}E_m^{(N,j;r)}=mE_{m-1}^{(N,j;r)}
\]
and
\[
 E_m^{(N,j;r)}(x+y)=\sum_{q=0}^{m}\binom mqE_q^{(N,j;r)}(x)y^{m-q}.
\]
\end{theorem}
\begin{proof}
Multiply the block EGF by $e^{xz}$ and compare coefficients.  Differentiation and translation follow from the exponential factor.
\end{proof}

\subsection{Coefficient asymptotics from denominator zeros}
If the nearest zeros of $\mathcal D_N$ are simple and isolated from the next zero modulus, residues at those zeros give the leading coefficient asymptotics of $\mathcal D_N^{-r}$.  This is a standard singularity mechanism, now parameterized by $N$ and $r$.

\begin{boundary}
The formal reciprocal families and their coefficient recurrences are exact at the origin.  Global analytic continuation, zero multiplicities, and derivative-weighted zero-zeta sums require separate analytic hypotheses and are not consequences of formal sequence realization alone.
\end{boundary}

\subsection{Why this layer matters for integer sequences}
The analytic seeds supply more than generating functions.  They produce:
\begin{enumerate}
\item finite differential systems;
\item block coefficient recurrences;
\item parameter-polynomial convolution semigroups;
\item generalized Euler inverse matrices;
\item Appell polynomial orbits;
\item local arithmetic coordinate systems after prime-dependent refactorization.
\end{enumerate}
The next section develops that local arithmetic layer explicitly.

% ===== END sections/12_special_function_families.tex =====

% ===== BEGIN sections/13_local_arithmetic.tex =====
\part{Local, Adelic, and Polyphase Transfer}
\section{Odd-prime and dyadic local transfer calculus}

The local arithmetic layer separates two tasks that are often conflated. A residue address records where a coefficient lies. A valuation weight records how accurately that coefficient is visible in a local lattice. Chinese remainder factorization controls the first task. Odd-prime, dyadic, Mahler, and transfer-kernel analysis control the second.

The principal structural point is that the odd primes and the dyadic prime share one convolutional reconstruction law but not one integral coordinate system. At an odd prime the torsion part of the unit group has order $p-1$. At $p=2$ the torsion part is only the sign packet and the principal-unit filtration begins one level deeper. The resulting transfer kernels have different exact valuation profiles.

\subsection{Valuation-filtered triangular systems}
Let $R$ be a complete valuation ring with uniformizer $\pi$. Consider a unit lower-triangular coordinate change
\[
 S_n=\sum_{m=0}^{n}u_{n,m}A_m,
 \qquad u_{n,n}\in R^\times,
\]
with inverse
\[
 A_n=\sum_{m=0}^{n}v_{n,m}S_m.
\]

\begin{theorem}[Uniform filtered truncation]
Assume that a nondecreasing weight function $w:\bbN\to\bbN$ satisfies
\[
 v_{n,m}\in \pi^{w(n-m)}R
 \qquad(0\leq m\leq n).
\]
Then the depth-$M$ truncation
\[
 A_n^{\langle M\rangle}
 =\sum_{n-M\leq m\leq n}v_{n,m}S_m
\]
satisfies
\[
 \boxed{A_n-A_n^{\langle M\rangle}\in \pi^{w(M+1)}R}
\]
uniformly in $n$.
\end{theorem}

\begin{proof}
Every omitted term has depth at least $M+1$, hence lies in $\pi^{w(M+1)}R$. A finite sum remains in that ideal.
\end{proof}

\begin{principle}[Visible weight versus accidental cancellation]
The weight attached to a transfer coefficient gives a uniform certified lower bound. Additional divisibility produced by the source value, a binomial coefficient, or cancellation among terms is real but row-dependent. A proof-producing local compiler records both the visible kernel weight and the sharper realized valuation.
\end{principle}

\subsection{The two unit geometries}
For an odd prime $p$,
\[
 \bbZ_p^\times\cong\mu_{p-1}\times(1+p\bbZ_p).
\]
At the dyadic prime,
\[
 \bbZ_2^\times\cong\{\pm1\}\times(1+4\bbZ_2).
\]
Thus the natural odd-prime character address is taken modulo $p-1$, whereas the first dyadic character address is only parity. After that finite character split, the principal-unit coordinate is handled by Mahler expansion and progressively finer $p$-power polyphase refinement.

\begin{boundary}[No substitution principle at $p=2$]
The dyadic theory is not obtained by inserting $p=2$ into an odd-prime formula. The sign packet, the depth of the principal-unit subgroup, and the exact transfer weights all change. The correct unification occurs at the level of source-kernel factorization, not at the level of identical coefficients.
\end{boundary}

\subsection{Odd-power coordinates}
Write
\[
 \sech t=\sum_{n\geq0}E_n\frac{t^n}{n!},
\]
so $E_n=0$ for odd $n$. For an odd prime $p$, define
\[
 \boxed{S_n(p)=\sum_{a=0}^{p-1}(-1)^a(2a+1)^n.}
\]

\begin{lemma}[Odd-power exponential generating function]
\[
 \boxed{\sum_{n\geq0}S_n(p)\frac{t^n}{n!}=\frac{1+e^{2pt}}2\sech t.}
\]
Its even part is
\[
 \boxed{\sum_{r\geq0}S_{2r}(p)\frac{t^{2r}}{(2r)!}=\cosh^2(pt)\sech t.}
\]
\end{lemma}

\begin{proof}
The finite geometric sum gives
\[
 \sum_{a=0}^{p-1}(-1)^ae^{(2a+1)t}
 =e^t\frac{1-(-e^{2t})^p}{1+e^{2t}}
 =\frac{1+e^{2pt}}2\sech t,
\]
because $p$ is odd. Averaging the identity at $t$ and $-t$ gives the even part.
\end{proof}

Let $Z_n$ denote the Euler zigzag numbers,
\[
 \sec x+\tan x=\sum_{n\geq0}Z_n\frac{x^n}{n!}.
\]
Then
\[
 \sech^2x=\sum_{m\geq0}(-1)^mZ_{2m+1}\frac{x^{2m}}{(2m)!}.
\]

\begin{theorem}[Odd-prime Euler reconstruction]
For every odd prime $p$ and every even $n\geq0$,
\[
 \boxed{E_n=\sum_{m=0}^{n/2}(-1)^mZ_{2m+1}\binom n{2m}p^{2m}S_{n-2m}(p).}
\]
The transformation is unit lower triangular in the even index.
\end{theorem}

\begin{proof}
The even source identity implies
\[
 \sech t=\sech^2(pt)\sum_{r\geq0}S_{2r}(p)\frac{t^{2r}}{(2r)!}.
\]
Insert the divided-power expansion of $\sech^2(pt)$ and compare coefficients.
\end{proof}

\begin{corollary}[Uniform odd-prime ladder]
For every $M\geq0$,
\[
 \boxed{E_n\equiv\sum_{m=0}^{\min(M,n/2)}(-1)^mZ_{2m+1}\binom n{2m}p^{2m}S_{n-2m}(p)\pmod{p^{2M+2}}.}
\]
\end{corollary}

\begin{proof}
Every omitted term has $m\geq M+1$ and contains $p^{2m}$. All remaining factors are integers.
\end{proof}

\subsection{The weighted dyadic source}
For even $n\geq0$, define
\[
 \boxed{R_n=\frac{3^{n+2}-1}{4(n+1)(n+2)}.}
\]
Its even exponential generating function is
\[
 \boxed{\sum_{n\text{ even}}R_n\frac{t^n}{n!}=\frac{\sinh(2t)\sinh t}{2t^2}.}
\]
Define $\beta_m$ by
\[
 \boxed{\frac{z^2}{\sinh^2z}=\sum_{m\geq0}\beta_mz^{2m}.}
\]
The first values are
\[
 1,-\frac13,\frac1{15},-\frac2{189},\frac1{675},\ldots.
\]

\begin{lemma}[Dyadic integrality and exact source valuation]
For every $m\geq0$ and every even $n\geq0$,
\[
 \boxed{\beta_m\in\bbZ_2,\qquad R_n\in\bbZ_2^\times.}
\]
In particular,
\[
 \boxed{v_2(R_n)=0.}
\]
\end{lemma}

\begin{proof}
The series
\[
 \left(\frac{\sinh z}{z}\right)^2
 =\sum_{m\geq0}\frac{2^{2m+1}}{(2m+2)!}z^{2m}
\]
has coefficients in $\bbZ_2$ and constant term one, so its inverse has coefficients in $\bbZ_2$.

For $R_n$, put $k=n+2$. Since $k$ is even, lifting the exponent gives
\[
 v_2(3^k-1)=v_2(3-1)+v_2(3+1)+v_2(k)-1=2+v_2(k).
\]
Because $n+1$ is odd,
\[
 v_2\bigl(4(n+1)(n+2)\bigr)=2+v_2(k).
\]
The valuations are equal, so $R_n$ is a unit.
\end{proof}

\begin{theorem}[Exact weighted dyadic Euler reconstruction]
For every even $n\geq0$,
\[
 \boxed{E_n=\sum_{m=0}^{n/2}\beta_m2^{2m}\fall n{2m}R_{n-2m}.}
\]
\end{theorem}

\begin{proof}
Substitute $z=2t$ into the inverse series and multiply:
\[
 \frac{4t^2}{\sinh^2(2t)}\frac{\sinh(2t)\sinh t}{2t^2}
 =\frac{2\sinh t}{\sinh(2t)}=\sech t.
\]
Coefficient extraction turns multiplication by $t^{2m}$ into $\fall n{2m}$.
\end{proof}

\subsection{One local Euler-transfer theorem}
Define
\[
 A_{p,n}=\begin{cases}
 S_n(p),&p\text{ odd},\\
 R_n,&p=2\text{ and }n\text{ even},\\
 0,&p=2\text{ and }n\text{ odd},
 \end{cases}
\]
and
\[
 \kappa_{p,m}=\begin{cases}
 (-1)^mZ_{2m+1}p^{2m},&p\text{ odd},\\
 \beta_m2^{2m}(2m)!,&p=2.
 \end{cases}
\]

\begin{theorem}[Unified local Euler-transfer theorem]
For every prime $p$ and every even $n\geq0$,
\[
 \boxed{E_n=\sum_{m=0}^{n/2}\binom n{2m}\kappa_{p,m}A_{p,n-2m}.}
\]
Equivalently, if
\[
 H_p(t)=\sum_{r\geq0}A_{p,2r}\frac{t^{2r}}{(2r)!},
 \qquad
 K_p(t)=\sum_{m\geq0}\kappa_{p,m}\frac{t^{2m}}{(2m)!},
\]
then
\[
 \boxed{K_p(t)H_p(t)=\sech t.}
\]
For odd $p$,
\[
 H_p(t)=\cosh^2(pt)\sech t,
 \qquad K_p(t)=\sech^2(pt).
\]
For $p=2$,
\[
 H_2(t)=\frac{\sinh(2t)\sinh t}{2t^2},
 \qquad K_2(t)=\frac{4t^2}{\sinh^2(2t)}.
\]
\end{theorem}

\begin{proof}
The two source-kernel products were proved above. Comparing divided-power coefficients gives the common convolution law.
\end{proof}

\subsection{Exact odd-prime weight ledger}
For odd $p$,
\[
 v_p(\kappa_{p,m})=2m+v_p(Z_{2m+1}).
\]
The tangent-number identity
\[
 Z_{2m+1}=(-1)^m\frac{2^{2m+2}(2^{2m+2}-1)B_{2m+2}}{2m+2}
\]
gives
\[
 \boxed{v_p(Z_{2m+1})=v_p(2^{2m+2}-1)+v_p(B_{2m+2})-v_p(2m+2).}
\]
For a nonzero reconstruction term,
\[
 \boxed{\begin{aligned}
 W_p(n,m)={}&2m+v_p(Z_{2m+1})+v_p\binom n{2m}\\
 &+v_p\bigl(S_{n-2m}(p)\bigr).
 \end{aligned}}
\]
Kummer's theorem gives
\[
 v_p\binom n{2m}=\frac{s_p(2m)+s_p(n-2m)-s_p(n)}{p-1},
\]
so the binomial factor counts base-$p$ carries.

\subsection{Exact dyadic kernel weights}
\begin{lemma}[Bernoulli form of the dyadic inverse]
For every $m\geq1$,
\[
 \boxed{\beta_m=-\frac{2^{2m}(2m-1)B_{2m}}{(2m)!}.}
\]
Consequently,
\[
 \boxed{v_2(\beta_m)=s_2(m)-1.}
\]
\end{lemma}

\begin{proof}
Differentiate
\[
 \coth z=\frac1z+\sum_{m\geq1}\frac{2^{2m}B_{2m}}{(2m)!}z^{2m-1}
\]
and use $-(\coth z)'=\sinh(z)^{-2}$. The coefficient formula follows. The von Staudt--Clausen theorem gives $v_2(B_{2m})=-1$, and Legendre's formula gives
\[
 v_2((2m)!)=2m-s_2(2m)=2m-s_2(m).
\]
\end{proof}

For $m\geq1$, define
\[
 \boxed{\kappa_{2,m}=\beta_m2^{2m}(2m)!=-2^{4m}(2m-1)B_{2m}.}
\]

\begin{theorem}[Exact dyadic transfer depth]
For every $m\geq1$,
\[
 \boxed{v_2(\kappa_{2,m})=4m-1.}
\]
Therefore
\[
 \boxed{v_2\left(\kappa_{2,m}\binom n{2m}R_{n-2m}\right)=4m-1+v_2\binom n{2m}.}
\]
Equivalently,
\[
 \boxed{W_2(n,m)=4m-1+s_2(m)+s_2(n-2m)-s_2(n).}
\]
\end{theorem}

\begin{proof}
The first identity follows from $v_2(B_{2m})=-1$. The source $R_{n-2m}$ is a unit. Kummer's binary carry formula gives the last expression.
\end{proof}

\begin{theorem}[Optimal dyadic truncation]
Let
\[
 T_{2,M}(n)=\sum_{m=0}^{\min(M,n/2)}\kappa_{2,m}\binom n{2m}R_{n-2m}.
\]
Then
\[
 \boxed{E_n\equiv T_{2,M}(n)\pmod{2^{4M+3}}.}
\]
More precisely,
\[
 \boxed{v_2\bigl(E_n-T_{2,M}(n)\bigr)\geq\min_{M+1\leq m\leq n/2}\left(4m-1+v_2\binom n{2m}\right).}
\]
The uniform exponent $4M+3$ is best possible.
\end{theorem}

\begin{proof}
Every omitted term has valuation at least $4(M+1)-1=4M+3$. The rowwise bound retains the exact carry contribution. For optimality, take $n=2M+2$. There is one omitted term, with $m=M+1$, $R_0=1$, and binomial coefficient one, so its valuation is exactly $4M+3$.
\end{proof}

The first three sharp truncations are
\[
 \begin{aligned}
 E_n&\equiv R_n&&\pmod 8,\\
 E_n&\equiv R_n-\frac83\binom n2R_{n-2}&&\pmod{128},\\
 E_n&\equiv R_n-\frac83\binom n2R_{n-2}+\frac{128}{5}\binom n4R_{n-4}&&\pmod{2048}.
 \end{aligned}
\]

\begin{figure}[htbp]
\centering
\includegraphics[width=.86\textwidth]{20_local_kernel_weight_profiles.pdf}
\caption{The local Euler kernels share one convolution law but have different valuation geometry. The dyadic positive-depth kernel has exact slope four with offset minus one.}
\end{figure}

\section{Finite-field standard windows and Fourier transfer}

The congruential Euler matrices contain standard windows whose first residual inverse is independent of packet address. The remaining period question can be settled exactly by finite-field character transfer.

\subsection{The standard-window operator}
Let $p$ be odd, let $m\mid p-1$, and put
\[
 h=\frac{p-1}{m}.
\]
Define
\[
 \boxed{(T_{m,p}f)_n=\sum_{q=0}^{n}\binom{mn}{mq}f_q\quad\text{in }\bbF_p.}
\]
Let
\[
 V_h=\{f:\bbN\to\bbF_p:f_{n+h}=f_n\text{ for every }n\geq1\}.
\]

\begin{theorem}[Finite-field standard-window transfer]
\[
 \boxed{T_{m,p}(V_h)=V_h.}
\]
Moreover, $T_{m,p}$ restricts to an automorphism of $V_h$.
\end{theorem}

\begin{proof}
Every $f\in V_h$ can be written $f=d+c\delta_0$, where $d$ is fully $h$-periodic. Since $(T_{m,p}\delta_0)_n=1$, it is enough to treat $d$.

Expand
\[
 d_q=\sum_{\chi\in\mu_h}c_\chi\chi^q.
\]
The $m$th-power map on $\bbF_p^\times$ has image $\mu_h$, so choose $y$ with $y^m=\chi$. The root-of-unity filter gives
\[
 T_{m,p}(\chi^\bullet)_n=\frac1m\sum_{\zeta\in\mu_m}(1+\zeta y)^{mn}.
\]
If $1+\zeta y=0$, that term vanishes for $n\geq1$. Otherwise
\[
 \left((1+\zeta y)^m\right)^h=(1+\zeta y)^{p-1}=1.
\]
Each summand is $h$-periodic for $n\geq1$. The operator is unit lower triangular, hence injective. Since $V_h$ is finite-dimensional, the restriction is bijective.
\end{proof}

\begin{corollary}[Exact standard-window period]
Let $e^{(m,p)}$ be defined by
\[
 \sum_{q=0}^{n}\binom{mn}{mq}e_q^{(m,p)}=\delta_{n0}
 \quad\text{in }\bbF_p.
\]
Then
\[
 \boxed{e_{n+h}^{(m,p)}=e_n^{(m,p)}\qquad(n\geq1).}
\]
Thus the infinite inverse is determined by $e_0,e_1,\ldots,e_h$.
\end{corollary}

\begin{proof}
The sequence $\delta=(1,0,0,\ldots)$ lies in $V_h$. Its inverse image under the automorphism $T_{m,p}$ also lies in $V_h$.
\end{proof}

\begin{theorem}[Lucas standard-window reduction]
For $N=mp$, $0\leq j<p$, and $0\leq q\leq n$,
\[
 \boxed{\binom{mpn+j}{mpq}\equiv\binom{mn}{mq}\pmod p.}
\]
The first residual inverse is independent of $j$ and is periodic with period $h$ from index one onward.
\end{theorem}

\begin{proof}
Lucas' theorem gives
\[
 \binom{mpn+j}{mpq}\equiv\binom j0\binom{mn}{mq}\pmod p.
\]
The reduced matrix has unit diagonal and a unique inverse. Apply the preceding corollary.
\end{proof}

\begin{figure}[htbp]
\centering
\includegraphics[width=.91\textwidth]{21_standard_window_fourier_transfer.pdf}
\caption{Finite-field character transfer for a standard window. Each character becomes a finite sum of $h$-periodic modes, forcing the inverse window to have a finite periodic signature.}
\end{figure}

\section{Carmichael--Mahler packet transfer}

Finite modular dynamics lift into exact local analytic expansions. After a sufficiently long jump, each unit mode or invertible state matrix lies in a principal-unit neighborhood of the identity. Binomial expansion then becomes a Mahler expansion with one visible valuation gain at every depth.

\subsection{Finite Mahler coordinates}
For $f:\bbN\to R$, define $\Delta f(n)=f(n+1)-f(n)$. Since
\[
 \Delta\binom xr=\binom x{r-1},
\]
the binomial polynomials are the natural finite-difference coordinates.

\begin{theorem}[Finite Mahler reconstruction]
For known values $f(0),\ldots,f(N)$, put $c_r=\Delta^rf(0)$. Then, for $0\leq n\leq N$,
\[
 \boxed{f(n)=\sum_{r=0}^{n}c_r\binom nr.}
\]
More generally,
\[
 \boxed{\Delta^sf(n)=\sum_{r\geq s}c_r\binom n{r-s}.}
\]
\end{theorem}

\begin{proof}
This is Newton interpolation followed by the binomial difference identity.
\end{proof}

\begin{corollary}[Difference valuation bound]
If $c_r\in\bbZ_p$ and $v_p(c_r)\geq w(r)$, then
\[
 \boxed{v_p(\Delta^sf(n))\geq\min_{r\geq s}\left(w(r)+v_p\binom n{r-s}\right).}
\]
\end{corollary}

\subsection{Scalar unit modes}
For $L\geq1$ and odd $p$, put $\Lambda_{p,L}=(p-1)p^{L-1}$. For $p=2$ and $L\geq3$, put $\Lambda_{2,L}=2^{L-2}$.

\begin{theorem}[Carmichael--Mahler transfer for one unit mode]
Let $u\in\bbZ_p^\times$, let $H=\Lambda_{p,L}$, and put $q_{u,L}=u^H-1\in p^L\bbZ_p$. For every address $j$ and every $x\in\bbZ_p$,
\[
 \boxed{u^{Hx+j}=u^j\sum_{r\geq0}q_{u,L}^{\,r}\binom xr.}
\]
The Mahler coefficients satisfy
\[
 \boxed{v_p\left(u^jq_{u,L}^{\,r}\right)\geq Lr.}
\]
\end{theorem}

\begin{proof}
Use $u^{Hx+j}=u^j(1+q_{u,L})^x$ and the convergent $p$-adic binomial series.
\end{proof}

At the first stable level,
\[
 (H,L)=\begin{cases}(p-1,1),&p\text{ odd},\\(2,3),&p=2.\end{cases}
\]

\subsection{Polynomial-exponential modes}
\begin{theorem}[Polynomial-exponential Mahler transfer]
Let $f(n)=P(n)u^n$, where $P\in\bbZ_p[n]$ has degree $d$. Let $H=\Lambda_{p,L}$, fix $j$, and put
\[
 Q_j(x)=P(Hx+j),\qquad y=u^H,\qquad q=y-1.
\]
Then
\[
 f(Hx+j)=\sum_{r\geq0}c_{j,r}\binom xr,
\]
where
\[
 \boxed{c_{j,r}=u^j\sum_{k=0}^{\min(d,r)}\binom rk\Delta^kQ_j(0)y^kq^{r-k}.}
\]
Consequently,
\[
 \boxed{v_p(c_{j,r})\geq L\max(r-d,0).}
\]
For $R\geq d$,
\[
 \boxed{f(Hx+j)\equiv\sum_{r=0}^{R}c_{j,r}\binom xr\pmod{p^{L(R+1-d)}}}
\]
uniformly on $\bbZ_p$.
\end{theorem}

\begin{proof}
On the polynomial factor, the difference operator acts as $q+y\Delta$. The two summands commute. Expanding $(q+y\Delta)^r$ and evaluating at zero gives the coefficient formula. Since $\Delta^kQ_j=0$ for $k>d$, each surviving term contains at least $q^{r-d}$.
\end{proof}

\subsection{Matrix state systems}
Let
\[
 v_{n+1}=Cv_n,\qquad a_n=w^\top C^nv_0,
\]
with $C\in M_d(\bbZ_p)$ invertible modulo $p$.

\begin{theorem}[Matrix Carmichael--Mahler transfer]
Assume that $C^H\equiv I\pmod{p^L}$ and write $C^H=I+p^LB$. For every nonnegative integer $x$,
\[
 \boxed{a_{Hx+j}=\sum_{r=0}^{x}p^{Lr}\binom xr w^\top C^jB^rv_0.}
\]
The right side extends to $x\in\bbZ_p$ as the convergent infinite Mahler series with the same coefficients, and
\[
 \boxed{v_p(c_{j,r})\geq Lr.}
\]
Equivalently, with $b_j(x)=a_{Hx+j}$,
\[
 \boxed{\Delta^rb_j(x)=p^{Lr}w^\top C^j(I+p^LB)^xB^rv_0,}
\]
so $\Delta^rb_j(x)\equiv0\pmod{p^{Lr}}$.
\end{theorem}

\begin{proof}
Expand $C^{Hx+j}=C^j(I+p^LB)^x$ by the matrix binomial theorem. One difference multiplies the expansion by $p^LB$.
\end{proof}

\begin{example}[Dyadic Fibonacci packet]
For
\[
 C=\begin{pmatrix}1&1\\1&0\end{pmatrix},
\]
one has
\[
 C^{12}=\begin{pmatrix}233&144\\144&89\end{pmatrix}
 =I+8\begin{pmatrix}29&18\\18&11\end{pmatrix}.
\]
Therefore
\[
 \boxed{\Delta_{12}^rF_n\equiv0\pmod{8^r}.}
\]
\end{example}

\subsection{Adelic synchronization}
\begin{theorem}[Finite adelic Mahler transfer]
Let $Q=\prod_{p\in S}p^{L_p}$, let $C\in M_d(\bbZ)$ be invertible modulo $Q$, and choose $H$ with $C^H\equiv I\pmod Q$. Write $C^H=I+QB$. Then, for every nonnegative integer $x$,
\[
 \boxed{a_{Hx+j}=\sum_{r=0}^{x}Q^r\binom xr w^\top C^jB^rv_0,}
\]
and
\[
 \boxed{\Delta_H^ra_n\equiv0\pmod{Q^r}.}
\]
\end{theorem}

\begin{proof}
Use the same matrix binomial expansion over $\bbZ$, then reduce at each selected prime.
\end{proof}

\section{Local Euler interpolation}

\subsection{Odd-prime depletion}
The odd source contains one nonunit base, $2a+1=p$. Put
\[
 \sigma_p=(-1)^{(p-1)/2},\qquad S_n^\times(p)=S_n(p)-\sigma_pp^n.
\]

\begin{theorem}[Depleted odd-prime Euler reconstruction]
For every even $n$,
\[
 \boxed{(1-\sigma_pp^n)E_n=\sum_{m=0}^{n/2}(-1)^mZ_{2m+1}\binom n{2m}p^{2m}S_{n-2m}^\times(p).}
\]
\end{theorem}

\begin{proof}
The removed source mode contributes $\sigma_p\cosh(pt)$. Multiplication by $\sech^2(pt)$ gives $\sigma_p\sech(pt)$. Subtracting from $\sech t$ gives coefficient $(1-\sigma_pp^n)E_n$.
\end{proof}

Fix an even residue $j\in\{0,2,\ldots,p-3\}$. For $x\in\bbZ_p$, define
\[
 \begin{aligned}
 \mathscr E_{p,j}^\times(x)=\sum_{m\geq0}
 &(-1)^mZ_{2m+1}p^{2m}\binom{(p-1)x+j}{2m}\\
 &\times\sum_{\substack{0\leq a<p\\2a+1\neq p}}
 (-1)^a(2a+1)^{j-2m}\left((2a+1)^{p-1}\right)^x.
 \end{aligned}
\]

\begin{theorem}[Odd-prime analytic Euler branch]
The series converges uniformly on $\bbZ_p$ and defines a $p$-adic analytic function. For every nonnegative integer $x$,
\[
 \boxed{\mathscr E_{p,j}^\times(x)=\left(1-\sigma_pp^{(p-1)x+j}\right)E_{(p-1)x+j}.}
\]
\end{theorem}

\begin{proof}
Every retained base is a unit and its $(p-1)$st power lies in $1+p\bbZ_p$. The binomial factor is integral, and $p^{2m}$ forces uniform convergence. At nonnegative integer $x$, the series terminates and becomes the depleted reconstruction formula.
\end{proof}

\begin{boundary}[The depleted factor is an interpolation identity on integers]
The analytic function is defined by the unit-source series. The displayed factor $p^{(p-1)x+j}$ is evaluated only when $x$ is a nonnegative integer. No analytic function $x\mapsto p^x$ on all of $\bbZ_p$ is assumed or required.
\end{boundary}

\subsection{Direct dyadic interpolation}
For $x\in\bbZ_2$, define
\[
 \boxed{\mathscr R(x)=\frac{9^{x+1}-1}{8(x+1)(2x+1)}.}
\]
The apparent singularity at $x=-1$ is removable, with value $-\log(9)/8$.

\begin{lemma}[Integral analytic dyadic source]
The function $\mathscr R$ is $2$-adic analytic and takes values in $\bbZ_2$. For every $n\geq0$,
\[
 \boxed{\mathscr R(n)=R_{2n}.}
\]
\end{lemma}

\begin{proof}
Put $y=x+1$. The quotient
\[
 \frac{9^y-1}{8y}=\sum_{k\geq1}\frac{8^{k-1}}{k}\binom{y-1}{k-1}
\]
is integral and uniformly convergent because $3(k-1)-v_2(k)\geq0$. The factor $2x+1$ is always a unit.
\end{proof}

Define
\[
 \boxed{\mathscr E_2(x)=\sum_{m\geq0}\kappa_{2,m}\binom{2x}{2m}\mathscr R(x-m).}
\]

\begin{theorem}[Direct dyadic Euler interpolation]
The series converges uniformly and analytically on $\bbZ_2$. For every nonnegative integer $n$,
\[
 \boxed{\mathscr E_2(n)=E_{2n}.}
\]
The truncation after depth $M$ is correct modulo $2^{4M+3}$ uniformly on $\bbZ_2$.
\end{theorem}

\begin{proof}
The binomial and source factors are integral. The coefficient $\kappa_{2,m}$ has valuation $4m-1$. At integer $n$, terms with $m>n$ vanish and the exact dyadic reconstruction remains.
\end{proof}

\begin{figure}[htbp]
\centering
\includegraphics[width=.93\textwidth]{22_odd_dyadic_local_geometry.pdf}
\caption{Odd and dyadic Euler branches share source-kernel reconstruction but differ in torsion address, principal-unit depth, nonunit correction, and precision growth.}
\end{figure}

\subsection{Address factorization and local visibility}
Let $N=2^aM$ with $M$ odd. For every residue $r$,
\[
 \boxed{P_{N,r}=P_{2^a,r\bmod2^a}P_{M,r\bmod M}.}
\]

\begin{principle}[CRT address versus local lattice]
CRT determines the discrete address of a coefficient. Local transfer kernels determine the precision at which that address becomes visible. A complete arithmetic packet records both.
\end{principle}

The combined pipeline is
\[
 \boxed{\begin{aligned}
 \text{global coefficient}&\longrightarrow\text{character address}\longrightarrow\text{principal-unit coordinate}\\
 &\longrightarrow\text{weighted local transfer}\longrightarrow\text{finite congruence truncation}\\
 &\longrightarrow\text{polyphase refinement}\longrightarrow\text{CRT or adelic reconstruction}.
 \end{aligned}}
\]
The next section develops polyphase refinement as an exact calculus.

% ===== END sections/13_local_arithmetic.tex =====

% ===== BEGIN sections/14_polyphase_refinement.tex =====
\section{Polyphase refinement and local packet cascades}

Polyphase refinement replaces one coarse packet by a finite family of finer packets while preserving exact synthesis. In the discrete sequence channel this is residue refinement. In the local channel it becomes a digit-tree decomposition of $\bbZ_p$. In Mahler coordinates it is represented by an explicit finite triangular transfer kernel. These descriptions can be compiled into one refinement graph.

\subsection{Discrete polyphase analysis and synthesis}
For $M\geq2$ and a sequence $a=(a_n)_{n\geq0}$, define
\[
 \boxed{a_r^{[M]}(q)=a_{Mq+r},\qquad0\leq r<M.}
\]
The analysis operator is
\[
 \cP_Ma=(a_0^{[M]},\ldots,a_{M-1}^{[M]}),
\]
and synthesis interleaves the components:
\[
 \boxed{a_{Mq+r}=a_r^{[M]}(q).}
\]
If
\[
 A(z)=\sum_{n\geq0}a_nz^n,
 \qquad A_r^{[M]}(w)=\sum_{q\geq0}a_{Mq+r}w^q,
\]
then
\[
 \boxed{A(z)=\sum_{r=0}^{M-1}z^rA_r^{[M]}(z^M).}
\]

\begin{theorem}[Exact nested polyphase factorization]
For $M,N\geq2$,
\[
 \boxed{\left(a_r^{[M]}\right)_s^{[N]}=a_{r+Ms}^{[MN]}.}
\]
Equivalently,
\[
 \boxed{\cP_{MN}=\Pi_{M,N}\left(\bigoplus_{r=0}^{M-1}\cP_N\right)\cP_M,}
\]
where $\Pi_{M,N}$ only reorders the phase labels.
\end{theorem}

\begin{proof}
\[
 \left(a_r^{[M]}\right)_s^{[N]}(q)
 =a_r^{[M]}(Nq+s)=a_{M(Nq+s)+r}=a_{MNq+Ms+r}.
\]
\end{proof}

For a mixed-radix path
\[
 \mathbf M=(M_1,\ldots,M_\ell),\qquad0\leq s_i<M_i,
\]
the final address is
\[
 \boxed{r(\mathbf s)=s_1+M_1s_2+M_1M_2s_3+\cdots+\left(\prod_{i=1}^{\ell-1}M_i\right)s_\ell.}
\]
The refined branch is
\[
 a_{\mathbf s}(q)=a_{(\prod_iM_i)q+r(\mathbf s)}.
\]
Only the product radix and accumulated address matter.

\subsection{The carry polyphase algebra}
Let $c=a*b$ be Cauchy convolution and write $A_r,B_r,C_r$ for the $M$-phase OGFs. Since two phase digits produce at most one carry,
\[
 \boxed{C_r(w)=\sum_{u+v=r}A_u(w)B_v(w)+w\sum_{u+v=r+M}A_u(w)B_v(w).}
\]
The first sum has zero phase carry. The second has one coarse-index delay.

\begin{theorem}[Polyphase shift matrix]
Let $E$ be the scalar shift $(Ea)_n=a_{n+1}$ and let $E_q$ shift the coarse index. Then
\[
 \boxed{\cP_M(Ea)=\mathsf S_M(E_q)\cP_Ma,}
\]
where
\[
 \mathsf S_M(E_q)=
 \begin{pmatrix}
 0&1&0&\cdots&0\\
 0&0&1&\cdots&0\\
 \vdots&&&\ddots&\vdots\\
 0&0&0&\cdots&1\\
 E_q&0&0&\cdots&0
 \end{pmatrix}.
\]
Moreover,
\[
 \boxed{\mathsf S_M(E_q)^M=E_qI_M.}
\]
\end{theorem}

\begin{proof}
For $r<M-1$,
\[
 (Ea)_{Mq+r}=a_{Mq+r+1}=a_{r+1}^{[M]}(q).
\]
For $r=M-1$,
\[
 (Ea)_{Mq+M-1}=a_{M(q+1)}=E_qa_0^{[M]}(q).
\]
After $M$ shifts every phase advances by one coarse index.
\end{proof}

\subsection{\texorpdfstring{$p$-adic}{p-adic} polyphase branches}
For a prime $p$, $a\geq1$, and $0\leq s<p^a$, define
\[
 \boxed{f_{a,s}(x)=f(p^ax+s).}
\]
Because
\[
 \bbZ_p=\bigsqcup_{s=0}^{p^a-1}(s+p^a\bbZ_p),
\]
the family $\{f_{a,s}\}$ reconstructs $f$ uniquely on $\bbZ_p$.

\begin{theorem}[Nested local branches]
For $a,b\geq1$,
\[
 \boxed{(f_{a,s})_{b,t}=f_{a+b,s+p^at}.}
\]
Thus local polyphase refinement is the $p$-adic digit tree.
\end{theorem}

\begin{proof}
\[
 (f_{a,s})_{b,t}(x)=f(p^a(p^bx+t)+s)=f(p^{a+b}x+p^at+s).
\]
\end{proof}

\subsection{The affine Mahler transfer kernel}
Let
\[
 \Phi_M(z)=(1+z)^M-1.
\]
For $M\geq1$, $s\geq0$, and $r,k\geq0$, define
\[
 \boxed{\Theta_{M,s}(r,k)=[z^r](1+z)^s\Phi_M(z)^k.}
\]
Equivalently,
\[
 \boxed{\Theta_{M,s}(r,k)=\sum_{j=0}^{k}(-1)^{k-j}\binom kj\binom{Mj+s}{r}.}
\]
Since $\ord_z\Phi_M=1$ and $\deg\Phi_M=M$,
\[
 \boxed{\Theta_{M,s}(r,k)=0\quad\text{unless}\quad k\leq r\leq Mk+s.}
\]

\begin{theorem}[Mahler polyphase transfer]
Let
\[
 f(x)=\sum_{r\geq0}c_r\binom xr
\]
be a convergent Mahler expansion. For $g_s(q)=f(Mq+s)$,
\[
 \boxed{g_s(q)=\sum_{k\geq0}d_{s,k}\binom qk,}
\]
where
\[
 \boxed{d_{s,k}=\sum_{r=k}^{Mk+s}c_r\Theta_{M,s}(r,k).}
\]
Each refined coefficient uses a finite window of the original Mahler coordinates.
\end{theorem}

\begin{proof}
\[
 \begin{aligned}
 (1+z)^{Mq+s}
 &=(1+z)^s\left((1+z)^M\right)^q\\
 &=(1+z)^s(1+\Phi_M(z))^q\\
 &=\sum_{k\geq0}\binom qk(1+z)^s\Phi_M(z)^k.
 \end{aligned}
\]
Taking $[z^r]$ gives the affine expansion of $\binom{Mq+s}{r}$. Substitute it into the Mahler series.
\end{proof}

\begin{theorem}[Polyphase kernel composition]
For $M,N\geq1$,
\[
 \boxed{\sum_{k\geq0}\Theta_{M,s}(r,k)\Theta_{N,t}(k,\ell)=\Theta_{MN,Mt+s}(r,\ell).}
\]
\end{theorem}

\begin{proof}
Apply the Mahler transfer first to $x=Mq+s$ and then to $q=Ny+t$. The composite substitution is $x=MNy+Mt+s$. Uniqueness of Mahler coefficients gives the identity.
\end{proof}

\subsection{Valuation transfer through refinement}
\begin{theorem}[Exact polyphase valuation ledger]
For every refined coefficient,
\[
 \boxed{v_p(d_{s,k})\geq\min_{k\leq r\leq Mk+s}\left(v_p(c_r)+v_p\bigl(\Theta_{M,s}(r,k)\bigr)\right).}
\]
If $v_p(c_r)\geq w(r)$ for a nondecreasing function $w$, then
\[
 \boxed{v_p(d_{s,k})\geq w(k).}
\]
Every monotone Mahler precision ladder survives arbitrary polyphase refinement without loss.
\end{theorem}

\begin{proof}
Apply the non-Archimedean triangle inequality to the finite transfer formula. Since $r\geq k$, monotonicity gives $w(r)\geq w(k)$, and $\Theta_{M,s}(r,k)$ is integral.
\end{proof}

Define the branchwise weight transform
\[
 \boxed{(\mathfrak W_{M,s}w)(k)=\min_{k\leq r\leq Mk+s}\left(w(r)+v_p(\Theta_{M,s}(r,k))\right).}
\]
A mixed-radix refinement tree propagates certified precision by applying these min-plus transforms along each path. Direct refinement can be sharper when cancellation increases the valuation of a composed kernel, but the stagewise ledger is always valid.

\subsection{Frobenius sparsification at prime-power radix}
\begin{theorem}[Prime-power polyphase sparsification]
Let $M=p^a$. Then
\[
 \boxed{\Theta_{p^a,s}(r,k)\equiv\binom{s}{r-p^ak}\pmod p,}
\]
with the binomial coefficient interpreted as zero outside $0\leq r-p^ak\leq s$. Consequently,
\[
 \boxed{d_{s,k}\equiv\sum_{u=0}^{s}\binom su c_{p^ak+u}\pmod p.}
\]
In particular,
\[
 \boxed{d_{0,k}\equiv c_{p^ak}\pmod p.}
\]
\end{theorem}

\begin{proof}
In characteristic $p$,
\[
 (1+z)^{p^a}-1\equiv z^{p^a}.
\]
Therefore
\[
 (1+z)^s\Phi_{p^a}(z)^k\equiv(1+z)^sz^{p^ak}\pmod p.
\]
Take coefficients and substitute into the Mahler transfer formula.
\end{proof}

This is the coefficient-space counterpart of Lucas reduction. At the first residual level, a prime-power refinement selects a short binomial window around $p^ak$ rather than mixing all Mahler coordinates.

\begin{figure}[htbp]
\centering
\includegraphics[width=.9\textwidth]{23_mahler_polyphase_kernel.pdf}
\caption{Support and residual sparsity of the affine Mahler transfer kernel. Each refined coefficient uses a finite triangular window, which collapses to a short Frobenius window modulo $p$.}
\end{figure}

\subsection{Precision amplification under \texorpdfstring{$p$-power}{p-power} refinement}
\begin{theorem}[Scalar polyphase precision amplification]
Let $u\in\bbZ_p^\times$ and $a,r\geq0$. For odd $p$,
\[
 \boxed{\Delta_{(p-1)p^a}^{\,r}u^n=u^n\left(u^{(p-1)p^a}-1\right)^r,}
\]
and
\[
 \boxed{v_p\left(\Delta_{(p-1)p^a}^{\,r}u^n\right)\geq(a+1)r.}
\]
For $p=2$ and odd $u$,
\[
 \boxed{\Delta_{2^{a+1}}^{\,r}u^n=u^n\left(u^{2^{a+1}}-1\right)^r,}
\]
and
\[
 \boxed{v_2\left(\Delta_{2^{a+1}}^{\,r}u^n\right)\geq(a+3)r.}
\]
\end{theorem}

\begin{proof}
One step difference multiplies a unit mode by $u^H-1$. For odd $p$, Fermat reduction and lifting the exponent give
\[
 v_p(u^{(p-1)p^a}-1)\geq a+1.
\]
For odd $u$,
\[
 v_2(u^{2^{a+1}}-1)=v_2(u^2-1)+a\geq a+3.
\]
\end{proof}

\begin{lemma}[Matrix precision amplification]
If $X=I+p^LB$ with $B\in M_d(\bbZ_p)$, then for every $a\geq0$ there is $B_a\in M_d(\bbZ_p)$ such that
\[
 \boxed{X^{p^a}=I+p^{L+a}B_a.}
\]
\end{lemma}

\begin{proof}
Induct on $a$. If $X_a=I+p^{L+a}B_a$, expand $X_a^p$. The linear term is divisible by $p^{L+a+1}$ and every higher term has at least that valuation.
\end{proof}

\begin{theorem}[Matrix polyphase precision amplification]
Assume $C^H=I+p^LB$. Put $H_a=Hp^a$ and write
\[
 C^{H_a}=I+p^{L+a}B_a.
\]
Then every refined packet $b_{a,j}(x)=a_{H_ax+j}$ has Mahler coefficients satisfying
\[
 \boxed{v_p(c_{a,j,r})\geq(L+a)r.}
\]
Equivalently,
\[
 \boxed{\Delta_{H_a}^{\,r}a_n\equiv0\pmod{p^{(L+a)r}}.}
\]
\end{theorem}

\begin{proof}
Apply the matrix Carmichael--Mahler theorem to the amplified congruence for $C^{H_a}$.
\end{proof}

\begin{example}[Refined dyadic Fibonacci ladder]
Since $C^{12}=I+8B$ for the Fibonacci companion matrix,
\[
 C^{12\cdot2^a}\equiv I\pmod{2^{a+3}}.
\]
Therefore
\[
 \boxed{\Delta_{12\cdot2^a}^{\,r}F_n\equiv0\pmod{2^{(a+3)r}}.}
\]
Each dyadic refinement doubles the step and adds one guaranteed binary digit per Mahler depth.
\end{example}

\begin{figure}[htbp]
\centering
\includegraphics[width=.91\textwidth]{24_polyphase_precision_tree.pdf}
\caption{A local precision tree. Refining a principal-unit packet by one $p$-power digit increases the guaranteed Mahler slope from $L$ to $L+1$.}
\end{figure}

\subsection{Character-polyphase coordinates}
For odd $p$, write
\[
 \boxed{n=j+(p-1)(s+p^aq),}
\]
with $0\leq j<p-1$ and $0\leq s<p^a$. The address $j$ is the finite character packet. The address $s$ is the principal-unit polyphase digit. For a unit mode,
\[
 u^n=u^{j+(p-1)s}\left(u^{(p-1)p^a}\right)^q.
\]
At $p=2$, write
\[
 \boxed{n=j+2(s+2^aq),}
\]
with $j\in\{0,1\}$ and $0\leq s<2^a$. Then
\[
 u^n=u^{j+2s}\left(u^{2^{a+1}}\right)^q.
\]

\begin{principle}[Local phase lattice]
The finite character coordinate resolves torsion. The polyphase coordinate resolves principal-unit digits. The remaining coarse variable carries a rapidly convergent Mahler expansion:
\[
 \boxed{\text{character phase}\longrightarrow\text{$p$-power polyphase address}\longrightarrow\text{principal-unit Mahler depth}.}
\]
\end{principle}

\subsection{Application to the Euler branches}
The odd-prime depleted branch $\mathscr E_{p,j}^\times(x)$ and the direct dyadic branch $\mathscr E_2(x)$ are analytic functions on their local domains. They can be refined by $x=p^aq+s$. The Mahler polyphase theorem gives each branch an exact finite coefficient-transfer matrix. The valuation theorem preserves the original truncation ladder, while precision amplification strengthens it whenever the branch is represented by unit modes or an invertible finite state.

For the odd-prime Euler branch, the two packet axes are
\[
 n=j+(p-1)(s+p^aq).
\]
For the dyadic even-index branch, the natural coordinate is
\[
 n=2(s+2^aq).
\]
Thus the odd and dyadic interpolation formulas fit into one mixed character-polyphase tree even though their source kernels and exact transfer weights remain different.

\subsection{First polyphase research targets}
\begin{problem}[Exact higher valuation of the polyphase kernel]
Determine a closed formula or sharp combinatorial bound for
\[
 v_p\bigl(\Theta_{p^a,s}(r,k)\bigr)
\]
in terms of base-$p$ carries and the digit geometry of $(r,k,s)$.
\end{problem}

\begin{problem}[Optimal mixed-radix refinement]
Given a target precision and coefficient range, choose a factorization $M=M_1\cdots M_\ell$ that minimizes packet width, transfer cost, and precision loss.
\end{problem}

\begin{problem}[Spectral-local signature transfer]
Combine the spectral torsion signature of a C-finite sequence with its local character-polyphase signature. The target invariant should record complex root-of-unity collision, finite-field character phase, and principal-unit Mahler slope simultaneously.
\end{problem}

\begin{problem}[Nonlinear polyphase trees]
Develop an analogue of the Mahler transfer kernel for compositional and algebraic fixed-point systems. The expected objects are packet-labeled rooted trees whose edges carry phase addresses and local valuation weights.
\end{problem}

\begin{boundary}
Polyphase refinement is exact analysis and synthesis. It does not by itself prove that every refined branch is simpler. Simplification requires a transferred invariant such as spectral collision, Frobenius sparsity, valuation gain, or nonlinear tree cancellation.
\end{boundary}

% ===== END sections/14_polyphase_refinement.tex =====


% ===== BEGIN sections/14_compiler_complexity.tex =====
\part{A Proof-Producing Sequence Compiler}
\section{Compiler architecture, normal forms, and complexity profiles}

The realization theorem becomes operational when a finite sequence specification is represented as a typed abstract syntax tree.  The compiler does not guess an intended rule from terms.  It accepts a declared rule, checks its legality, constructs exact evaluators and packet systems, and emits verifiable certificates.

\subsection{Typed abstract syntax}
A sequence specification is built from nodes of the following kinds.
\begin{center}
\begin{tabularx}{\textwidth}{L{.2\textwidth}Y L{.25\textwidth}}
\toprule
Node type&Required data&Semantic channel\\
\midrule
\texttt{Seed}&finite constants or initial vectors&all\\
\texttt{LinearState}&matrix, seed, output map&sequence/OGF\\
\texttt{SpectralTransfer}&active modes or minimal linear state&packet recurrence/signature\\
\texttt{Recurrence}&coefficient functions, initial data&sequence\\
\texttt{Differential}&operator and initial jet&OGF/EGF\\
\texttt{DigitSystem}&radix, digit matrices, seed, output&sequence\\
\texttt{Sum/Product}&typed child nodes&OGF/EGF/Dirichlet\\
\texttt{Hadamard}&two coefficient-compatible children&sequence\\
\texttt{Inverse}&unit child and topology&OGF/EGF\\
\texttt{Compose}&outer child, inner child of zero constant&series\\
\texttt{FixedPoint}&finite map, branch/filtration condition&series/state\\
\texttt{Riordan}&unit $g$, compositional unit $f$&OGF/EGF\\
\texttt{FilteredProduct}&factor rule and weight function&completed series\\
\texttt{PrimePower}&local rule and synthesis mode&arithmetic function\\
\texttt{LocalTransform}&source kernel, inverse kernel, valuation weights&local coefficients\\
\texttt{StandardWindow}&prime, section width, finite-field character data&finite-field inverse signature\\
\texttt{MahlerState}&jump matrix, local modulus, output map&local analytic packet\\
\texttt{PolyphaseRefine}&radix, phase path, affine Mahler kernel&sequence/local branches\\
\texttt{Packet}&modulus and packet address&all coefficient channels\\
\bottomrule
\end{tabularx}
\end{center}

Every node records its coefficient ring, topology, normalization, and legal coercions.  For example, EGF multiplication over $\bbZ$ is interpreted in the divided-power basis; a modular inverse is legal only when the constant term is a unit; and a recurrence division must be justified in the target ring.

\subsection{Compilation stages}
\begin{figure}[htbp]
\centering
\includegraphics[width=.96\textwidth]{14_sequence_compiler_pipeline.pdf}
\caption{A proof-producing sequence compiler.  It begins from a structural specification and ends with exact values plus a certificate, not with an unsupported pattern match.}
\end{figure}

\begin{algorithm}[Projector-seed compilation]
Given a typed syntax tree $\mathfrak P$, a requested degree $D$, and optional packet modulus $N$:
\begin{enumerate}
\item \textbf{Type check.} Verify rings, channels, units, branch conditions, and topology.
\item \textbf{Normalize.} Simplify explicit transforms, reduce linear states, and canonicalize declared subgrammars.
\item \textbf{Build dependencies.} For every output coefficient through $D$, compute the finite set of child coefficients required.
\item \textbf{Compile packets.} If $N$ is supplied, replace each node by its exact packet conjugate or packet operation law.
\item \textbf{Compile local coordinates.} Separate finite character phase from principal-unit dynamics, select odd-prime or dyadic source kernels, and attach a valuation ledger to every transfer edge.
\item \textbf{Compile polyphase refinement.} Factor a requested mixed radix, construct the affine Mahler kernels $\Theta_{M,s}$, propagate precision by the min-plus weight transforms, and retain a direct-refinement certificate for comparison.
\item \textbf{Minimize linear packets.} On a C-finite node, choose either exact reachable--observable reduction of $C^N$ or the spectral torsion--Fourier compiler.  Emit the reduced recurrence and, when requested, the finite torsion packet signature.
\item \textbf{Choose engine.} Select exact integer/rational, finite-field, $p$-adic truncation, symbolic, or mixed execution.
\item \textbf{Evaluate.} Process the dependency DAG in topological or coefficient order; iterate contractive blocks until the requested prefix stabilizes.
\item \textbf{Certify.} Emit residual identities, recurrence checks, transform inverses, standard-window signatures, local tail bounds, polyphase-kernel composition checks, branch data, and content hashes.
\end{enumerate}
\end{algorithm}

\begin{theorem}[Compiler soundness]
If every primitive node is interpreted according to its declared semantics and every implicit block passes its well-foundedness condition, then the compiler output through degree $D$ equals the first $D+1$ coefficients of the unique semantic object.  Packet or polyphase compilation followed by synthesis gives the same output as unsplit compilation. Local compilation returns the same coefficient modulo the certified precision ideal.
\end{theorem}
\begin{proof}
Induct over the normalized dependency graph.  Explicit nodes follow their coefficient identities.  Linear and triangular blocks are exact finite solves.  Contractive blocks stabilize by the formal fixed-point theorem.  Filtered products omit only factors invisible through degree $D$.  Packet equivalence follows from exact conjugacy and reconstruction at every node. Mahler polyphase equivalence follows from the affine kernel identity and its composition law. Local truncations omit only terms whose certified valuation exceeds the requested precision.
\end{proof}

\subsection{Execution channels}
\paragraph{Exact integer and rational.}
Use arbitrary-precision integers and normalized rational numbers.  This is the reference channel for formulas involving Bernoulli coefficients or triangular denominators.

\paragraph{Finite field and modular ring.}
Reduce only after verifying that every required denominator is a unit, or use a fraction-free representation whose final value can be reduced safely.

\paragraph{$p$-adic precision.}
Represent a value together with a guaranteed precision ideal $p^L\Zp$.  Every operation propagates both the residue and its precision.  Weighted triangular truncation can improve efficiency by discarding terms whose visible valuation exceeds $L$. Character packets, Carmichael--Mahler jumps, and polyphase refinements retain separate precision ledgers so that dyadic and odd-prime gains are never conflated.

\paragraph{Symbolic.}
Retain parameters, recurrences, matrices, generating functions, and proof identities.  Symbolic output is not accepted as certified merely because it simplifies numerically; the compiler records the transformation rule used.

\paragraph{Hybrid.}
A proof may be symbolic while a large coefficient table is computed modularly and reconstructed by CRT.  The certificate then includes modulus bounds and a reconstruction range.

\subsection{Restricted normal forms}
No canonical form exists for arbitrary programs.  Several important subgrammars do have exact reductions.

\begin{theorem}[Linear normal form]
A finite linear realization reduces canonically, up to similarity, to the reachable-observable shift representation on the Hankel column space.  Its dimension is the Hankel rank.
\end{theorem}

\begin{theorem}[Complete C-finite packet normal form]
For a persistent C-finite sequence, the data
\[
 (\chi_a,\{P_\lambda\}_{\lambda\in\Lambda},L_a,\mathfrak T_a)
\]
form a finite exact normal form for arithmetic-subsequence complexity: for every modulus and address they determine the reduced packet recurrence, rational denominator, Hankel rank, collision deficit, and phase-cancellation deficit.  The same outputs can be certified without factorization by exact reachable--observable reduction of the jump realization.
\end{theorem}
\begin{proof}
Apply the spectral torsion--Fourier packet transfer theorem and its exact periodicity statement.
\end{proof}

\begin{proposition}[Rational generating-function normal form]
Over a field, a rational OGF has a unique reduced fraction $P(z)/Q(z)$ with $\gcd(P,Q)=1$ and a chosen normalization of $Q$.  Equality of two rational-sequence specifications is decidable by cross multiplication.
\end{proposition}

\begin{proposition}[Finite automaton normal form]
A deterministic finite automaton with output has a unique minimal quotient up to state renaming, obtained by merging states with identical future output behavior.
\end{proposition}

\begin{proposition}[Triangular normal form]
A declared lower-triangular transform with invertible diagonal has a unique inverse on the same indexed coefficient module.  Equality of two such finite-rule transforms is reducible to equality of their entry formulas under the declared symbolic theory.
\end{proposition}

\begin{boundary}
Equivalence of unrestricted finite grammars is not decidable in general.  The compiler may prove equality by reduction to a supported subgrammar, a common annihilator plus initial values, a common fixed-point uniqueness theorem, or an exact transform diagram.  Failure to find such a proof is not proof of inequality.
\end{boundary}

\subsection{Equivalence certificates}
The compiler recognizes several certificate patterns.
\begin{enumerate}
\item \textbf{Recurrence certificate}: both sequences satisfy the same order-$r$ recurrence and agree on $r$ initial values.
\item \textbf{Spectral packet certificate}: an active-mode table, torsion partition, phase-polynomial table, and reduced recurrence agree with direct modal substitution or an exact jump-state minimization.
\item \textbf{Differential certificate}: both series satisfy the same normalized differential equation and initial jet.
\item \textbf{Rational certificate}: reduced rational OGFs agree.
\item \textbf{Algebraic branch certificate}: both satisfy the same polynomial system and the same uniqueness condition.
\item \textbf{Transform certificate}: $b=Ta$ and $a=T^{-1}b$ under exact triangular or Riordan identities.
\item \textbf{Product certificate}: logarithmic derivatives and constant terms agree.
\item \textbf{Local certificate}: exact equality modulo a stated power, with a proven tail valuation bound.
\item \textbf{Standard-window certificate}: finite character decomposition, unit-triangular inverse, and a verified period signature.
\item \textbf{Polyphase certificate}: affine kernel table, mixed-radix composition identity, synthesis address, and propagated valuation profile.
\item \textbf{Finite-state certificate}: minimized automata or linear representations are isomorphic.
\end{enumerate}

\subsection{Complexity profile}
A single scalar ``complexity'' hides different resources.  Associate to a presentation the vector
\[
 \boxed{\kappa(\mathfrak P)=
 (d_{\rm lin},s_{\rm syn},a_{\max},w_{\rm pkt},h_{\rm dep},\ell_{\rm loc},r_{\rm poly},c_{\rm cert}).}
\]
Here:
\begin{itemize}
\item $d_{\rm lin}$ is minimal known linear-state dimension;
\item $s_{\rm syn}$ is normalized syntax size;
\item $a_{\max}$ is maximal nonlinear arity;
\item $w_{\rm pkt}$ is the number of active packet states after minimization;
\item $h_{\rm dep}$ is coefficient dependency horizon or recurrence order;
\item $\ell_{\rm loc}$ is required local precision depth;
\item $r_{\rm poly}$ is refinement depth together with maximal active phase width;
\item $c_{\rm cert}$ is the cost of independently checking the emitted proof object.
\end{itemize}

\begin{figure}[htbp]
\centering
\includegraphics[width=.96\textwidth]{15_sequence_complexity_spectrum.pdf}
\caption{One useful spectrum of sequence grammars.  Greater expressive breadth generally reduces canonical identifiability from finite data.}
\end{figure}

\begin{definition}[Packet profile]
For a sequence $a$, let
\[
 \pi_a(N)=\bigl(\kappa(P_{N,0}a),\ldots,\kappa(P_{N,N-1}a)\bigr)
\]
under a fixed grammar and minimization convention.  The profile records rank drops, periodic packets, zero packets, local gains, and transform symmetries.
\end{definition}

\begin{corollary}[Finite linear packet profile]
If $a$ is persistent C-finite with torsion period $L_a$, the linear-rank coordinate of every profile $\pi_a(N)$ is generated by the finite table $\mathfrak T_a$.  Counts of addresses having any fixed rank are degree-one quasipolynomials in $N$.
\end{corollary}

\begin{proposition}[Submultiplicative inherited packet dimension]
If $a$ has a $d$-dimensional linear realization, every packet subsequence has a realization of dimension at most $d$.  The total inherited state for all $N$ packets can share one jump operator and need not be represented as $Nd$ unrelated dimensions.
\end{proposition}

\subsection{Discovery as ranked hypothesis search}
Although finite data cannot determine an intended rule, it can guide a ranked search inside a declared grammar.

\begin{algorithm}[Certified hypothesis search]
Given terms $a_0,\ldots,a_D$ and a grammar budget:
\begin{enumerate}
\item test periodic and quasipolynomial models;
\item compute exact Hankel ranks and search C-finite recurrences;
\item factor certified C-finite recurrences when feasible, compute torsion signatures, and inspect packet-rank collapse before moving to broader grammars;
\item search bounded-degree P-recursive recurrences;
\item test hypergeometric ratios and factorial normalizations;
\item test standard transforms and inverse transforms;
\item search algebraic equations and low-arity fixed points;
\item inspect digit kernels and residue packet complexity;
\item test local character packets, standard-window periods, Mahler slopes, and low-cost polyphase refinements;
\item test Euler products, Dirichlet multiplicativity, and local congruence patterns;
\item rank surviving models by declared complexity and predictive validation;
\item promote a model only after an independent symbolic certificate is obtained.
\end{enumerate}
\end{algorithm}

\begin{boundary}[No finite-data certainty]
A candidate can be certified as a consequence of a proposed rule.  The finite data cannot certify that the proposed rule was the unique intended source of the data.
\end{boundary}

\subsection{Proof object design}
A portable sequence certificate should include:
\begin{enumerate}
\item a machine-readable typed specification;
\item normalization and version-independent mathematical formulas;
\item initial data and branch conditions;
\item exact output hashes or coefficient tables for smoke tests;
\item symbolic identities sufficient to verify the infinite claim;
\item packet, mixed-radix, and affine Mahler refinement maps;
\item local precision and denominator ledgers;
\item a claim-status tag: theorem, conditional theorem, finite certificate, or conjecture.
\end{enumerate}

The certificate checker should be smaller and more conservative than the search engine.  This separates discovery complexity from verification complexity.

\subsection{Sequence database integration}
A database entry can carry several projector-seed presentations rather than one privileged formula.  A possible record schema is
\[
 \{
 \texttt{terms},\texttt{recurrences},\texttt{ogf},\texttt{egf},\texttt{automata},
 \texttt{transforms},\texttt{packets},\texttt{polyphase},\texttt{local},\texttt{certificates}
 \}.
\]
Cross-links are then proof diagrams: a recurrence-to-OGF conversion, a transform conjugacy, a packet-jump law, or a local reconstruction theorem.  This would turn a catalog of formulas into a graph of mathematically checked realizations.

\subsection{Implementation invariants}
The reference implementation follows five invariants.
\begin{enumerate}
\item No floating-point number decides an exact identity.
\item Every truncation carries a theorem explaining why omitted terms are invisible.
\item Every division states the ring in which the divisor is a unit.
\item Every packet or polyphase computation synthesizes back to the unsplit result in tests.
\item Every conjectural pattern is stored separately from theorem-backed compilation.
\end{enumerate}

% ===== END sections/14_compiler_complexity.tex =====

% ===== BEGIN sections/15_sequence_atlas.tex =====
\section{Sequence-family atlas}

The following atlas is not a claim that every listed family is new.  It records how heterogeneous classical sequences enter the unified grammar and which mechanism is structurally primary.  Many rows admit several alternative realizations.

\begin{figure}[htbp]
\centering
\includegraphics[width=.97\textwidth]{13_sequence_family_coverage.pdf}
\caption{Representative sequence families against active projector-seed mechanisms.  The important feature is overlap: one family often has several exact presentations.}
\end{figure}

\small
\begin{longtable}{L{.15\textwidth}L{.20\textwidth}L{.28\textwidth}L{.18\textwidth}}
\caption{Atlas of structured integer-sequence realizations.}\label{tab:atlas}\\
\toprule
Family&Seed specification&Primary mechanism&Status in this paper\\
\midrule
\endfirsthead
\toprule
Family&Seed specification&Primary mechanism&Status in this paper\\
\midrule
\endhead
Fibonacci&$F_{n+2}=F_{n+1}+F_n$&two-dimensional companion orbit; jump packets&proved instance\\
Lucas&same companion, different output seed&trace of companion powers; common jump law&proved instance\\
General Lucas $U(P,Q)$&$U_{n+2}=PU_{n+1}-QU_n$&quadratic companion and $N$th-power trace&proved family\\
Pell&$P_{n+2}=2P_{n+1}+P_n$&linear state&direct instance\\
Jacobsthal&$J_{n+2}=J_{n+1}+2J_n$&linear state&direct instance\\
Tribonacci&order-three constant recurrence&three-dimensional companion&direct instance\\
$k$-bonacci&order-$k$ constant recurrence&finite companion&direct family\\
Persistent C-finite class&active exponential-polynomial modes&complete torsion--Fourier packet signature&proved transfer class\\
Torsion-free recurrences&no root-of-unity ratio&all arithmetic subsequences preserve minimal rank&proved corollary\\
Torsion-degenerate recurrences&root-of-unity mode ratios&exact periodic packet-rank collapse&proved classification\\
Mersenne&$2^n-1$&two exponential modes&direct instance\\
Polynomial values&$a_n=P(n)$&unipotent Newton state&proved class\\
Quasi-polynomials&polynomial on each residue class&residue packets plus Newton states&proved class\\
Periodic sequences&finite cyclic state&residue/address automaton&proved class\\
Central binomial&$\binom{2n}{n}$&hypergeometric ratio; algebraic OGF&proved mechanism\\
Catalan&$C=1+zC^2$&contractive fixed point; binary trees; hypergeometric&proved instance\\
Fuss--Catalan&$C=1+zC^{r+1}$&Lagrange inversion; $(r+1)$-ary trees&proved family\\
Motzkin&$M=1+zM+z^2M^2$&polynomial fixed point&proved instance\\
Large Schr\"oder&$S=1+zS+zS^2$&polynomial fixed point&direct instance\\
Little Schr\"oder&normalization of large Schr\"oder&linear transform of algebraic seed&direct instance\\
Narayana&bivariate Catalan refinement&marked tree grammar and coefficient projection&direct extension\\
Bell&$e^{e^t-1}$&composition; differential system; set partitions&proved instance\\
Touchard&$e^{x(e^t-1)}$&Stirling transform and block count&proved family\\
Stirling second kind&$(e^t-1)^k/k!$&composition and triangular transform&direct family\\
Stirling first kind&powers of $\log(1+t)$&inverse composition transform&direct family\\
Ordered Bell&$(2-e^t)^{-1}$&unit inverse of EGF seed&direct instance\\
Derangements&$e^{-t}/(1-t)$&Appell product; first-order recurrence&direct instance\\
Involutions&$e^{t+t^2/2}$&finite exponential composition&direct instance\\
Set partitions with block weights&$\exp(\sum w_kt^k/k!)$&exponential formula&direct family\\
Euler signed secant&$\sech t$&congruential Euler boundary; odd/dyadic coordinates&proved family\\
Euler zigzag&$\sec t+\tan t$&reciprocal/logarithmic differential seed&used explicitly\\
Secant numbers&$\sec t$&transfer reciprocal at $N=2$&proved boundary\\
Tangent numbers&$\tan t$&transfer logarithmic derivative&proved boundary\\
Bernoulli&$t/(e^t-1)$&unit inverse and Appell family&direct boundary\\
Genocchi&$2t/(e^t+1)$&Bernoulli/Euler transform orbit&direct boundary\\
Euler polynomials&$2e^{xt}/(e^t+1)$&Appell orbit&proved mechanism\\
Bernoulli polynomials&$te^{xt}/(e^t-1)$&Appell orbit&proved mechanism\\
Hermite-type&$g(t)e^{x f(t)}$&Sheffer system&proved mechanism\\
Congruential Euler $C_{N,j}$&normalized reciprocal of exponential section&triangular packet inverse&proved family\\
Transfer secants&$\mathcal D_N^{-r}$&reciprocal semigroup and block convolution&proved family\\
Transfer tangents&$-\mathcal D_N'/\mathcal D_N$&higher Riccati system&proved family\\
Ordinary partitions&$\prod_{m\ge1}(1-z^m)^{-1}$&filtered product; Euler transform&proved instance\\
Distinct partitions&$\prod(1+z^m)$&filtered product and product identity&proved instance\\
Odd-part partitions&$\prod_{m\text{ odd}}(1-z^m)^{-1}$&filtered product&proved equivalent instance\\
Colored partitions&$\prod(1-z^m)^{-r}$&Euler transform with constant exponent seed&direct family\\
Overpartitions&$\prod(1+z^m)/(1-z^m)$&finitely ruled product&direct instance\\
Plane partitions&$\prod(1-z^m)^{-m}$&filtered product with $b_m=m$&direct instance\\
Euler-transform families&$\prod(1-z^m)^{-b_m}$&filtered product and divisor recurrence&proved class\\
Ap\'ery numbers&hypergeometric sum or cubic recurrence&nonautonomous rational state&proved instance\\
Factorials&$a_{n+1}=(n+1)a_n$&one-dimensional hypergeometric state&direct instance\\
Double factorials&rational step ratio&hypergeometric state&direct instance\\
Binomial rays&$\binom{an+b}{cn+d}$&hypergeometric state when admissible&direct family\\
Domb-type sums&finite hypergeometric sums&creative telescoping / P-recursive state&grammar inclusion\\
Diagonal sequences&diagonal of rational function&constant-term projector and multiplication operator&grammar inclusion\\
Thue--Morse&binary digit parity&two-state automatic system&proved instance\\
Rudin--Shapiro&parity of $11$ occurrences&four-state automaton&direct instance\\
Period-doubling&uniform substitution&digit operators&direct instance\\
Paperfolding&finite radix automaton&digit operators&direct instance\\
Sum of binary digits&$s_2(2n)=s_2(n)$, $s_2(2n+1)=s_2(n)+1$&two-dimensional affine digit state&proved mechanism\\
Automatic sequences&finite $k$-kernel&finite digit state&proved class\\
Regular sequences&finitely generated $k$-kernel module&linear digit state&proved class\\
M\"obius $\mu$&local seed $1,-1,0,\ldots$&multiplicative prime-power synthesis&proved mechanism\\
Liouville $\lambda$&local seed $(-1)^e$&completely multiplicative local mode&direct instance\\
Divisor count $\tau$&local seed $e+1$&multiplicative polynomial exponent seed&direct instance\\
General $\tau_r$&local seed $\binom{e+r-1}{r-1}$&Dirichlet convolution powers&direct family\\
Sum of divisors $\sigma_k$&local geometric sum&prime-power synthesis&direct family\\
Euler totient $\varphi$&$p^{e-1}(p-1)$&prime-power synthesis&direct instance\\
Jordan totient $J_k$&$p^{ek}(1-p^{-k})$&prime-power synthesis&direct family\\
Squarefree indicator&local bounded-exponent seed&prime-power synthesis&direct instance\\
$\omega(n)$&one per distinct prime&additive prime-power synthesis&direct instance\\
$\Omega(n)$&local exponent $e$&additive prime-power synthesis&direct instance\\
Dirichlet inverses&triangular divisor recursion&Dirichlet unit inverse&proved class\\
Binomial transforms&$b_n=\sum\binom nk a_k$&invertible triangular/Appell transform&proved class\\
Inverse binomial&alternating triangular inverse&transform conjugacy&proved class\\
Riordan arrays&$g(z)A(f(z))$&invertible substitution-multiplication&proved class\\
Sheffer sequences&$g(t)e^{xf(t)}$&exponential Riordan action&proved class\\
Mahler coefficient families&$f(x)=\sum c_r\binom xr$&finite differences and local completion&proved finite theory\\
Standard-window inverses&$T_{m,p}e=\delta_0$&finite-field Fourier modes and periodic signature&proved family\\
Carmichael--Mahler packets&$C^H=I+p^LB$&matrix binomial expansion and local analytic state&proved family\\
Affine polyphase branches&$f(Mq+s)$&finite Mahler kernel $\Theta_{M,s}$&proved transfer class\\
Prime-power polyphase&$M=p^a$&Frobenius-sparse residual kernel&proved corollary\\
Adelic state packets&$C^H\equiv I\pmod Q$&simultaneous Mahler congruence ladder&proved family\\
Odd-prime Euler coordinates&alternating odd-power sums&weighted triangular transfer&proved family\\
Dyadic Euler coordinates&$R_n$, $z^2/\sinh^2z$&separate $2$-adic weighted inverse&proved family\\
CRT packet families&coprime residue refinements&commuting coefficient projectors&proved class\\
\bottomrule
\end{longtable}
\normalsize

\subsection{Coverage by mechanism rather than by name}
The atlas is more useful when read vertically.  A handful of closure mechanisms covers many named rows:
\begin{enumerate}
\item companion evolution covers all rational OGFs;
\item spectral torsion--Fourier transfer classifies the exact minimal recurrence of every arithmetic subsequence in the persistent C-finite class;
\item nonautonomous finite states cover P-recursive and hypergeometric families;
\item composition and divided-power products cover assemblies and set partitions;
\item polynomial fixed points cover broad tree families;
\item filtered products cover partition-like families;
\item finite digit paths cover automatic and regular families;
\item prime-power synthesis covers multiplicative arithmetic functions;
\item triangular reciprocal and local coordinates cover Euler/Bernoulli-type families;
\item Carmichael--Mahler jumps convert modular state periods into analytic local packets;
\item affine polyphase kernels refine local coordinates through exact mixed-radix branches.
\end{enumerate}

\subsection{Multiple placements are expected}
Catalan numbers are simultaneously hypergeometric, algebraic, P-recursive, and tree-generated.  Bell numbers are compositional, differential, triangular-transform, and partition-count sequences.  Euler numbers are reciprocal analytic coefficients, triangular inverse coefficients, and local weighted reconstructions.  A unified theory should preserve these multiple placements rather than force every sequence into one canonical narrative.

\subsection{What an OEIS-scale audit would measure}
A serious large-database experiment should report at least:
\begin{enumerate}
\item the fraction of entries with a recognized finite grammar;
\item the grammar class and normalized syntax size;
\item whether the realization was imported from a definition or inferred from terms;
\item whether an infinite certificate was proved or only a finite match obtained;
\item complete torsion signatures for C-finite entries and packet complexity profiles for broader classes;
\item transform-orbit links between entries;
\item local congruence certificates and finite standard-window signatures;
\item polyphase branch kernels, valuation profiles, and mixed-radix composition certificates;
\item unresolved entries and failure modes.
\end{enumerate}
A raw percentage without these distinctions would overstate what was achieved.

% ===== END sections/15_sequence_atlas.tex =====

% ===== BEGIN sections/16_boundaries_open_problems.tex =====
\section{Boundaries, solved transfer targets, and the forward research program}

The sequence branch now contains both a broad realization theorem and one complete invariant-transfer theorem.  The distinction remains essential.  Realization explains how a declared structure generates coefficients.  Transfer explains how a nontrivial invariant behaves under a systematic family of refinements.

\subsection{What has been established}
Within the declared grammar, the paper establishes:
\begin{enumerate}
\item exact lossless residue decomposition and arbitrary compatible refinement;
\item a finite realization theorem for well-founded specifications;
\item an exact packet compiler commuting with synthesis;
\item minimal linear dimension through Hankel rank;
\item common jump operators for all residue subsequences of a linear state;
\item the exact spectral torsion--Fourier transfer theorem for every persistent C-finite sequence;
\item exact minimal packet recurrences, reduced rational denominators, and Hankel ranks;
\item a complete root-of-unity rank-drop criterion and collision--phase deficit decomposition;
\item exact double periodicity, finite torsion packet signatures, and quasipolynomial rank counts;
\item classwise Fourier information conservation and a packet uncertainty principle;
\item finite-support transient reduction for arbitrary C-finite sequences;
\item P-recursive and automatic packet closure;
\item conjugate transport of projectors under invertible transforms;
\item packet Fa\`a di Bruno and packet-labeled rooted-tree expansions;
\item coefficientwise completion of finitely ruled filtered products;
\item prime-power reconstruction for multiplicative functions;
\item exponential-section, reciprocal, tangent, congruential Euler, and block-Appell families;
\item exact odd-prime and separate weighted dyadic Euler reconstructions;
\item Mahler difference coordinates, Lucas standard-window reduction, and CRT packet factorization;
\item a proof-producing compiler design and finite verification suite.
\end{enumerate}

\subsection{Two previously open linear targets now closed}
The earlier version isolated two conjectures:
\begin{enumerate}
\item an exact packet-rank spectrum formula for C-finite sequences;
\item eventual periodic control of C-finite packet-rank drops.
\end{enumerate}
Both are resolved in stronger form by \cref{thm:spectral-torsion-fourier}.  The rank formula is exact, includes confluent Jordan-polynomial cancellation, and produces the reduced recurrence itself.  The periodicity is not merely eventual: after removing the explicit finite nilpotent transient, it is exactly periodic in both the sampling modulus and packet address.

This changes the forward program.  The linear stratum no longer needs more examples of aliasing.  It has a complete classification theorem.  The next task is to identify analogous native invariants in larger strata.

\subsection{What has not been established}
The paper does not prove:
\begin{enumerate}
\item that every integer sequence has a short presentation in the restricted grammar;
\item that almost every database entry belongs to one specified class;
\item that terms alone identify an intended rule;
\item that a canonical minimal presentation exists across all grammar types;
\item that every broad representation yields a new theorem or faster algorithm;
\item historical priority for every generalized packet family or every formulation in the transfer section;
\item a complete packet-minimality theorem for P-recursive, algebraic, automatic, product, or nonlinear fixed-point systems;
\item every conjectural local periodicity or valuation bound suggested by computation.
\end{enumerate}

\subsection{Canonical sequence normal forms}
\begin{conjecture}[Canonical modular normal form]
There exists a practically broad subgrammar containing finite linear states, finite automata, algebraic fixed points with selected branches, restricted Riordan transforms, and effectively filtered products for which every object admits a finite canonical decomposition into irreducible seed blocks and explicit transform links.
\end{conjecture}

A plausible route is stratified rather than global:
\[
 \text{linear}\to\text{holonomic}\to\text{algebraic/compositional}\to
 \text{digital/product/local},
\]
with canonical forms inside strata and certified morphisms between them.  The C-finite spectral signature supplies the first complete stratum-level invariant.

\subsection{The next transfer target: holonomic packet minimality}
\begin{problem}[Minimal multisection operator]
Given a minimal differential or polynomial-recurrence operator annihilating $A(z)$, construct the minimal annihilator of $P_{N,j}A$ with quantitative order and degree bounds.  Separate unavoidable monodromy aliasing from address-dependent cancellation, analogously to the torsion--phase decomposition in the C-finite theorem.
\end{problem}

The existing D-finite closure theorem proves that an annihilator exists.  It does not determine the minimal one.  A successful theorem would require a replacement for the finite set of characteristic roots, likely involving local exponents, monodromy, singularity orbits under $z\mapsto\omega_Nz$, and differential-module cyclicity.

\begin{conjecture}[Holonomic finite signature under finite monodromy torsion]
For a D-finite series whose relevant local monodromy ratios have finite order, the minimal multisection orders are controlled by a finite torsion signature after explicitly separating singular transients and apparent singularities.
\end{conjecture}

This is intentionally narrower than a universal D-finite claim.  It mirrors the precise hypothesis that made the linear signature finite.

\subsection{Algebraic packet degree and branch transfer}
\begin{conjecture}[Packet algebraicity bound]
If $A(z)$ is algebraic of degree $d$ over $K(z)$, then the joint field generated by its $N$ residue packets has degree bounded effectively in terms of $d$ and $N$, and each packet satisfies an explicitly constructible algebraic equation whose degree can often be reduced by the symmetry of the root filter.
\end{conjecture}

The existence of an algebraic equation follows from closure under $z\mapsto\omega_N^rz$ and finite sums.  The research problem is a sharp degree bound, a minimal elimination procedure, and a branch-collision invariant analogous to the spectral torsion classes.

\subsection{Transform-orbit classification}
\begin{problem}[Restricted transform orbits]
Classify sequence families under rational or algebraic Riordan, Appell, Sheffer, and valuation-preserving triangular transforms.  Determine complete or computable invariants for important subclasses.
\end{problem}

The linear theorem suggests one concrete direction: determine how a restricted transform acts on the torsion signature, whether it preserves the global torsion period, and which transforms can create or remove packet-rank cancellation without altering global Hankel rank.

\begin{problem}[Transfer of torsion signatures through transforms]
For a declared invertible transform group $\mathfrak T$, characterize the relation between $\mathfrak T_a$ and $\mathfrak T_{Ta}$.  Identify subgroups preserving the signature, subgroups acting by an explicit finite rule, and transforms that can change the active root-ratio torsion.
\end{problem}

A successful theory would identify when two database entries are genuinely different seeds and when they are coordinate shadows of one object.

\subsection{Composition complexity}
\begin{conjecture}[Compressed packet Fa\`a di Bruno]
Packet-labeled partition sums admit orbit compression under equal block sizes and equal packet addresses, yielding algorithms whose cost depends on colored integer partitions rather than all set partitions.
\end{conjecture}

The deeper transfer target is to find an invariant of a finite compositional seed system that determines the minimal complexity of every packet fixed point, rather than merely expanding it exactly.

\subsection{Nonlinear tree signatures}
\begin{problem}[Tree-class transfer theorem]
For a contractive polynomial fixed-point system, classify packet-labeled rooted trees by the finite residue action on node arities and prove a finite signature theorem for a declared complexity invariant, such as algebraic degree, minimal system dimension, or dominant singularity orbit.
\end{problem}

The linear transfer theorem shows the desired pattern: a global finite object, an exact action of refinement, and a finite table classifying infinitely many packets.  The nonlinear analogue should not be expected to use characteristic roots, but it should preserve that architecture.

\subsection{Filtered-product acceleration}
\begin{problem}[Optimal product compiler]
For a finitely ruled product, choose adaptively among direct truncated multiplication, logarithmic-derivative recurrence, sparse functional identities, Newton iteration, and modular CRT reconstruction.  Prove complexity bounds in terms of output degree and coefficient height.
\end{problem}

The formal theorem guarantees finite work, not optimal work.  Compiler-level choice can make the difference between a conceptual representation and a competitive algorithm.

\subsection{Additive--multiplicative bridge}
\begin{problem}[Two-monoid sequence calculus]
Develop a common categorical or bialgebraic framework for Cauchy convolution on additive indices and Dirichlet convolution on multiplicative indices, including residue packets, prime-power coordinates, and transforms that genuinely intertwine the two products.
\end{problem}

This remains one of the clearest opportunities for a second semantic transfer theorem.  A useful result must move an invariant or obstruction between the additive and multiplicative channels rather than merely hosting both structures in parallel.

\subsection{Local arithmetic and polyphase extensions}
The previous standard-window period conjecture is resolved in this edition.  Finite-field Fourier transfer proves that the inverse window has exact period $h=(p-1)/m$ from index one onward.  The former search for an unspecified dyadic weight is also resolved for the Euler kernel: its positive-depth coefficient has exact valuation $4m-1$, and the uniform truncation exponent $4M+3$ is optimal.

\begin{problem}[Higher polyphase-kernel valuations]
Determine a sharp formula for
\[
 v_p\bigl(\Theta_{p^a,s}(r,k)\bigr)
\]
beyond its residual Frobenius window.  A complete answer should identify the precise base-$p$ carry statistic controlling every higher layer.
\end{problem}

\begin{problem}[Full prime-power standard-window lift]
Lift the finite residual period signature of the standard window from $\bbF_p$ to compatible signatures modulo $p^L$.  Determine whether the lift is governed by a finite matrix Mahler system and identify the minimal period and precision slope at every depth.
\end{problem}

\begin{problem}[Beyond Euler]
Construct odd/dyadic source-kernel factorizations for Bernoulli, tangent, Genocchi, Ap\'ery-type, partition, and other arithmetic families.  The target is a reusable local transfer theorem with an exact kernel-weight profile, not unrelated congruences family by family.
\end{problem}

\begin{problem}[Local normal forms]
For analytic sequence branches on $\bbZ_p$, determine when the character address, polyphase tree, and Mahler coefficient profile admit a canonical minimal representation.  The desired invariant should be stable under refinement and compatible with finite-state and spectral normal forms.
\end{problem}

A particularly strong next theorem would associate one finite spectral-local signature to a global arithmetic seed and use it to classify all recurrence ranks, character packets, and congruence-depth profiles under simultaneous residue and polyphase refinement.

\subsection{Minimal grammar and identifiability}
\begin{problem}[Restricted minimum-description theorem]
For a fixed decidable subgrammar, determine when a shortest presentation exists, whether it is unique up to declared equivalences, and how much initial data suffices to identify it among presentations below a fixed complexity budget.
\end{problem}

This is the correct formal replacement for the impossible unrestricted question ``Can the intended OEIS rule always be recovered from terms?''  Finite identification may become possible after the hypothesis class and complexity budget are fixed.

\subsection{Database-scale experimental program}
A rigorous coverage study should proceed in phases.
\begin{enumerate}
\item Parse explicit definitions, recurrences, and generating functions rather than infer from terms first.
\item Compile each parsed definition into a typed seed graph.
\item For every certified C-finite entry, compute its minimal recurrence, torsion period, finite packet signature, and exact degeneracy set.
\item Prove equivalence between multiple definitions when possible.
\item Generate packet profiles and transform links.
\item Search for low-complexity alternative presentations.
\item Reserve a held-out coefficient range for predictive testing.
\item Distinguish theorem certificates from empirical matches.
\item Publish failures and ambiguous classifications.
\end{enumerate}

A successful audit would quantify the phrase ``a very large fraction of structured sequences'' without pretending that a heterogeneous database is a probability space with a natural notion of ``almost every.''  The linear transfer theorem makes the C-finite part of this audit fully exact rather than experimental.

\subsection{Structural universality after the first transfer theorem}
\begin{conjecture}[Structural universality for natural sequence specifications]
There is a restricted, finitely presented, proof-producing projector-seed grammar that contains the standard finite mechanisms used to define natural integer sequences, is closed under their principal transforms and refinements, admits canonical or stratified normal forms on broad subclasses, and transfers nontrivial invariants between those subclasses.
\end{conjecture}

The current paper proves realization and closure for an explicit broad grammar, complete minimality transfer for the persistent C-finite stratum, and exact local/polyphase transfer laws for standard windows, Euler kernels, Mahler state systems, and affine refinement. Canonicality, quantitative database coverage, and corresponding minimality theorems for nonlinear and broader holonomic strata remain unfinished.

\subsection{Conclusion}
The sequence calculus is not merely the observation that every list can be projected into residue classes.  Its proved logical chain is now
\[
\boxed{
\begin{aligned}
\text{finite structural seed}
&\longrightarrow\text{typed operator grammar}
\longrightarrow\text{exact coefficient semantics}\\
&\longrightarrow\text{spectral and character packets}
\longrightarrow\text{exact finite signatures}\\
&\longrightarrow\text{principal-unit Mahler coordinates}
\longrightarrow\text{polyphase precision refinement}\\
&\longrightarrow\text{local or adelic completion}
\longrightarrow\text{proof-certified reconstruction}.
\end{aligned}}
\]

Linear recurrences, Bell compositions, Catalan trees, partition products, automatic digit paths, multiplicative prime-power rules, and Euler local coordinates are not forced into one elementary formula.  They are unified at the level of how finite information generates, refines, transfers, completes, and reconstructs an infinite coefficient object.

The spectral torsion--Fourier theorem supplies the first complete minimality transfer result. The local and polyphase theorems add a second transfer pattern: finite character data and modular state dynamics determine analytic branches, valuation ladders, and refinement gains. The central forward question is now precise: which finite global invariants control minimality in the holonomic, nonlinear, transform, and product strata?

% ===== END sections/16_boundaries_open_problems.tex =====

\appendix
% ===== BEGIN sections/A_worked_examples.tex =====
\section{Worked derivations and hand-checkable examples}

This appendix resolves representative families from their finite seeds to explicit coefficients, packet equations, and certificates.  The examples are deliberately heterogeneous: the point is to show the same realization pipeline across different native mechanisms.

\subsection{Fibonacci packets at modulus three}
Take
\[
 C=\begin{pmatrix}1&1\\1&0\end{pmatrix},
 \qquad
 v_0=\begin{pmatrix}1\\0\end{pmatrix},
 \qquad
 \lambda=\begin{pmatrix}0&1\end{pmatrix}.
\]
Then $F_n=\lambda C^nv_0$.  Direct multiplication gives
\[
 C^2=\begin{pmatrix}2&1\\1&1\end{pmatrix},
 \qquad
 C^3=\begin{pmatrix}3&2\\2&1\end{pmatrix}.
\]
The jump matrix has
\[
 \operatorname{tr}(C^3)=4,
 \qquad
 \det(C^3)=-1,
\]
so its characteristic polynomial is $x^2-4x-1$.  Every packet therefore satisfies
\[
 u_{m+2}=4u_{m+1}+u_m.
\]
The initial values and first checks are
\[
\begin{array}{c|rrrrrr}
 j&u_0&u_1&u_2&u_3&u_4&u_5\\
\hline
0&0&2&8&34&144&610\\
1&1&3&13&55&233&987\\
2&1&5&21&89&377&1597
\end{array}
\]
For the $j=2$ packet,
\[
 21=4\cdot5+1,
 \qquad
 89=4\cdot21+5,
 \qquad
 377=4\cdot89+21.
\]
Synthesis interlaces the three packet rows back into
\[
 0,1,1,2,3,5,8,13,21,34,55,89,\ldots.
\]

\subsection{A complete order-two torsion signature}
Consider
\[
 a_n=2^n+(-2)^n.
\]
The unsplit sequence has two active roots, $2$ and $-2$, each with a constant modal polynomial.  Its minimal recurrence polynomial is
\[
 (X-2)(X+2)=X^2-4,
\]
so $\rho(a)=2$.  Their ratio is $-1$, of order two, hence $L_a=2$.

If $N$ is odd, the jump roots are $2^N$ and $-2^N$, which remain distinct.  Every packet has rank two and satisfies
\[
 u_{m+2}=4^Nu_m.
\]
If $N$ is even, both roots map to $2^N$.  Choosing $2$ as representative gives the phase polynomial
\[
 R_j=1+(-1)^j.
\]
Therefore
\[
 a_{Nm+j}=
 \begin{cases}
 2^{j+1}(2^N)^m,&j\text{ even},\\
 0,&j\text{ odd}.
 \end{cases}
\]
The exact rank signature is
\[
\begin{array}{c|cc}
N\bmod2\backslash j\bmod2&0&1\\
\hline
0&1&0\\
1&2&2
\end{array}
\]
The rank drop is not merely an upper-bound phenomenon.  The even packet has nonzero $1\times1$ Hankel determinant and vanishing $2\times2$ determinant, while the odd packet is identically zero.

\subsection{Leading Jordan-layer cancellation}
Let
\[
 a_n=n2^n+(n+1)(-2)^n.
\]
The two modal polynomials are
\[
 P_+(x)=x,
 \qquad
 P_-(x)=x+1.
\]
The global recurrence polynomial is
\[
 (X-2)^2(X+2)^2,
\]
so $\rho(a)=4$.

For even $N$, both roots merge.  The two packet phase polynomials are
\[
 R_0(x)=P_+(x)+P_-(x)=2x+1,
\]
and
\[
 R_1(x)=P_+(x)-P_-(x)=-1.
\]
Thus the even-address packet keeps a degree-one polynomial mode and has rank two, while the odd-address packet cancels the entire degree-one layer and has rank one.

For example, at $N=2$:
\[
 a_{2m}=2m4^m+(2m+1)4^m=(4m+1)4^m,
\]
whereas
\[
 a_{2m+1}=(2m+1)2\cdot4^m-(2m+2)2\cdot4^m=-2\cdot4^m.
\]
The packet recurrences are therefore
\[
 (E-4)^2a_{2m}=0,
 \qquad
 (E-4)a_{2m+1}=0.
\]
Explicitly,
\[
 a_{2m}:1,20,144,832,4352,\ldots,
\]
and
\[
 a_{2m+1}:-2,-8,-32,-128,-512,\ldots.
\]
For the odd packet,
\[
 a_{2(m+1)+1}=4a_{2m+1}.
\]
For the even packet,
\[
 a_{2(m+2)}=8a_{2(m+1)}-16a_{2m}.
\]
The rank deficit decomposes as follows.  The merged class has global contribution $2+2=4$ and maximal modal rank two.  Hence the unavoidable collision deficit is two.  At even address the phase deficit is zero, while at odd address it is one:
\[
 \delta(2,0)=2+0=2,
 \qquad
 \delta(2,1)=2+1=3.
\]

\subsection{Hand reconstruction by finite Fourier inversion}
The same example gives a complete two-point Fourier analysis.  The forward packet-phase transform is
\[
 \begin{pmatrix}R_0\\R_1\end{pmatrix}
 =
 \begin{pmatrix}1&1\\1&-1\end{pmatrix}
 \begin{pmatrix}P_+\\P_-\end{pmatrix}.
\]
Since the inverse matrix is one half of the same matrix,
\[
 \begin{pmatrix}P_+\\P_-\end{pmatrix}
 =
 \frac12
 \begin{pmatrix}1&1\\1&-1\end{pmatrix}
 \begin{pmatrix}R_0\\R_1\end{pmatrix}.
\]
Substitution gives
\[
 P_+(x)=\frac{(2x+1)+(-1)}2=x,
\]
and
\[
 P_-(x)=\frac{(2x+1)-(-1)}2=x+1.
\]
An individual packet has lower rank, but the pair of packet phase polynomials reconstructs both global Jordan modes exactly.

\subsection{A complete mixed period-six signature}
Let $\omega^3=1$, $\omega\neq1$, and define
\[
 a_n=2^n+(-2)^n+3^n(\omega^n+\omega^{2n}).
\]
Because
\[
 \omega^n+\omega^{2n}
 =\begin{cases}2,&3\mid n,\\-1,&3\nmid n,\end{cases}
\]
the sequence is integer-valued.  It has four active roots:
\[
 2,-2,3\omega,3\omega^2.
\]
The first pair has torsion ratio of order two and the second pair has torsion ratio of order three.  Cross-pair ratios are not roots of unity because their absolute values differ.  Hence
\[
 L_a=\operatorname{lcm}(2,3)=6.
\]

The order-two class contributes
\[
 \begin{cases}
 2,&N\text{ odd},\\
 1,&N\text{ even and }j\text{ even},\\
 0,&N\text{ even and }j\text{ odd}.
 \end{cases}
\]
The order-three class contributes
\[
 \begin{cases}
 2,&3\nmid N,\\
 1,&3\mid N,
 \end{cases}
\]
because $1+\omega^j$ never vanishes.  Adding the two independent class contributions gives
\[
\begin{array}{c|rrrrrr}
N\bmod6\backslash j\bmod6&0&1&2&3&4&5\\
\hline
0&2&1&2&1&2&1\\
1&4&4&4&4&4&4\\
2&3&2&3&2&3&2\\
3&3&3&3&3&3&3\\
4&3&2&3&2&3&2\\
5&4&4&4&4&4&4
\end{array}
\]
This finite table gives every packet rank.  For example,
\[
 \rho_a(26,17)=\rho_a(2,5)=2,
\]
while
\[
 \rho_a(30,14)=\rho_a(0,2)=2.
\]
Here the first coordinate is always interpreted modulo six after choosing a positive representative for the actual sampling modulus.

\subsection{Rank-count quasipolynomial from the signature}
For the mixed example, take $N=6t+2$.  The relevant signature row is
\[
 3,2,3,2,3,2.
\]
Every complete address period contributes three packets of rank three and three packets of rank two.  The remaining prefix has addresses $0,1$, contributing one of each rank.  Hence
\[
 \boxed{
 c_3(6t+2)=3t+1,
 \qquad
 c_2(6t+2)=3t+1.
 }
\]
For $N=20=6\cdot3+2$, there are ten packets of rank three and ten of rank two.  No new recurrence minimization is required at modulus twenty; the finite signature already determines the count.

\subsection{Bell numbers from a two-equation seed system}
Start from
\[
 E'=E,
 \qquad
 B'=EB,
 \qquad
 E(0)=B(0)=1.
\]
Write
\[
 E(t)=\sum_{n\geq0}\frac{t^n}{n!},
 \qquad
 B(t)=\sum_{n\geq0}B_n\frac{t^n}{n!}.
\]
The coefficient of $t^n/n!$ in $EB$ is
\[
 \sum_{k=0}^{n}\binom nk B_k.
\]
Therefore
\[
 B_{n+1}=\sum_{k=0}^{n}\binom nk B_k.
\]
Starting from $B_0=1$:
\[
\begin{aligned}
B_1&=1,\\
B_2&=1+1=2,\\
B_3&=1+2\cdot1+2=5,\\
B_4&=1+3\cdot1+3\cdot2+5=15,\\
B_5&=1+4+6\cdot2+4\cdot5+15=52.
\end{aligned}
\]
The composition solution is recovered by integration:
\[
 \frac{B'}B=e^t,
 \qquad
 \log B=e^t-1,
 \qquad
 B=e^{e^t-1}.
\]

\subsection{Bell coefficients from set partitions}
At derivative order four, the set partitions of $\{1,2,3,4\}$ have block-count distribution
\[
 1,7,6,1
\]
for one, two, three, and four blocks.  Thus
\[
 T_4(x)=x+7x^2+6x^3+x^4,
\]
and
\[
 B_4=T_4(1)=15.
\]
This agrees with the recurrence calculation.  The same coefficient has been reconstructed through composition derivatives, a differential system, a triangular Stirling transform, and direct partition counting.

\subsection{Bell residue packets modulo three}
The coefficient packets begin
\[
\begin{array}{c|rrrrr}
 j&0&1&2&3&4\\
\hline
0&B_0=1&B_3=5&B_6=203&B_9=21147&B_{12}=4213597\\
1&B_1=1&B_4=15&B_7=877&B_{10}=115975&B_{13}=27644437\\
2&B_2=2&B_5=52&B_8=4140&B_{11}=678570&B_{14}=190899322
\end{array}
\]
They are not independent.  If $E_r=P_{3,r}e^t$ and $B_s=P_{3,s}B$, then
\[
 (B_{k+1})'=\sum_{r+s\equiv k\,(3)}E_rB_s.
\]
For example,
\[
 B_1'=E_0B_0+E_1B_2+E_2B_1.
\]
Synthesis of the three equations is exactly $B'=e^tB$.

\subsection{Catalan fixed point and parity packets}
Let $u=z+u^2$.  Write
\[
 u=z+c_2z^2+c_3z^3+c_4z^4+c_5z^5+\cdots.
\]
Coefficient comparison gives
\[
 c_2=1,
 \quad c_3=2,
 \quad c_4=5,
 \quad c_5=14,
\]
so
\[
 u=z+z^2+2z^3+5z^4+14z^5+42z^6+\cdots.
\]
Split into even and odd degree packets:
\[
\begin{aligned}
 u_0&=z^2+5z^4+42z^6+429z^8+\cdots,\\
 u_1&=z+2z^3+14z^5+132z^7+\cdots.
\end{aligned}
\]
The packet equations are
\[
 u_0=u_0^2+u_1^2,
 \qquad
 u_1=z+2u_0u_1.
\]
Through degree four,
\[
 u_1^2=z^2+4z^4+O(z^6),
 \qquad
 u_0^2=z^4+O(z^6),
\]
so the first equation gives $u_0=z^2+5z^4+O(z^6)$.  The second gives
\[
 z+2(z^2)(z)+2(5z^4\cdot z+z^2\cdot2z^3)+\cdots
 =z+2z^3+14z^5+\cdots.
\]

\subsection{Fuss--Catalan by Lagrange inversion}
Let $C=1+zC^3$ and put $u=C-1$.  Then
\[
 u=z(1+u)^3.
\]
Lagrange inversion gives
\[
 [z^n]u=\frac1n[w^{n-1}](1+w)^{3n}
 =\frac1n\binom{3n}{n-1}
 =\frac1{2n+1}\binom{3n}{n}.
\]
Thus
\[
 C(z)=1+z+3z^2+12z^3+55z^4+273z^5+\cdots.
\]
The ternary-tree grammar and the coefficient formula are exact translations of the same seed equation.

\subsection{Motzkin from a polynomial system}
The Motzkin equation is
\[
 M=1+zM+z^2M^2.
\]
Let $M=\sum_{n\geq0}M_nz^n$.  Then
\[
 M_0=1,
 \qquad
 M_n=M_{n-1}+\sum_{k=0}^{n-2}M_kM_{n-2-k}.
\]
Therefore
\[
\begin{aligned}
M_1&=1,\\
M_2&=M_1+M_0^2=2,\\
M_3&=M_2+2M_0M_1=4,\\
M_4&=M_3+(M_0M_2+M_1^2+M_2M_0)=9,\\
M_5&=M_4+(M_0M_3+M_1M_2+M_2M_1+M_3M_0)=21.
\end{aligned}
\]

\subsection{Partition product stabilization}
To compute $p(n)$ through $n=6$, factors with $m>6$ are invisible.  Multiply successively modulo $z^7$:
\[
 \prod_{m=1}^{6}(1-z^m)^{-1}
 \equiv
 1+z+2z^2+3z^3+5z^4+7z^5+11z^6
 \pmod{z^7}.
\]
The logarithmic recurrence gives the same values.  Since
\[
 (\sigma_1(1),\ldots,\sigma_1(6))=(1,3,4,7,6,12),
\]
we obtain
\[
 6p(6)=1p(5)+3p(4)+4p(3)+7p(2)+6p(1)+12p(0)
 =7+15+12+14+6+12=66.
\]
Hence $p(6)=11$.

\subsection{Distinct parts and odd parts}
Through degree eight,
\[
 \prod_{m=1}^{8}(1+z^m)
 \equiv
 \prod_{\substack{1\leq m\leq8\\m\text{ odd}}}(1-z^m)^{-1}
 \pmod{z^9}.
\]
Both sides have coefficients
\[
 1,1,1,2,2,3,4,5,6.
\]
The equality is not a finite coincidence; it follows factorwise from
\[
 \prod_{m\geq1}(1+z^m)
 =\prod_{m\geq1}\frac{1-z^{2m}}{1-z^m}
 =\prod_{m\text{ odd}}(1-z^m)^{-1}.
\]

\subsection{Binomial transform of Fibonacci}
Define
\[
 b_n=\sum_{k=0}^{n}\binom nkF_k.
\]
At $n=4$,
\[
 b_4=4F_1+6F_2+4F_3+F_4
 =4+6+8+3=21=F_8.
\]
In general, Binet modes $\alpha^n$ and $\beta^n$ are sent to $(1+\alpha)^n$ and $(1+\beta)^n$.  Since
\[
 1+\alpha=\alpha^2,
 \qquad
 1+\beta=\beta^2,
\]
we obtain
\[
 \boxed{b_n=F_{2n}.}
\]
The transformed characteristic polynomial is
\[
 x^2-3x+1,
\]
so $b_{n+2}=3b_{n+1}-b_n$.

\subsection{Thue--Morse digit path}
Take $n=13=(1101)_2$.  Starting from $t_0=0$, append digits from left to right using
\[
 x\xmapsto{0}x,
 \qquad
 x\xmapsto{1}1-x.
\]
The path is
\[
 0\xmapsto{1}1\xmapsto{1}0\xmapsto{0}0\xmapsto{1}1.
\]
Thus $t_{13}=1$, agreeing with the odd parity of three one-bits.  Fixing the final two binary digits gives the four radix packets
\[
\begin{aligned}
 t_{4n}&=t_n,\\
 t_{4n+1}&=1-t_n,\\
 t_{4n+2}&=1-t_n,\\
 t_{4n+3}&=t_n.
\end{aligned}
\]
Only two semantic kernel states remain.

\subsection{Ap\'ery recurrence checks}
With $A_0=1$ and $A_1=5$, the recurrence at $n=1$ gives
\[
 8A_2=117\cdot5-1=584,
 \qquad A_2=73.
\]
At $n=2$,
\[
 27A_3=535\cdot73-8\cdot5=39015,
 \qquad A_3=1445.
\]
At $n=3$,
\[
 64A_4=1463\cdot1445-27\cdot73=2{,}112{,}064,
 \qquad A_4=33001.
\]
The rational state recurrence and the finite binomial sum give independent exact presentations.

\subsection{Odd-prime Euler reconstruction at \texorpdfstring{$n=8$, $p=5$}{n=8, p=5}}
The coordinate values are
\[
 S_0=1,
 \quad S_2=49,
 \quad S_4=4705,
 \quad S_6=428689,
 \quad S_8=37665985.
\]
The zigzag coefficients entering the row are
\[
 Z_1=1,
 \quad Z_3=2,
 \quad Z_5=16,
 \quad Z_7=272,
 \quad Z_9=7936.
\]
Thus
\[
\begin{array}{c|r}
 m&(-1)^mZ_{2m+1}\binom8{2m}5^{2m}S_{8-2m}(5)\\
\hline
0&37{,}665{,}985\\
1&-600{,}164{,}600\\
2&3{,}293{,}500{,}000\\
3&-5{,}831{,}000{,}000\\
4&3{,}100{,}000{,}000\\
\hline
\text{sum}&1{,}385
\end{array}
\]
Truncation through $m=2$ gives $2{,}731{,}001{,}385$.  The difference from $E_8$ is
\[
 -2{,}731{,}000{,}000,
\]
whose $5$-adic valuation is six, exactly the guaranteed depth-$2$ modulus $5^6$.

\subsection{Weighted dyadic Euler row at \texorpdfstring{$n=8$}{n=8}}
The exact terms are
\[
 \frac{7381}{45},
 \quad-\frac{6560}{3},
 \quad\frac{163072}{15},
 \quad-\frac{204800}{9},
 \quad\frac{229376}{15}.
\]
With common denominator $45$, their numerators are
\[
 7381,-98400,489216,-1024000,688128,
\]
which sum to $62325=45\cdot1385$.  Hence the row reconstructs $E_8=1385$ exactly despite rational intermediate coordinates.

After retaining only the first term, the error is
\[
 1385-\frac{7381}{45}=\frac{54944}{45},
\]
and $v_2(54944/45)=5$, stronger than the sharp uniform depth-zero bound three.  The theorem relies only on the exact visible kernel weight and does not assume such extra cancellation.


\subsection{Sharpness of the dyadic ladder}
At truncation depth $M=1$, the sharp theorem predicts a uniform error in $2^7\bbZ_2$.  Take the extremal row $n=2M+2=4$.  Retaining $m=0,1$ gives
\[
 T_{2,1}(4)=R_4-\frac83\binom42R_2.
\]
Since
\[
 R_4=\frac{3^6-1}{4\cdot5\cdot6}=\frac{91}{15},
 \qquad
 R_2=\frac{3^4-1}{4\cdot3\cdot4}=\frac53,
\]
we obtain
\[
 T_{2,1}(4)=\frac{91}{15}-\frac{80}{3}=-\frac{103}{5}.
\]
The exact value is $E_4=5$, so
\[
 E_4-T_{2,1}(4)=\frac{128}{5},
 \qquad
 v_2\!\left(\frac{128}{5}\right)=7.
\]
Thus the modulus $2^{4M+3}=2^7$ is attained exactly and cannot be improved uniformly.

\subsection{A complete standard-window signature}
Take $p=5$ and $m=2$.  Then
\[
 h=\frac{p-1}{m}=2,
\]
and the inverse coefficients satisfy
\[
 \sum_{q=0}^{n}\binom{2n}{2q}e_q=\delta_{n0}\pmod5.
\]
Triangular inversion gives
\[
 e_0=1,
 \qquad
 e_1=-\binom20e_0\equiv4\pmod5,
\]
\[
 e_2=-\left(\binom40e_0+\binom42e_1\right)
      \equiv-(1+6\cdot4)\equiv0\pmod5.
\]
The next rows are
\[
 e_3\equiv4,
 \qquad e_4\equiv0,
 \qquad e_5\equiv4.
\]
Hence
\[
 \boxed{(e_n)_{n\geq0}=1,4,0,4,0,4,0,\ldots\quad\text{in }\bbF_5,}
\]
which exhibits the exact period two from index one.  One finite signature $(4,0)$ determines the entire residual inverse.

\subsection{Affine Mahler polyphase transfer at radix two}
Start from
\[
 f(x)=x^2=\binom x1+2\binom x2.
\]
Refine to the odd branch $g(q)=f(2q+1)$.  For $M=2$ and $s=1$,
\[
 \Phi_2(z)=2z+z^2,
\]
and
\[
 (1+z)\Phi_2(z)=2z+3z^2+z^3,
\]
\[
 (1+z)\Phi_2(z)^2=4z^2+8z^3+5z^4+z^5.
\]
Thus the relevant kernel values are
\[
 \Theta_{2,1}(1,0)=1,
 \quad
 \Theta_{2,1}(1,1)=2,
 \quad
 \Theta_{2,1}(2,1)=3,
 \quad
 \Theta_{2,1}(2,2)=4.
\]
With $c_0=0,c_1=1,c_2=2$, the refined coefficients are
\[
 d_0=c_1\Theta_{2,1}(1,0)=1,
\]
\[
 d_1=c_1\Theta_{2,1}(1,1)+c_2\Theta_{2,1}(2,1)=2+6=8,
\]
\[
 d_2=c_2\Theta_{2,1}(2,2)=8.
\]
Therefore
\[
 \boxed{(2q+1)^2=1+8\binom q1+8\binom q2.}
\]
At $q=2$, this gives $1+16+8=25$, exactly.  The branch coefficients also display the residual Frobenius sparsity: modulo two only the constant coordinate survives.

\subsection{Refined dyadic Fibonacci precision}
For the Fibonacci companion matrix,
\[
 C^{12}=I+8B,
 \qquad
 B=\begin{pmatrix}29&18\\18&11\end{pmatrix}.
\]
Squaring gives
\[
 C^{24}=(I+8B)^2=I+16(B+4B^2),
\]
so every step-$24$ branch has Mahler coefficients divisible by $2^{4r}$.  In finite-difference form,
\[
 \boxed{\Delta_{24}^{\,r}F_n\equiv0\pmod{2^{4r}}.}
\]
For $r=1$ and $n=0$,
\[
 F_{24}-F_0=46368,
 \qquad v_2(46368)=5\geq4.
\]
For the preceding step $12$, the guaranteed slope is only three:
\[
 F_{12}-F_0=144,
 \qquad v_2(144)=4\geq3.
\]
The extra dyadic polyphase digit doubles the step and raises the certified Mahler slope by one.

\subsection{Mahler reconstruction of \texorpdfstring{$n^3$}{n cubed}}
For $f(n)=n^3$,
\[
 \Delta f(0)=1,
 \quad
 \Delta^2f(0)=6,
 \quad
 \Delta^3f(0)=6,
 \quad
 \Delta^rf(0)=0\ (r\geq4).
\]
Therefore
\[
 \boxed{n^3=\binom n1+6\binom n2+6\binom n3.}
\]
At $n=5$ this gives
\[
 5+6\cdot10+6\cdot10=125.
\]
One further difference shifts the coefficients:
\[
 \Delta(n^3)=1+6\binom n1+6\binom n2=3n^2+3n+1.
\]

\subsection{Multiplicative reconstruction at \texorpdfstring{$n=360$}{n=360}}
Factor
\[
 360=2^3\,3^2\,5.
\]
Then
\[
\begin{aligned}
 \mu(360)&=0,\\
 \lambda(360)&=(-1)^{3+2+1}=1,\\
 \tau(360)&=(3+1)(2+1)(1+1)=24,\\
 \sigma_1(360)&=(1+2+4+8)(1+3+9)(1+5)=1170,\\
 \varphi(360)&=2^{2}(2-1)\,3^{1}(3-1)\,(5-1)=96.
\end{aligned}
\]
Each result is synthesized from three local prime-power seeds.

\subsection{Dirichlet convolution at a prime power}
Since the constant-one function $\mathbf1$ has local seed $1,1,1,\ldots$, its $r$-fold Dirichlet convolution has local series
\[
 (1-T)^{-r}=\sum_{e\geq0}\binom{e+r-1}{r-1}T^e.
\]
Therefore
\[
 \tau_r(p^e)=\binom{e+r-1}{r-1}.
\]
For $r=3$ and $e=4$,
\[
 \tau_3(p^4)=\binom62=15,
\]
the number of triples of nonnegative exponents summing to four.

\subsection{CRT packet address}
For modulus $40=8\cdot5$ and residue $27$,
\[
 27\equiv3\pmod8,
 \qquad
 27\equiv2\pmod5.
\]
Hence
\[
 P_{40,27}=P_{8,3}P_{5,2}.
\]
An index such as $107$ is selected because $107\equiv3\pmod8$ and $107\equiv2\pmod5$.  An index $67$ has the same pair and is also $27$ modulo $40$.  No other residue modulo $40$ has that ordered pair.

\subsection{A complete compiler trace}
Consider the specification
\[
 A(z)=\frac{1}{1-z-z^2},
 \qquad
 b_n=\sum_{k=0}^{n}\binom nka_k,
 \qquad
 u_m=b_{3m+1}.
\]
The compiler proceeds as follows.
\begin{enumerate}
\item Reduce $A$ to the Fibonacci linear state with $a_n=F_{n+1}$.
\item Conjugate by the binomial transform; characteristic roots shift by one.
\item Obtain the transformed recurrence $b_{n+2}=3b_{n+1}-b_n$.
\item Form the jump matrix of step three.  If the transformed roots are $\alpha^2,\beta^2$, their cubes are $\alpha^6,\beta^6$.
\item The packet recurrence has trace $L_6=18$ and determinant $1$:
\[
 u_{m+2}=18u_{m+1}-u_m.
\]
\item Initial values are $u_0=b_1=2$ and $u_1=b_4=F_9=34$ because $b_n=F_{2n+1}$ for this shifted seed.
\item The next value is $18\cdot34-2=610=b_7=F_{15}$.
\end{enumerate}
This trace combines rational OGF normalization, linear realization, transform transport, residue jumping, and recurrence certification.

% ===== END sections/A_worked_examples.tex =====

% ===== BEGIN sections/B_formula_registry.tex =====
\section{Formula registry}

The registry assigns stable identifiers to the principal formulas used by the paper.  ``Derived'' means proved here from earlier displayed identities; ``classical input'' means the statement is used as a standard theorem; ``conjecture'' means it is not used to support a proved claim.

\small
\begin{longtable}{L{.08\textwidth}L{.22\textwidth}L{.36\textwidth}L{.12\textwidth}}
\caption{Core formula registry.}\\
\toprule
ID&Object&Formula or defining relation&Status\\
\midrule
\endfirsthead
\toprule
ID&Object&Formula or defining relation&Status\\
\midrule
\endhead
PS-01&Residue projector&$P_{N,j}a=(a_n\mathbf1_{n\equiv j\,(N)})_n$&definition\\
PS-02&Reconstruction&$\sum_jP_{N,j}=I$&derived\\
PS-03&Refinement&$P_{N,j}=\sum_{r\equiv j\,(N)}P_{M,r}$ for $N\mid M$&derived\\
PS-04&Fourier projector&$P_{N,j}A=N^{-1}\sum_r\omega^{-jr}A(\omega^rz)$&derived\\
PS-05&CRT packet&$P_{MN,r}=P_{M,r}P_{N,r}$ for $(M,N)=1$&derived\\
PS-06&Packet product&$(AB)_k=\sum_{r+s=k}A_rB_s$&derived\\
PS-07&Nonlinear conjugacy&$F^{[N]}=\Analysis_NF\Synthesis_N$&derived\\
RL-01&Linear realization&$a_n=\lambda C^nv$&definition\\
RL-02&Cayley--Hamilton recurrence&$\chi_C(C)=0$ applied to $C^nv$&classical input\\
RL-03&Hankel rank&$\min d=\rank(a_{i+j})$&proved\\
RL-04&Jump state&$a_{Nm+j}=\lambda(C^N)^mC^jv$&proved\\
RL-05&Fibonacci packet&$u_{m+2}=L_Nu_{m+1}-(-1)^Nu_m$&proved\\
ST-01&Active modal form&$a_n=\sum_\lambda P_\lambda(n)\lambda^n$&classical input / normalized here\\
ST-02&Torsion equivalence&$\lambda\sim_N\eta\iff(\lambda/\eta)^N=1$&definition\\
ST-03&Phase polynomial&$R_{C,j}(x)=\sum_{\lambda\in C}(\lambda/\lambda_C)^jP_\lambda(x)$&definition\\
ST-04&Modal packet transfer&$a_{Nm+j}=\sum_C\lambda_C^jR_{C,j}(Nm+j)(\lambda_C^N)^m$&proved\\
ST-05&Minimal packet polynomial&$\chi_{N,j}=\prod_{C:R_{C,j}\neq0}(X-\lambda_C^N)^{1+\deg R_{C,j}}$&proved\\
ST-06&Exact packet rank&$\rho_a(N,j)=\sum_{C:R_{C,j}\neq0}(1+\deg R_{C,j})$&proved\\
ST-07&Rank-drop criterion&rank drops iff $\lambda\mapsto\lambda^N$ is noninjective on active roots&proved\\
ST-08&Torsion period&$L_a=\operatorname{lcm}\{\operatorname{ord}(\lambda/\eta)\}$&definition\\
ST-09&Double periodicity&$\rho(N+L_a,j)=\rho(N,j)=\rho(N,j+L_a)$&proved\\
ST-10&Torsion signature&$\mathfrak T_a:(\bbZ/L_a\bbZ)^2\to\bbZ_{\ge0}$&proved invariant\\
ST-11&Deficit split&global deficit $=$ collision deficit $+$ phase deficit&proved\\
ST-12&Rank-count law&$c_q(tL_a+s)=tA_{q,s}+B_{q,s}$&proved\\
ST-13&Fourier reconstruction&$P_\zeta=L^{-1}\sum_j\zeta^{-j}R_j$&proved\\
ST-14&Packet uncertainty&$s_dt_d\ge L$&proved\\
ST-15&Finite transient split&$a=h+b$ with $h$ finite support and $b$ persistent&proved\\
HL-01&P-recursive state&$v_{n+1}=C(n)v_n$&derived\\
HL-02&Multisection closure&$P_{N,j}$ preserves D-finiteness&classical closure\\
HL-03&Hypergeometric packet ratio&$u_{m+1}/u_m=\prod_{s=0}^{N-1}R(Nm+j+s)$&proved\\
HL-04&Ap\'ery recurrence&$(n+1)^3A_{n+1}=(34n^3+51n^2+27n+5)A_n-n^3A_{n-1}$&classical anchor\\
DG-01&Digit realization&$a_n=u^TT_{d_\ell}\cdots T_{d_0}v$&definition\\
DG-02&Thue--Morse&$t_{2n}=t_n$, $t_{2n+1}=1-t_n$&proved from bit parity\\
DG-03&Affine lift&$x\mapsto A_dx+b_d$ lifted by one coordinate&proved\\
TR-01&Riordan action&$\mathcal R_{g,f}A=gA\circ f$&definition\\
TR-02&Riordan product&$(g,f)(h,\ell)=(g(h\circ f),\ell\circ f)$&proved\\
TR-03&Transported projector&$Q_{N,j}=TP_{N,j}T^{-1}$&proved\\
TR-04&Binomial transform&$b_n=\sum_k\binom nka_k$&definition\\
TR-05&Stirling transform&$\widehat B(t)=\widehat A(e^t-1)$&derived\\
NL-01&Packet Fa\`a di Bruno&set partitions with packet-labeled derivative factors&proved\\
NL-02&Bell EGF&$\mathcal B=e^{e^t-1}$&derived seed\\
NL-03&Bell recurrence&$B_{n+1}=\sum_k\binom nkB_k$&proved\\
NL-04&Catalan fixed point&$u=z+u^2$&definition\\
NL-05&Catalan parity&$u_0=u_0^2+u_1^2$, $u_1=z+2u_0u_1$&proved\\
NL-06&Lagrange inversion&$[z^n]u^k=(k/n)[w^{n-k}]\Phi(w)^n$&classical input\\
NL-07&Motzkin&$M=1+zM+z^2M^2$&definition\\
FP-01&Filtered product&$F_m\equiv1\pmod{z^{w(m)}}$, $w(m)\to\infty$&proved completion\\
FP-02&Euler transform&$A=\prod_m(1-z^m)^{-b_m}$&definition\\
FP-03&Euler recurrence&$na_n=\sum_{k=1}^nc_ka_{n-k}$&proved\\
FP-04&Exponent inversion&$b_n=n^{-1}\sum_{d\mid n}\mu(n/d)c_d$&proved\\
AR-01&Dirichlet product&$(f\Dir g)(n)=\sum_{d\mid n}f(d)g(n/d)$&definition\\
AR-02&Prime-power synthesis&$f(n)=\prod_{p^e\parallel n}f(p^e)$&proved\\
AR-03&Local convolution&$(f\Dir g)(p^e)=\sum_{r=0}^ef(p^r)g(p^{e-r})$&proved\\
AR-04&Multiplicative packet&$P_k(f\Dir g)=\sum_{rs=k}f_r\Dir g_s$&proved\\
SF-01&Exponential section&$\mathcal E_{N,j}=\sum_nz^{Nn+j}/(Nn+j)!$&definition\\
SF-02&Signed denominator&$\mathcal D_N=\sum_n(-1)^nz^{Nn}/(Nn)!$&definition\\
SF-03&Signed ODE&$\mathcal D_{N,j}^{(N)}=-\mathcal D_{N,j}$&proved\\
SF-04&Transfer secant&$\operatorname{Sec}_N^{[r]}=\mathcal D_N^{-r}$&definition\\
SF-05&Secant semigroup&$P_{N,n}(r+s)=\sum_q\binom{Nn}{Nq}P_{N,q}(r)P_{N,n-q}(s)$&proved\\
SF-06&Transfer tangent&$\operatorname{Tan}_N=-\mathcal D_N'/\mathcal D_N$&definition\\
SF-07&Higher Riccati&$Y_N(-T,-T',\ldots)=-1$&proved\\
SF-08&Congruential Euler&$\mathcal C_{N,j}=(z^j/j!)/\mathcal E_{N,j}$&definition\\
SF-09&Triangular inverse&$\sum_q\binom{Nn+j}{Nq}C_{N,j}(q)=\delta_{n0}$&proved\\
SF-10&Block Appell&$\mathcal C_{N,j}^re^{xz}=\sum_mE_m^{(N,j;r)}(x)z^m/m!$&proved\\
LA-01&Odd-power coordinate&$S_n(p)=\sum_{a=0}^{p-1}(-1)^a(2a+1)^n$&definition\\
LA-02&Odd Euler inverse&$E_n=\sum_m(-1)^mZ_{2m+1}\binom n{2m}p^{2m}S_{n-2m}$&proved\\
LA-03&Odd tail&error after $M$ lies in $p^{2M+2}\bbZ_p$&proved\\
LA-04&Dyadic source&$R_n=(3^{n+2}-1)/(4(n+1)(n+2))$&definition\\
LA-05&Dyadic inverse&$z^2/\sinh^2z=\sum_m\beta_mz^{2m}$&definition\\
LA-06&Unified local kernel&$E_n=\sum_m\binom n{2m}\kappa_{p,m}A_{p,n-2m}$&proved\\
LA-07&Exact dyadic kernel&$v_2(\kappa_{2,m})=4m-1$ for $m\geq1$&proved\\
LA-08&Sharp dyadic tail&error after $M$ lies in $2^{4M+3}\bbZ_2$&proved optimal\\
LA-09&Standard-window period&$e_{n+h}^{(m,p)}=e_n^{(m,p)}$ for $n\geq1$&proved\\
LA-10&Mahler reconstruction&$f(n)=\sum_r\Delta^rf(0)\binom nr$&classical input\\
LA-11&Scalar Carmichael--Mahler&$u^{Hx+j}=u^j\sum_r(u^H-1)^r\binom xr$&proved\\
LA-12&Matrix Carmichael--Mahler&$a_{Hx+j}=\sum_rp^{Lr}\binom xr w^TC^jB^rv_0$&proved\\
LA-13&Adelic transfer&$\Delta_H^ra_n\equiv0\pmod{Q^r}$&proved\\
LA-14&Odd depleted Euler&$(1-\sigma_pp^n)E_n=K_p*S_n^\times$&proved\\
LA-15&Dyadic interpolation&$\mathscr E_2(n)=E_{2n}$ on $\bbN$&proved\\
PP-01&Discrete polyphase&$a_r^{[M]}(q)=a_{Mq+r}$&definition\\
PP-02&Nested refinement&$(a_r^{[M]})_s^{[N]}=a_{r+Ms}^{[MN]}$&proved\\
PP-03&Carry convolution&$C_r=\sum_{u+v=r}A_uB_v+w\sum_{u+v=r+M}A_uB_v$&proved\\
PP-04&Mahler kernel&$\Theta_{M,s}(r,k)=[z^r](1+z)^s((1+z)^M-1)^k$&definition\\
PP-05&Mahler transfer&$d_{s,k}=\sum_{r=k}^{Mk+s}c_r\Theta_{M,s}(r,k)$&proved\\
PP-06&Kernel composition&$\Theta_{M,s}\Theta_{N,t}=\Theta_{MN,Mt+s}$&proved\\
PP-07&Valuation preservation&$v_p(c_r)\geq w(r)\Rightarrow v_p(d_{s,k})\geq w(k)$&proved\\
PP-08&Frobenius sparsity&$\Theta_{p^a,s}(r,k)\equiv\binom{s}{r-p^ak}\pmod p$&proved\\
PP-09&Precision amplification&$C^{Hp^a}\equiv I\pmod{p^{L+a}}$&proved\\
CP-01&Compiler soundness&finite evaluator agrees with semantic object&proved\\
CP-02&Packet compiler&synthesis after packet compilation equals scalar compilation&proved\\
CP-03&Finite-prefix identifiability&nonunique continuations $P+c\binom n{m+1}$&proved negative theorem\\
\bottomrule
\end{longtable}
\normalsize

\subsection{Dependency map}
The registry has a short logical spine:
\[
\begin{aligned}
\text{coefficient semantics}&\Longrightarrow\text{exact projectors}\Longrightarrow\text{typed closure laws}\\
&\Longrightarrow\text{family realizations}\Longrightarrow\text{packet/local refinements}\\
&\Longrightarrow\text{certified algorithms}.
\end{aligned}
\]
The spectral transfer formulas depend on the active exponential-polynomial decomposition and exact finite Fourier grouping, not on numerical rank experiments.  The odd and dyadic Euler formulas depend on their own generating-function factorizations, not on conjectural local periodicity.  The compiler theorem depends on well-foundedness and exact closure, not on any database-coverage claim.

% ===== END sections/B_formula_registry.tex =====

% ===== BEGIN sections/C_algorithms.tex =====
\section{Reference algorithms and source listings}

The companion code is intentionally small enough to audit.  It is not a complete implementation of the proposed compiler.  It provides exact certificates for the central realization, spectral-transfer, nonlinear, digital, product, local-arithmetic, adelic, and polyphase examples and regenerates every figure in this paper.

\subsection{Reproduction layout}
\begin{verbatim}
Projector_Seed_Integer_Sequence_Calculus/
  main.tex
  Projector_Seed_Calculus_for_Integer_Sequences_Local_Polyphase_Extension.tex
  sections/
  code/verify_sequences.py
  code/generate_figures.py
  code/flatten_latex.py
  figures/
  results/verification_summary.json
  results/verification_run.txt
  tables/verification_ledger.csv
\end{verbatim}

Run from the project root:
\begin{verbatim}
python code/verify_sequences.py
python code/generate_figures.py
python code/flatten_latex.py
pdflatex -interaction=nonstopmode -halt-on-error main.tex
pdflatex -interaction=nonstopmode -halt-on-error main.tex
\end{verbatim}

\subsection{Finite standard-window inversion}
For fixed $p$ and $m\mid p-1$, the residual inverse is computed by the unit-triangular recursion
\[
 e_0=1,
 \qquad
 e_n=-\sum_{q=0}^{n-1}\binom{mn}{mq}e_q\pmod p.
\]
Only the first $h=(p-1)/m$ positive-index entries need to be retained; the Fourier-transfer theorem certifies their periodic repetition.

\subsection{Affine Mahler polyphase compilation}
Given a Mahler prefix $c_0,\ldots,c_D$, radix $M$, and phase $s$, compute
\[
 \Theta_{M,s}(r,k)=[z^r](1+z)^s((1+z)^M-1)^k
\]
for $0\leq k\leq D$ and $k\leq r\leq\min(D,Mk+s)$, then return
\[
 d_{s,k}=\sum_{r=k}^{\min(D,Mk+s)}c_r\Theta_{M,s}(r,k).
\]
The certificate records both this direct table and the mixed-radix composition identity when refinement is staged.

\subsection{Exact verification program}
The program includes direct packet reconstruction, exact Hankel-rank verification for four torsion-transfer families, exact double-periodicity checks, the complete mixed period-six signature, rank-count quasipolynomials, Fourier reconstruction, and the nonlinear, sharp dyadic, standard-window, Carmichael--Mahler, and affine polyphase certificates.

\lstinputlisting[caption={Exact sequence, spectral-transfer, and arithmetic verification suite.}]{code/verify_sequences.py}

\subsection{Figure-generation program}
The plotting layer uses floating-point coordinates only for visualization.  Every discrete value displayed in a figure is generated from exact integer or rational formulas before conversion for plotting.  The transfer figures visualize spectral packet maps, finite torsion signatures, local kernel weights, standard-window Fourier transfer, odd/dyadic geometry, Mahler polyphase kernels, and precision refinement trees.

\lstinputlisting[caption={Program generating the twenty-four figures used in the local and polyphase extension.}]{code/generate_figures.py}

% ===== END sections/C_algorithms.tex =====

% ===== BEGIN sections/D_verification.tex =====
\section{Verification protocol and certificate summary}

\subsection{Exact test policy}
The test suite uses Python arbitrary-precision integers, exact rational numbers, symbolic polynomial identities, and exact Hankel-matrix ranks.  Floating-point comparison is never used to decide an identity, recurrence, congruence, rank, or valuation bound.

The deterministic random seed is
\[
 20260812.
\]
Randomized tests are used only to sample exact algebraic identities over bounded integer inputs; every sampled assertion is exact.

\subsection{Summary}
\begin{certificate}[Unified verification run]
The reference suite completed
\[
 \boxed{87{,}915\text{ exact checks}}
\]
with
\[
 \boxed{0\text{ failures}.}
\]
\end{certificate}

\noindent\begin{tabularx}{\linewidth}{@{}L{.23\linewidth}L{.18\linewidth}Y@{}}
\toprule
Mechanism&Scale&Checked property\\
\midrule
Residue packet algebra&multiple moduli and random inputs&reconstruction and cyclic product laws\\
CRT refinement&four coprime modulus pairs&equivalence of joint and factored addresses\\
Fibonacci/Lucas jumps&steps $1$ through $9$&common jump recurrence for every packet\\
Spectral torsion ranks&four exact sequence families, $N\leq18$, twelve addresses&predicted minimal packet rank equals exact finite Hankel rank\\
Torsion periodicity&mixed period-six family, $N\leq30$&exact periodicity in both $N$ and $j$\\
Finite torsion signature&complete $6\times6$ table&all rank entries and rank-count quasipolynomials\\
Fourier reconstruction&order-two Jordan family&exact recovery of both modal polynomials\\
Bell&through degree $18$&initial values and binomial recurrence\\
Catalan&through degree $19$&quadratic convolution recurrence\\
Motzkin&through degree $20$&reference initial segment\\
Partitions&through degree $50$&divisor recurrence and product stabilization\\
Ap\'ery&through degree $14$&reference terms and cubic recurrence\\
Binomial transform&100 random vectors&exact forward/inverse recovery\\
Thue--Morse&indices below $4096$&bit parity and both radix recursions\\
Odd-prime Euler&$p=3,5,7,11$, even rows through $24$&exact reconstruction and every truncation modulus\\
Sharp dyadic Euler&even rows through $24$ plus optimal rows&exact reconstruction, exact kernel valuations, unit sources, and $2^{4M+3}$ tails\\
Standard-window transfer&eight odd primes and every divisor $m\mid p-1$&inverse equations and exact period $h=(p-1)/m$\\
Carmichael--Mahler&four odd primes and refined Fibonacci matrices&scalar expansions, coefficient weights, matrix amplification, and finite differences\\
Mahler polyphase&radices through six and mixed cascades&affine binomial identity and exact kernel composition\\
Frobenius sparsity&four primes and two prime-power levels&residual short-window formula for $\Theta_{p^a,s}$\\
Polyphase valuation&three primes and four radices&preservation of monotone $p$-adic coefficient ladders\\
Finite-prefix ambiguity&five perturbations&agreement on prefix and divergence afterward\\
\bottomrule
\end{tabularx}

\subsection{Spectral-transfer certificate design}
The new transfer checks are deliberately independent of the printed closed formulas at the final comparison stage.  For each sampled modulus and address, the program:
\begin{enumerate}
\item generates packet terms directly from the unsplit defining formula;
\item forms an exact finite Hankel matrix over the integers;
\item computes its exact rank symbolically;
\item compares that rank with the torsion-phase prediction;
\item separately checks the finite period-six signature and its rank-count formula.
\end{enumerate}
For the tested families, the Hankel blocks are larger than the predicted ranks, so the comparison detects both overestimation and underestimated recurrence order within the sampled range.

\subsection{Reference coefficient anchors}
The run records the following initial segments.
\[
\begin{aligned}
B_n:&\quad1,1,2,5,15,52,203,877,4140,21147,115975,\\
C_n:&\quad1,1,2,5,14,42,132,429,1430,4862,16796,\\
M_n:&\quad1,1,2,4,9,21,51,127,323,835,2188,\\
p(n):&\quad1,1,2,3,5,7,11,15,22,30,42,56,77,101,135,176,\\
A_n:&\quad1,5,73,1445,33001,819005,21460825,584307365.
\end{aligned}
\]
The even signed Euler anchors are
\[
\begin{aligned}
1,&-1,5,-61,1385,-50521,2702765,-199360981,\\
&19391512145,-2404879675441,370371188237525.
\end{aligned}
\]

\subsection{What the suite certifies}
The suite supplies strong regression and finite-instance evidence that:
\begin{itemize}
\item implementations match the printed formulas;
\item packet synthesis and operation laws are exact on sampled data;
\item the transfer theorem's four anchor rank profiles match independently computed exact Hankel ranks;
\item the mixed torsion signature and its double periodicity agree on the declared finite range;
\item the classwise Fourier reconstruction example is exact;
\item representative sequence values and recurrences agree;
\item every tested odd and dyadic truncation meets its proved modulus, including the sharp dyadic exponent;
\item standard-window inverses satisfy the proved finite period on all tested prime-section pairs;
\item scalar and matrix Mahler packets satisfy their predicted valuation ladders;
\item affine polyphase kernels satisfy direct substitution, composition, Frobenius sparsity, and valuation preservation;
\item figures are generated from the same formulas as the text.
\end{itemize}

\subsection{What the suite does not certify}
Finite tests do not replace proofs of universal statements.  In particular, the suite does not establish:
\begin{itemize}
\item the master realization theorem independently of its proof;
\item the spectral torsion--Fourier theorem independently of its proof;
\item untested historical novelty;
\item canonical minimality outside finite linear systems and automata;
\item any claim about the fraction of an external database covered.
\end{itemize}

\subsection{Independent checking protocol}
An independent auditor can proceed in increasing order of trust.
\begin{enumerate}
\item Inspect the formulas and proof dependencies without executing code.
\item Re-derive the classwise transfer identity and polynomial-mode annihilator lemma.
\item Run the verification suite in a clean environment.
\item Replace the exact Hankel-rank and arithmetic implementations with independent ones.
\item Recompute selected torsion signatures from minimal jump realizations rather than modal formulas.
\item Recompute standard-window signatures and polyphase kernels with independent finite-field and polynomial code.
\item Check local congruence tails at primes and rows not included in the default suite.
\item Audit each family and transfer formulation for historical priority separately from internal correctness.
\end{enumerate}

\subsection{Machine-readable result}
The file \texttt{results/verification\_summary.json} contains status, check count, failure count, family labels, and coefficient anchors.  The CSV ledger contains the compact run metrics.  The text transcript records the exact output produced by the suite.

% ===== END sections/D_verification.tex =====

% ===== BEGIN sections/E_claim_ledger.tex =====
\section{Claim-status and dependency ledger}

\subsection{Proved internal statements}
The following are proved in the paper under their stated hypotheses:
\begin{enumerate}
\item exact residue reconstruction, refinement, Fourier representation, and CRT factorization;
\item exact packet operation laws and nonlinear conjugacy;
\item formal fixed-point existence under strict coefficient contraction;
\item finite-prefix and filtered-product closure;
\item the master realization and exact packet-compiler theorems;
\item Cayley--Hamilton recurrence, Hankel-rank minimality, and jump realization;
\item active exponential-polynomial normalization for persistent C-finite sequences;
\item the spectral torsion--Fourier packet transfer theorem;
\item the exact minimal recurrence, reduced rational denominator, Hankel rank, and state dimension of every persistent C-finite arithmetic subsequence;
\item the complete root-of-unity rank-drop criterion, collision--phase deficit decomposition, exact double periodicity, finite torsion signature, quasipolynomial rank counts, Fourier information conservation, and packet uncertainty;
\item the finite-support transient extension to arbitrary C-finite sequences;
\item Fibonacci and general Lucas packet laws, with exact minimality in torsion-free cases;
\item hypergeometric packet ratios, finite nonautonomous jumps, automatic digit refinement, and affine lifts;
\item Riordan composition, inverse, transform-transported projectors, Bell differential and partition realizations, packet-labeled tree expansion, and Catalan parity equations;
\item Euler-transform recurrence, exponent inversion, prime-power synthesis, and Dirichlet packet reconstruction;
\item signed transfer differential systems, reciprocal semigroup, higher Riccati law, congruential Euler inverse, and block-Appell formulas;
\item exact odd-prime Euler reconstruction and its uniform $p^{2M+2}$ ladder;
\item the unified local Euler source-kernel theorem;
\item exact dyadic source units, $v_2(\beta_m)=s_2(m)-1$, exact kernel depth $v_2(\kappa_{2,m})=4m-1$, and the optimal $2^{4M+3}$ truncation theorem;
\item finite-field standard-window transfer and the exact period $e_{n+h}^{(m,p)}=e_n^{(m,p)}$;
\item finite Mahler reconstruction and difference valuation bounds;
\item scalar, polynomial-exponential, matrix, and finite adelic Carmichael--Mahler transfer;
\item odd-prime depleted Euler reconstruction and analytic interpolation;
\item integral dyadic source interpolation and direct analytic interpolation of $E_{2n}$;
\item exact discrete polyphase analysis, synthesis, nested mixed-radix factorization, carry convolution, and the shift matrix law;
\item the affine Mahler polyphase kernel, its finite support, and its exact composition law;
\item polyphase valuation preservation, prime-power Frobenius sparsification, and scalar and matrix precision amplification;
\item compiler soundness for legal typed specifications, including packet, local, and polyphase compilation;
\item completeness of the spectral packet compiler on the persistent C-finite stratum.
\end{enumerate}

\subsection{Classical inputs}
The exposition uses the following standard inputs without claiming them as new:
\begin{itemize}
\item Cayley--Hamilton, Jordan decomposition, and finite-dimensional realization theory;
\item equivalence of rational OGFs and constant recurrences;
\item exponential-polynomial descriptions of linear recurrence sequences;
\item finite Fourier inversion, character orthogonality, and root-of-unity filters;
\item P-recursive/D-finite correspondence and standard closure properties;
\item finite automata, $k$-kernels, and substitution realizations;
\item Fa\`a di Bruno and Lagrange inversion;
\item M\"obius inversion and multiplicative-function factorization;
\item Lucas' theorem, Kummer's theorem, lifting the exponent, Newton interpolation, and Mahler coordinates;
\item the von Staudt--Clausen denominator theorem and Legendre digit-sum formulas;
\item Carmichael exponents of prime-power unit groups and the $p$-adic binomial series;
\item standard generating functions for Fibonacci, Bell, Catalan, partitions, Euler, Bernoulli, and related anchor families.
\end{itemize}

\subsection{Derived synthesis and transfer contributions}
The paper's contribution is not the isolated rediscovery of every classical ingredient. It organizes them into one typed realization, transfer, local-completion, and refinement theory. Specific contributions include:
\begin{itemize}
\item one distinction between storage, structural generation, and intended-rule discovery;
\item one restricted grammar spanning linear, nonlinear, digital, product, multiplicative, local, and polyphase systems;
\item a common exact packet lift for all legal operations;
\item a complete global-to-packet minimality law for the persistent C-finite stratum;
\item the torsion packet signature as a finite invariant of infinitely many arithmetic-subsequence complexities;
\item one source-kernel formulation that unifies odd and dyadic Euler reconstruction while retaining distinct local weights;
\item a finite-field Fourier proof that closes the standard-window period problem;
\item the exact dyadic transfer profile and its optimal uniform precision ladder;
\item a state-level theorem transferring finite modular dynamics into local analytic Mahler packets and simultaneous adelic congruences;
\item the affine Mahler polyphase kernel as an exact coefficient-level refinement operator;
\item exact mixed-radix kernel composition, residual Frobenius sparsity, and valuation-preserving refinement;
\item the character-polyphase lattice separating torsion phase, local digit address, and Mahler depth;
\item a joint complexity profile connecting seed size, packet width, spectral torsion, local depth, polyphase width, and certificate cost;
\item a compiler architecture that retains equivalent presentations, local precision ledgers, and proof diagrams.
\end{itemize}
Historical priority for particular terminology, generalized families, or equivalent formulations remains an external audit task.

\subsection{Finite computational certificates}
The \textbf{87,915 exact checks} certify implementation-level agreement on the declared finite ranges. They include exact Hankel-rank comparisons, torsion signatures, Fourier reconstruction, odd and sharp dyadic Euler identities, standard-window periods, scalar and matrix Mahler transfer, affine polyphase identities, mixed-radix kernel composition, Frobenius sparsity, and valuation preservation. They do not elevate conjectures to theorems.

\subsection{Conjectural and open statements}
The following remain open within this paper:
\begin{enumerate}
\item canonical or stratified normal forms for a broad mixed grammar;
\item minimal multisection operators and finite signatures beyond the C-finite stratum;
\item effective transform-orbit invariants and transfer of torsion signatures under restricted transforms;
\item sharp algebraic packet degree and branch-collision bounds;
\item compressed packet Fa\`a di Bruno algorithms and nonlinear tree-class transfer signatures;
\item a nontrivial additive--multiplicative bridge theory;
\item exact higher valuations of $\Theta_{p^a,s}(r,k)$;
\item compatible prime-power lifts of standard-window signatures;
\item odd/dyadic source-kernel systems beyond the Euler family;
\item canonical local normal forms and a combined spectral-local signature;
\item optimal mixed-radix polyphase trees;
\item structural universality in the stronger canonical and multi-stratum minimality sense.
\end{enumerate}

\subsection{Resolved open statements from earlier editions}
Three formerly open targets are resolved in this edition sequence:
\begin{enumerate}
\item the exact C-finite packet-rank spectrum and periodic control of rank drops, by spectral torsion--Fourier transfer;
\item the standard-window period $e_{n+h}=e_n$ over $\bbF_p$, by finite-field Fourier transfer;
\item the optimal dyadic Euler transfer weight, by the exact formula $v_2(\kappa_{2,m})=4m-1$ and the sharp tail modulus $2^{4M+3}$.
\end{enumerate}

\subsection{No hidden implication}
The realization implication remains one-way:
\[
 \text{valid finite grammar}\Longrightarrow\text{exact sequence realization}.
\]
The transfer implications are stratified:
\[
 \text{persistent C-finite spectral data}
 \Longrightarrow
 \text{exact minimal complexity of every arithmetic subsequence},
\]
\[
 \text{finite modular state data}
 \Longrightarrow
 \text{local Mahler packets and congruence ladders},
\]
and
\[
 \text{Mahler coefficient profile}
 \Longrightarrow
 \text{exact polyphase branches with certified precision}.
\]
None asserts that an arbitrary sequence has a short legal grammar, that finite terms reveal an intended specification, or that every grammar stratum already has a minimality theorem.

\subsection{Relation to the wider projector-seed series}
This monograph consolidates sequence-specific material distributed across foundational operator presentations, refinement calculus, nonlinear packet trees, analytic special-function families, and arithmetic packet systems. It does not import the quantum, PDE, computational-hardware, local Langlands, or global automorphy claims of the wider series.

\subsection{Final status statement}
The resulting object is a self-contained sequence calculus with three established transfer layers:
\begin{enumerate}
\item one global spectral invariant determines the exact minimal recurrence and complexity of every persistent C-finite arithmetic subsequence;
\item finite character and modular state data determine local analytic branches, exact valuation ladders, and finite adelic congruences;
\item affine polyphase kernels determine every mixed-radix refinement of Mahler coordinates and preserve or amplify certified local precision.
\end{enumerate}
The next threshold is a minimality theorem for a nonlinear or holonomic stratum and a combined spectral-local signature.

% ===== END sections/E_claim_ledger.tex =====

\end{document}
